\documentclass[a4paper,11pt]{article}
        \usepackage[top=15mm,bottom=18mm,left=16mm,right=16mm]{geometry}
        \usepackage{fontspec}
        \usepackage{xeCJK}
        \setmainfont{Times New Roman}
        \setCJKmainfont{Yu Gothic}
        \setmonofont{Consolas}
        \setCJKmonofont{Yu Gothic}
        \usepackage{booktabs}
        \usepackage{longtable}
        \usepackage{tabularx}
        \usepackage{array}
        \usepackage{enumitem}
        \usepackage{hyperref}
        \usepackage{listings}
        \usepackage{xcolor}
        \usepackage{amsmath}
        \usepackage{amssymb}
        \hypersetup{
          colorlinks=true,
          linkcolor=blue,
          urlcolor=blue,
          pdftitle={offline フルレース simulation 本体},
          pdfauthor={Codex}
        }
        \lstset{
          basicstyle=\ttfamily\small,
          breaklines=true,
          columns=fullflexible,
          keepspaces=true,
          frame=single,
          framerule=0.2pt,
          rulecolor=\color{black},
          showstringspaces=false
        }
        \setlength{\parskip}{0.42em}
        \setlength{\parindent}{1em}
        \setlength{\emergencystretch}{3em}
        \renewcommand{\arraystretch}{1.15}
        \title{10. offline フルレース simulation 本体}
        \author{solar\_ws0129-main}
        \date{2026-07-29}
        \begin{document}
        \maketitle
        \tableofcontents
        \clearpage

        \section{対象}
        \begin{tabularx}{\linewidth}{>{\raggedright\arraybackslash}p{0.18\linewidth} X}
        \toprule
        項目 & 内容 \\
        \midrule
        ファイル & \path|scripts/solar_sim.py| \\
        ソースSHA-256 & \path|72dbec56f03ec2611bd7c68bfd8ac9df1c3a9b173b9a1b97e754631bce626eac| \\
        種別 & Python \\
        区分 & offline core \\
        役割 & profile、forecast、route、maps を使って全レースを逐次再生し、upper/lower 相当の実行を CSV と HTML に落とす。 \\
        起動文脈 & simulate や historical-simulate で直接呼ばれる同期版の基準実装。 \\
        \bottomrule
        \end{tabularx}

        \section{このファイルを読む理由}
        profile、forecast、route、maps を使って全レースを逐次再生し、upper/lower 相当の実行を CSV と HTML に落とす。

        \begin{itemize}[leftmargin=1.6em]
        \item SolarCarModel を直に持つ同期版なので、数理理解の最短入口。
\item full summary JSON、detail CSV、upper plan CSV、HTML report を生成する。
\item live の mpc\_node とかなり同型の距離上位計画ロジックを持つ。
        \end{itemize}

        \section{呼び出し関係}
        \subsection{どこから呼ばれるか}
        \begin{itemize}[leftmargin=1.6em]
        \item \path|scripts/solar_control.sh|
        \end{itemize}

        \subsection{次に読むと理解が進むファイル}
        \begin{itemize}[leftmargin=1.6em]
        \item \path|mpc_solarcar/model.py|
\item \path|mpc_solarcar/upper_horizon.py|
\item \path|mpc_solarcar/upper_solver.py|
        \end{itemize}

        \subsection{同一リポジトリ内の主要依存}
        \begin{itemize}[leftmargin=1.6em]
        \item \path|mpc_solarcar/model.py|
\item \path|mpc_solarcar/route_utils.py|
\item \path|mpc_solarcar/schedule_utils.py|
\item \path|mpc_solarcar/signal_utils.py|
\item \path|mpc_solarcar/solar_profile.py|
\item \path|mpc_solarcar/upper_cost.py|
\item \path|mpc_solarcar/upper_horizon.py|
\item \path|mpc_solarcar/upper_policy.py|
\item \path|mpc_solarcar/upper_solver.py|
        \end{itemize}

        \section{ソース冒頭の把握}
        モジュール docstring または先頭意図の要約:

        モジュール docstring は明示されていない。

        \subsection{代表コード断片}
        \begin{lstlisting}[language=Python]
os.makedirs(parent, exist_ok=True)


class DetailCsvStream:
    """Write the 1 Hz execution trace without retaining the full race in RAM."""

    def __init__(self, path):
        self.path = str(path)
        self._file = None
        self._writer = None
        self.fieldnames = []
        self.row_count = 0

    def write(self, row):
        if self._writer is None:
            ensure_parent_dir(self.path)
            self.fieldnames = list(row.keys())
            opener = gzip.open if self.path.lower().endswith('.gz') else open
            self._file = opener(self.path, 'wt' if opener is gzip.open else 'w', encoding='utf-8', newline='')
            self._writer = csv.DictWriter(self._file, fieldnames=self.fieldnames, extrasaction='ignore')
            self._writer.writeheader()
        self._writer.writerow(row)
\end{lstlisting}


        \section{主要構造の読み取り}
        主要クラスは DetailCsvStream。 主要関数は load\_yaml, sim\_log, select\_optimized\_vector, limit\_step\_duration\_to\_distance, terminal\_soc\_predictions, advance\_rate\_limiter\_to\_distance\_boundary, snap\_execution\_stop\_speed\_kmh, choose\_integration\_step\_seconds。 CLI 引数宣言は 96 件。

        \subsection{主要クラス・関数}
            \begin{longtable}{p{0.07\linewidth} p{0.16\linewidth} p{0.25\linewidth} p{0.40\linewidth}}
            \toprule
            行 & 種別 & 名前 & 補足 \\
            \midrule
            549 & class & \path|DetailCsvStream| &  \\
63 & function & \path|load_yaml| & args: path \\
73 & function & \path|sim_log| & args: message \\
77 & function & \path|select_optimized_vector| & args: res, x0 \\
115 & function & \path|limit_step_duration_to_distance| & args: step\_sec, speed\_kmh, s0\_km, race\_km \\
129 & function & \path|terminal_soc_predictions| & args: upper\_solve\_log \\
145 & function & \path|advance_rate_limiter_to_distance_boundary| & args: limiter, target\_kmh \\
195 & function & \path|snap_execution_stop_speed_kmh| & args: speed\_kmh, target\_kmh \\
212 & function & \path|choose_integration_step_seconds| & args: simulation\_step\_sec, weather\_step\_sec \\
218 & function & \path|choose_stationary_integration_step_seconds| & args: simulation\_step\_sec, weather\_step\_sec, requested\_step\_sec \\
230 & function & \path|write_model_snapshot| & args: detail\_csv, parameters, map\_paths \\
278 & function & \path|get_workspace_revision| & args: root\_dir \\
331 & function & \path|timestamp_ns| & args: value \\
335 & function & \path|load_stops| & args: path \\
365 & function & \path|load_profile| & args: path \\
373 & function & \path|forecast_distance_column| & args: df \\
383 & function & \path|load_forecast_dataframe| & args: path, tzname \\
410 & function & \path|merge_forecast_dataframes| & args: primary, fallback \\
423 & function & \path|build_forecast_grid_payload| & args: df \\
476 & function & \path|interp_forecast_grid| & args: payload, col, t\_utc, s\_km, default \\
523 & function & \path|load_progress_reference_dataframe| & args: path \\
543 & function & \path|ensure_parent_dir| & args: path \\
552 & function & \path|__init__| & args: self, path \\
559 & function & \path|write| & args: self, row \\
570 & function & \path|close| & args: self \\
576 & function & \path|__enter__| & args: self \\
579 & function & \path|__exit__| & args: self, exc\_type, exc, tb \\
583 & function & \path|_deep_copy_cfg| & args: cfg \\
587 & function & \path|_set_nested| & args: cfg, dotted\_key, value \\
601 & function & \path|apply_overrides| & args: cfg, overrides \\
618 & function & \path|build_config_tag| & args: profile\_path, profile\_cfg, overrides \\
633 & function & \path|tag_output_path| & args: path\_value, tag, default\_ext \\
645 & function & \path|apply_profile_cfg_to_args| & args: profile\_path, profile\_cfg, args \\
798 & function & \path|apply_profile_defaults| & args: args \\
806 & function & \path|resolve_config_path| & args: profile\_path, value \\
820 & function & \path|evaluate_model_validation_gate| & args: cfg, profile\_path \\
1032 & function & \path|get_profile_val| & args: df, s\_km, field, default \\
1040 & function & \path|parse_utc_arg| & args: raw\_value \\
1052 & function & \path|resolve_forecast_mode| & args: mode, df \\
1064 & function & \path|forecast_row_index| & args: df, sim\_t \\
1082 & function & \path|format_metric| & args: value, digits, default \\
1094 & function & \path|decimate_xy| & args: xs, ys, max\_points \\
1105 & function & \path|build_svg_chart| & args: xs, ys \\
1146 & function & \path|flatten_params_for_report| & args: cfg \\
1149 & function & \path|visit| & args: prefix, value \\
1166 & function & \path|write_simulation_report| & args: path, summary, detail\_df, params\_rows \\
1288 & function & \path|mpc_solve| & args: model, data, z0, Tb0, s0\_km, v0\_kmh, v\_max\_kmh, term\_soc\_min, w\_dv, w\_dv\_limit, dv\_max\_kmhps, w\_T, w\_speed\_limit, w\_current, speed\_profile, soc\_target, soc\_band, w\_soc\_target, w\_soc\_band, schedule, soc\_day\_end\_max, w\_soc\_day\_max, soc\_finish\_target, soc\_finish\_tol, w\_soc\_progress, w\_soc\_terminal, race\_km, soc\_day\_end\_target, soc\_day\_end\_tol, w\_soc\_day\_track \\
1305 & function & \path|quad_penalty| & args: x, cap \\
1312 & function & \path|cost| & args: v \\
1402 & function & \path|soc_guard_speed| & args: model, v\_kmh, z, Tb, s\_km, d0, route\_profile, mode, soc\_guard \\
1409 & function & \path|z_next_for| & args: v\_kmh\_local \\
1451 & function & \path|soc_day_guard_speed| & args: model, v\_kmh, z, Tb, s\_km, d0, route\_profile, target\_soc, tol \\
1458 & function & \path|z_next_for| & args: v\_kmh\_local \\
1485 & function & \path|main| &  \\
1852 & function & \path|integration_step_seconds| &  \\
1859 & function & \path|forecast_has_coverage| & args: t\_utc \\
1868 & function & \path|integrate_drive_time_between| & args: t\_start\_utc, t\_end\_utc \\
1885 & function & \path|reference_value_at_time| & args: t\_utc, field \\
1907 & function & \path|reference_distance_at_time| & args: t\_utc \\
1958 & function & \path|time_to_index| & args: dt\_utc \\
2095 & function & \path|remaining_day_budget| & args: sim\_t\_local, \_k\_idx \\
            \bottomrule
            \end{longtable}

        \subsection{ファイルを上から読んだときの定義順}
            \begin{longtable}{p{0.07\linewidth} p{0.84\linewidth}}
            \toprule
            行 & 内容 \\
            \midrule
            23 & \_numeric\_threads に str(os.environ.get('SOLAR\_NUMERIC\_THREADS', '1')) の結果を代入する。 \\
24 & ('OPENBLAS\_NUM\_THREADS', 'OMP\_NUM\_THREADS', 'MKL\_NUM\_THREADS', 'NUMEXPR\_NUM\_THREADS') を順に走査し、各要素を \_thread\_variable に入れて処理する。 \\
36 & ROOT に Path(\_\_file\_\_).resolve().parents[1] の結果を代入する。 \\
37 & 条件 str(ROOT) not in sys.path を判定し、真なら内部処理を行う。 \\
60 & DISTANCE\_EPS\_KM に 1e-06 の結果を代入する。 \\
63 & 関数 load\_yaml を定義する。 \\
73 & 関数 sim\_log を定義する。 \\
77 & 関数 select\_optimized\_vector を定義する。 \\
115 & 関数 limit\_step\_duration\_to\_distance を定義する。 \\
129 & 関数 terminal\_soc\_predictions を定義する。 \\
145 & 関数 advance\_rate\_limiter\_to\_distance\_boundary を定義する。 \\
195 & 関数 snap\_execution\_stop\_speed\_kmh を定義する。 \\
212 & 関数 choose\_integration\_step\_seconds を定義する。 \\
218 & 関数 choose\_stationary\_integration\_step\_seconds を定義する。 \\
230 & 関数 write\_model\_snapshot を定義する。 \\
278 & 関数 get\_workspace\_revision を定義する。 \\
331 & 関数 timestamp\_ns を定義する。 \\
335 & 関数 load\_stops を定義する。 \\
365 & 関数 load\_profile を定義する。 \\
373 & 関数 forecast\_distance\_column を定義する。 \\
383 & 関数 load\_forecast\_dataframe を定義する。 \\
410 & 関数 merge\_forecast\_dataframes を定義する。 \\
423 & 関数 build\_forecast\_grid\_payload を定義する。 \\
476 & 関数 interp\_forecast\_grid を定義する。 \\
523 & 関数 load\_progress\_reference\_dataframe を定義する。 \\
543 & 関数 ensure\_parent\_dir を定義する。 \\
549 & クラス DetailCsvStream を定義する。 \\
583 & 関数 \_deep\_copy\_cfg を定義する。 \\
587 & 関数 \_set\_nested を定義する。 \\
601 & 関数 apply\_overrides を定義する。 \\
618 & 関数 build\_config\_tag を定義する。 \\
633 & 関数 tag\_output\_path を定義する。 \\
645 & 関数 apply\_profile\_cfg\_to\_args を定義する。 \\
798 & 関数 apply\_profile\_defaults を定義する。 \\
806 & 関数 resolve\_config\_path を定義する。 \\
820 & 関数 evaluate\_model\_validation\_gate を定義する。 \\
1032 & 関数 get\_profile\_val を定義する。 \\
1040 & 関数 parse\_utc\_arg を定義する。 \\
1052 & 関数 resolve\_forecast\_mode を定義する。 \\
1064 & 関数 forecast\_row\_index を定義する。 \\
            \bottomrule
            \end{longtable}

        \subsection{import 群の詳細}
            \begin{longtable}{p{0.07\linewidth} p{0.33\linewidth} p{0.50\linewidth}}
            \toprule
            行 & import 文 & このファイルで読む理由 \\
            \midrule
            2 & \path|import argparse| & CLI 引数を宣言し、実行時パラメータを外から受け取るため。 このファイル内での主な使用位置は L1486, L1534。 \\
3 & \path|import csv| & CSV の逐次読込・逐次書込を行うため。 このファイル内での主な使用位置は L565。 \\
4 & \path|import copy| & 設定辞書や payload を安全に複製するため。 このファイル内での主な使用位置は L584。 \\
5 & \path|import gzip| & detail CSV などの圧縮出力を行うため。 このファイル内での主な使用位置は L563, L564。 \\
6 & \path|import hashlib| & snapshot ID や入力資産の digest を作るため。 このファイル内での主な使用位置は L238, L252, L624。 \\
7 & \path|import html| & HTML report の文字列を安全に埋め込むため。 このファイル内での主な使用位置は L1131, L1136, L1206, L1210, L1214, L1222, L1253, L1254。 \\
8 & \path|import json| & manifest、checkpoint、UDP payload をやり取りするため。 このファイル内での主な使用位置は L253, L1206, L2830, L2950, L4001, L4005。 \\
9 & \path|import math| & Python標準のスカラー数値関数と定数を使うため。実際に参照する名前と行は使用位置欄で確認する。 このファイル内での主な使用位置は L134, L138, L885, L886, L887, L888, L889, L890, ...。 \\
10 & \path|import time| & wall/monotonic time に基づく周期制御や freshness 判定を行うため。 このファイル内での主な使用位置は L2055, L3713。 \\
11 & \path|import os| & パス、環境変数、プロセス外部状態を扱うため。 このファイル内での主な使用位置は L23, L30, L544, L546, L626, L639, L739, L742, ...。 \\
12 & \path|import subprocess| & git 情報取得や外部コマンド実行を行うため。 このファイル内での主な使用位置は L286, L299, L308, L317。 \\
13 & \path|import sys| & sys モジュールを利用するため。 このファイル内での主な使用位置は L37, L38, L4028。 \\
14 & \path|import traceback| & 例外時に crash log を残すため。 このファイル内での主な使用位置は L4025。 \\
15 & \path|from collections import deque| & 固定長の時系列や遅延キューを効率よく保持するため。 このファイル内での主な使用位置は L2075。 \\
16 & \path|from pathlib import Path| & ファイルやディレクトリを安全に扱うため。 このファイル内での主な使用位置は L36, L235, L270, L278, L810, L813, L835, L901, ...。 \\
17 & \path|from zoneinfo import ZoneInfo| & zoneinfo から ZoneInfo を読み込み、このファイルの処理を組み立てるため。 このファイル内での主な使用位置は L393, L2030, L2032, L3736, L3756。 \\
18 & \path|from datetime import datetime, timedelta, timezone| & UTC 時刻や相対時間を扱うため。 このファイル内での主な使用位置は L476, L626, L628, L1046, L1048, L1049, L1064, L1364, ...。 \\
32 & \path|import numpy as np| & 数値配列、clip、補間、統計計算を行うため。 このファイル内での主な使用位置は L78, L93, L100, L135, L433, L470, L484, L486, ...。 \\
33 & \path|import pandas as pd| & CSV や時刻列の読込・整形・再サンプリングに使うため。 このファイル内での主な使用位置は L332, L369, L373, L374, L383, L384, L387, L410, ...。 \\
34 & \path|import yaml| & profile、stop、schedule、summary YAML を読み書きするため。 このファイル内での主な使用位置は L68, L271, L612, L623, L3785。 \\
40 & \path|from mpc_solarcar.model import SolarCarModel, Params| & 車体物理・電気モデル本体 から SolarCarModel, Params を読み込み、このファイルの内部処理を分担させるため。 実体ファイルは mpc\_solarcar/model.py。 このファイル内での主な使用位置は L1689, L1703。 \\
41 & \path|from mpc_solarcar.route_utils import average_profile, interpolate_profile| & route\_utils.py から average\_profile, interpolate\_profile を読み込み、このファイルの内部処理を分担させるため。 実体ファイルは mpc\_solarcar/route\_utils.py。 このファイル内での主な使用位置は L1035, L2460, L2656。 \\
42 & \path|from mpc_solarcar.schedule_utils import DriveSchedule| & schedule\_utils.py から DriveSchedule を読み込み、このファイルの内部処理を分担させるため。 実体ファイルは mpc\_solarcar/schedule\_utils.py。 このファイル内での主な使用位置は L1664。 \\
43 & \path|from mpc_solarcar.signal_utils import SmoothRateLimiter| & signal\_utils.py から SmoothRateLimiter を読み込み、このファイルの内部処理を分担させるため。 実体ファイルは mpc\_solarcar/signal\_utils.py。 このファイル内での主な使用位置は L2078。 \\
44 & \path|from mpc_solarcar.solar_profile import get_path, get_section, load_profile as load_workflow_profile| & profile YAML 読込と検証 から get\_path, get\_section, load\_profile as load\_workflow\_profile を読み込み、このファイルの内部処理を分担させるため。 実体ファイルは mpc\_solarcar/solar\_profile.py。 このファイル内での主な使用位置は L649, L650, L651, L659, L660, L661, L662, L663, ...。 \\
45 & \path|from mpc_solarcar.upper_cost import active_upper_cost_terms, load_upper_cost_config, quad_penalty, upper_stage_cost, upper_terminal_cost| & 上位MPC 目的関数 から active\_upper\_cost\_terms, load\_upper\_cost\_config, quad\_penalty, upper\_stage\_cost, upper\_terminal\_cost を読み込み、このファイルの内部処理を分担させるため。 実体ファイルは mpc\_solarcar/upper\_cost.py。 このファイル内での主な使用位置は L1350, L1360, L1362, L1366, L1370, L1376, L1378, L1379, ...。 \\
52 & \path|from mpc_solarcar.upper_horizon import build_upper_distance_horizon, plan_segment_index| & 上位MPC 距離メッシュ生成 から build\_upper\_distance\_horizon, plan\_segment\_index を読み込み、このファイルの内部処理を分担させるため。 実体ファイルは mpc\_solarcar/upper\_horizon.py。 このファイル内での主な使用位置は L2341, L3431, L3470, L3590。 \\
53 & \path|from mpc_solarcar.upper_policy import absolute_control_distances, shift_upper_policy_warm_start| & 上位速度計画の補間と warm start から absolute\_control\_distances, shift\_upper\_policy\_warm\_start を読み込み、このファイルの内部処理を分担させるため。 実体ファイルは mpc\_solarcar/upper\_policy.py。 このファイル内での主な使用位置は L2366, L2374。 \\
57 & \path|from mpc_solarcar.upper_solver import hybrid_bounded_minimize| & 上位探索ソルバ から hybrid\_bounded\_minimize を読み込み、このファイルの内部処理を分担させるため。 実体ファイルは mpc\_solarcar/upper\_solver.py。 このファイル内での主な使用位置は L2873。 \\
1394 & \path|from scipy.optimize import minimize| & 目的関数と制約・boundsに基づく連続数値最適化を解くため。 このファイル内での主な使用位置は L1396。 \\
            \bottomrule
            \end{longtable}

        \section{このファイルを読む前に必要な基礎知識}
        次の章は、構文やROS用語を既知と仮定しないための説明である。

        \subsection{プログラム、プロセス、メモリ、オブジェクトを区別する}
ソースファイルはディスク上の文字列であり、それ自体は走っていない。Python実行可能プログラムがソースを読み、OSがその実行を一つのプロセスとして管理する。プロセスには仮想メモリ、開いているファイル、スレッド、終了コードなどが対応する。


\begin{lstlisting}
ソースファイル -> Pythonインタプリタが読み込む -> OS上のプロセス -> Pythonオブジェクトをメモリに生成 -> 関数やcallbackを実行
\end{lstlisting}


「メモリ上に生成する」とは、実行中プロセスが使う記憶領域に、その値の型、属性、参照関係を表す実体を用意することである。変数は実体そのものというより、そのオブジェクトを指す名前として理解するとPythonの代入が読みやすい。


\begin{lstlisting}[language=python]
node = MPCNode()
alias = node
# nodeとaliasは同じオブジェクトを参照する。
\end{lstlisting}


一つのプロセスには複数スレッドを持たせられる。スレッドは同じプロセスのメモリを共有するため、受信callbackと最適化callbackが同じself.zを書き換える場合は実行順と排他を検討する必要がある。

\paragraph{根拠資料}
\begin{itemize}[leftmargin=1.6em]
\item \href{https://docs.python.org/3/tutorial/classes.html}{Python公式チュートリアル: Classes}
\end{itemize}

\subsection{名前、代入、型、参照}
Pythonの代入文は、右辺を先に評価し、その結果のオブジェクトへ左辺の名前を結び付ける。`x = f()`では、まずfを呼び、その戻り値が得られてからxが更新される。


`float(...)`、`int(...)`、`bool(...)`は型名を呼び出して値を変換する。外部CSV、YAML、ROS parameterから来た値は型が期待どおりとは限らないため、このリポジトリでは明示変換が多い。


`None`は値が無いことを示す単一のオブジェクト、`math.nan`は浮動小数点値ではあるが有効な数値ではないことを示す。両者は用途が異なるため、`is None`と`math.isfinite`を使い分ける。

\paragraph{根拠資料}
\begin{itemize}[leftmargin=1.6em]
\item \href{https://docs.python.org/3/tutorial/classes.html}{Python公式チュートリアル: Classes}
\end{itemize}

\subsection{関数、引数、戻り値、スコープ}
`def`は関数本体をその場で実行する文ではない。関数オブジェクトを作り、その名前を現在の名前空間へ登録する。`f`は関数そのもの、`f()`はその関数を呼んだ結果である。


仮引数は関数定義側の受け取り名、実引数は呼出側から渡す値である。位置引数、キーワード引数、既定値付き引数、`*args`、`**kwargs`は渡し方が異なる。


\begin{lstlisting}[language=python]
def clip_speed(value, minimum=0.0, maximum=110.0):
    return min(maximum, max(minimum, value))

v = clip_speed(120.0, maximum=90.0)
\end{lstlisting}


関数内で作った通常の名前はローカルスコープに属する。メソッドが`self.z`を更新すればオブジェクトの状態が残るが、単に`z`を更新しただけならその呼出しのローカル値である。

\paragraph{根拠資料}
\begin{itemize}[leftmargin=1.6em]
\item \href{https://docs.python.org/3/tutorial/classes.html}{Python公式チュートリアル: Classes}
\end{itemize}

\subsection{module、package、import、console\_scripts}
Pythonファイルはmoduleとして読み込める。複数moduleをディレクトリにまとめたものがpackageである。`from .model import SolarCarModel`の先頭の点は、現在と同じpackage内のmodel moduleを指す相対importである。


import時にはファイルのトップレベル文が上から一度実行される。`def`や`class`は関数・クラスを登録するが、その本体の通常処理は呼び出すまで走らない。


setuptoolsの`console\_scripts`は、端末で使う実行可能名と`package.module:function`を対応付けるインストール時メタデータである。生成される実行用ラッパーはmoduleをimportし、指定関数を呼び、戻り値を終了コードとして扱う小さな入口である。


\begin{lstlisting}
ROS launchのexecutable名 -> install済みconsole script -> mpc_solarcar.mpc_nodeをimport -> main()を呼ぶ
\end{lstlisting}

\paragraph{根拠資料}
\begin{itemize}[leftmargin=1.6em]
\item \href{https://setuptools.pypa.io/en/latest/userguide/entry_point.html}{setuptools公式: Entry Points / Console Scripts}
\item \href{https://docs.python.org/3/tutorial/classes.html}{Python公式チュートリアル: Classes}
\end{itemize}

\subsection{例外、try/finally、with、資源解放}
例外は通常の戻り値とは別経路で異常を呼出元へ伝える。`try/except`は想定した異常を処理し、`finally`は成功・失敗にかかわらず後始末を行う。


`with open(...) as f:`はcontext managerを使い、ブロックを出るとファイルを閉じる。CSVやログの破損を避けるため、開いた資源の所有者と閉じる場所を明確にする。


`except Exception: pass`はノードを止めない利点がある一方、入力異常を隠して原因追跡を難しくする。安全に関係する値では、少なくとも頻度制限付き警告、異常カウンタ、fallback状態のpublishを検討する。



\subsection{class、独自クラス、インスタンス、self、継承}
クラスは、保持するデータと、そのデータを扱う関数を一つの型としてまとめる設計単位である。「独自クラス」とはPythonやROS 2が最初から用意した型ではなく、このプロジェクトが目的に合わせて定義した新しい型を指す。


`class MPCNode(Node)`は、rclpyのNodeを基底クラスとして継承し、Nodeが持つpublisher、subscription、timer、parameterなどの機能にMPC固有機能を追加する。丸括弧内は関数引数ではなく基底クラス指定である。


`MPCNode()`はクラスオブジェクトを呼び出して新しいインスタンスを作る。Pythonは新しい実体を用意し、その実体を第1引数selfとして`\_\_init\_\_`へ渡す。`self`は予約語ではなく慣習的な名前だが、この慣習を守ることでコードの意味が共有される。


\begin{lstlisting}[language=python]
class Counter:
    def __init__(self):
        self.value = 0

    def increment(self):
        self.value += 1

counter = Counter()
counter.increment()
\end{lstlisting}


`super().\_\_init\_\_(...)`は継承元の初期化処理を呼ぶ。MPCNodeではこれを省くとROSノードとして必要な内部実体が作られず、`create\_publisher`などを正しく使えない。

\paragraph{根拠資料}
\begin{itemize}[leftmargin=1.6em]
\item \href{https://docs.python.org/3/tutorial/classes.html}{Python公式チュートリアル: Classes}
\item \href{https://docs.ros.org/en/humble/Tutorials/Beginner-CLI-Tools/Understanding-ROS2-Nodes/Understanding-ROS2-Nodes.html}{ROS 2 Humble公式: Understanding nodes}
\end{itemize}

\subsection{先頭アンダースコア、\_\_init\_\_、名前付けの作法}
`self.\_step\_solar`の先頭1個のアンダースコアは「クラス内部で使う実装詳細」という慣習であり、アクセスを強制禁止する機構ではない。外部コードから呼べるが、公開APIとして依存しない意思表示である。


`\_\_init\_\_`の前後2個のアンダースコアはPythonが定めた特殊メソッド名である。クラス生成時に自動で呼ばれる。任意の名前に前後2個を付けて独自仕様を作ることは避ける。


`\_`で始まるローカル変数は、未使用または外部へ見せない意図を示す場合がある。例えば`\_base\_csv`は戻り値として受け取る必要はあるが、その後の処理では使わないことを読者へ示す。

\paragraph{根拠資料}
\begin{itemize}[leftmargin=1.6em]
\item \href{https://docs.python.org/3/tutorial/classes.html}{Python公式チュートリアル: Classes}
\end{itemize}

\subsection{list、dict、deque、NumPy配列、pandas DataFrame}
listは順序付き可変列、tupleは順序付きで通常変更しない列、dictはキーから値を引く対応表、dequeは両端追加・削除が効率的なキューである。どの構造を選ぶかはアクセス方法と更新方法で決まる。


NumPy配列は同種数値を連続的に扱い、要素ごとの演算、clip、補間、線形代数を簡潔に書く。shapeは各次元の要素数であり、速度系列なら通常`(N,)`、候補集団なら`(population, N)`となる。


pandas DataFrameは列名を持つ表である。CSV読込後は列型、欠損、単位、timezone、並び順、重複を明示的に処理しなければ、数値計算が動いても意味が誤る。



\subsection{CLI、PowerShell、Bash、環境変数、終了コード}
CLIは端末からプログラム名と引数を渡す操作界面である。`argparse`は文字列として届く引数を名前、型、既定値、必須性に従って解析する。


PowerShellとBashは別のshellであり、変数記法、改行継続、引用、パス表記が異なる。このプロジェクトではWindows側のSolarSim.ps1がWSL側のsolar\_control.shへ処理を渡す。


環境変数は親プロセスから子プロセスへ受け渡される名前付き文字列である。ROS\_DOMAIN\_ID、RMW\_IMPLEMENTATION、Pythonの数値スレッド数などはコード外から動作を変えるため、実行記録へ残す必要がある。


終了コード0は一般に成功、0以外は失敗を示す。shellルータは子プロセスの終了コードを握り潰さず上位へ返すことで、自動運用が失敗を検知できる。

\paragraph{根拠資料}
\begin{itemize}[leftmargin=1.6em]
\item \href{https://setuptools.pypa.io/en/latest/userguide/entry_point.html}{setuptools公式: Entry Points / Console Scripts}
\end{itemize}

\subsection{目的関数、制約、L-BFGS-B、SHGO、有限grid証明}
数値最適化器は、利用者が与えた目的関数を複数の候補点で評価し、より小さい値を持つ候補を探す。solverが物理を理解するのではなく、物理と運用価値はcost関数へ書かれる。


L-BFGS-Bは変数ごとの上下限を扱える局所最適化法である。初期値の近くの谷へ収束し得るため、非凸問題では複数seedや大域探索と組み合わせる。successがFalseでも有限な候補が返る場合があるため、採用条件をコード側で決める。


SHGOは定めたsamplingと局所最適化を組み合わせる大域最適化法である。有限Cartesian gridの全列挙は、そのgrid上の最良を証明できるが、連続領域全体の最良を自動的に証明しない。資料ではこの証明範囲を区別する。

\paragraph{根拠資料}
\begin{itemize}[leftmargin=1.6em]
\item \href{https://docs.scipy.org/doc/scipy/reference/generated/scipy.optimize.minimize.html}{SciPy公式: scipy.optimize.minimize}
\item \href{https://docs.scipy.org/doc/scipy/reference/generated/scipy.optimize.shgo.html}{SciPy公式: scipy.optimize.shgo}
\end{itemize}

\subsection{Cross-Entropy Methodを式と実装で理解する}
CEMは候補を生成する確率分布を持ち、良かったelite候補から分布を更新する反復的な確率最適化である。このリポジトリのupper\_solver.pyは各制御点速度を独立正規分布で生成し、上下限へclipする。


\[
u_i^{(j)} \sim \mathcal{N}(\mu_i^{(g)},(\sigma_i^{(g)})^2),\qquad u_i^{(j)}\leftarrow\operatorname{clip}(u_i^{(j)},l_i,h_i)
\]

\[
\mu_i^{(g+1)}=\frac{1}{K}\sum_{j\in\mathcal{E}_g}u_i^{(j)},\qquad \sigma_i^{(g+1)}=\max\left(\operatorname{Std}_{j\in\mathcal{E}_g}u_i^{(j)},0.05(h_i-l_i)\right)
\]

ここでE\_gはcostが小さい上位K候補である。平均は良い領域へ移り、標準偏差は探索幅を表す。標準偏差の下限は探索が完全に潰れることを避ける。


現行hybrid\_bounded\_minimizeは、deterministic seedを評価し、上位候補をL-BFGS-Bで局所refineし、設定とseed間不一致に応じてCEMを実行し、最後に再度局所refineする。したがってCEM単独ではなくhybrid solverである。


CEMで落とした候補を永久保存しないこと自体は通常の最適化として自然だが、off-nominal状態からの再利用には別のpolicy library設計が必要である。状態を無制限に全組合せ保存する代わりに、SoC、進捗、時刻、予報誤差、停止状態などのscenarioを設計し、近傍policyを検索してMPCで再最適化する。

\paragraph{根拠資料}
\begin{itemize}[leftmargin=1.6em]
\item \href{https://link.springer.com/book/10.1007/978-1-4757-4321-0}{Rubinstein and Kroese: The Cross-Entropy Method}
\item \href{https://docs.scipy.org/doc/scipy/reference/generated/scipy.optimize.minimize.html}{SciPy公式: scipy.optimize.minimize}
\end{itemize}

\subsection{warm startは何を保存し、どう効くか}
warm startは前回またはoffline探索で得た入力系列を、次の最適化の初期候補として渡すことである。答えを固定するのではなく、探索開始点と候補libraryの一員を与える。


\[
u^{0,\mathrm{new}}_i=\operatorname{interp}\left(s_i^{\mathrm{new}};\,s_j^{\mathrm{old}},u_j^{\ast,\mathrm{old}}\right)
\]

距離基準計画では、現在より後ろの旧制御点を捨て、新しい絶対距離制御点へ補間してshiftする。初期状態や天候が変わればcost評価は新条件で行われるので、warm startが不適切でも他seed、CEM、安全fallbackが補う設計が必要である。


warm startの効き方は、局所法なら収束先と反復数、CEMなら初期meanまたは候補pool、receding horizonなら前回解の時間・距離shiftとして現れる。どの位置に渡しているかを呼出引数まで追う。

\paragraph{根拠資料}
\begin{itemize}[leftmargin=1.6em]
\item \href{https://docs.scipy.org/doc/scipy/reference/generated/scipy.optimize.minimize.html}{SciPy公式: scipy.optimize.minimize}
\item \href{https://link.springer.com/book/10.1007/978-1-4757-4321-0}{Rubinstein and Kroese: The Cross-Entropy Method}
\end{itemize}

\subsection{車両力学、電力収支、効率map}
\[
F_{\mathrm{aero}}=\frac{1}{2}\rho C_dA(v+v_{\mathrm{wind}})^2,\quad F_{\mathrm{roll}}\approx mgC_{rr}\cos\theta,\quad F_{\mathrm{grade}}=mg\sin\theta
\]

車輪機械powerは概ね`P\_mech = F\_total * v + P\_inertia`で、駆動時と回生時で異なる効率mapを通してDC側powerへ変換する。mapは速度・torqueなどを軸にした実測または同定tableであり、範囲外の補間・clip規則もモデルの一部である。


\[
P_{\mathrm{pack}}=P_{\mathrm{drive,dc}}-P_{\mathrm{regen,dc}}+P_{\mathrm{aux}}-P_{\mathrm{pv}}
\]

符号規約は必ずコードで確認する。本プロジェクトでは正のpack powerを放電側としてSoCを減らす処理が中心であり、発電と回生はpack負荷を下げる方向に働く。



\subsection{SoC、内部抵抗、端子電圧、温度、MHE}
\[
V=V_{\mathrm{oc}}(z,T)-IR_{\mathrm{total}}(z,T,I),\qquad z_{k+1}=z_k-\frac{\eta(I,T)I\Delta t}{Q_{\mathrm{eff}}}
\]

SoCは直接完全には観測できないため、電流積算、OCV、端子電圧、温度、容量、効率を組み合わせる。内部抵抗はSoC、温度、電流方向で変わり、発熱と電圧制約の両方へ影響する。


MHEは有限時間窓の状態と観測誤差をまとめて最適化し、SoCや温度を推定する。古い測定や欠損値を同じ重みで使うと推定が壊れるため、timestamp freshnessと観測可用性を入口で確認する。

\paragraph{根拠資料}
\begin{itemize}[leftmargin=1.6em]
\item \href{https://docs.scipy.org/doc/scipy/reference/generated/scipy.optimize.minimize.html}{SciPy公式: scipy.optimize.minimize}
\end{itemize}

\subsection{天候、route、補間、時刻、単位}
予報は時刻だけの系列か、時刻とroute距離の2次元gridになり得る。現在時刻tと距離sがgrid点の間なら、周囲の値から補間して日射、温度、風を求める。


UTC、現地timezone、naive datetimeを混ぜると数時間のずれが発生する。入力でtimezoneを確定し、内部比較はUTC、表示は現地時刻という役割分担が安全である。


route距離のkm、速度のkm/hとm/s、時間の秒、energyのWhとJを跨ぐため、変換係数3.6、3600、1000の意味を各式で確認する。



\subsection{freshness、filter、guard、fallback、fail-safe}
分散システムでは最後に受け取った値が現在も有効とは限らない。受信時刻とtimeoutからfreshnessを判定し、stale値を計画状態へ無条件に同期しない。


filterはnoiseと一時的な飛び値を抑えるが、遅れを生む。slew limitは指令変化率を制限する。安全guardはsolverのcost罰則とは別に、現在出力へ強制制約を適用する最後の防波堤である。


fallbackは失敗時の代替動作を事前に決める設計である。前回計画保持、物理に基づく決定論的入力、停止、低速制限などから、故障modeごとに選ぶ。fallback発生はstatusとlogへ残し、正常解と区別する。



\subsection{制御、状態、入力、モデル予測制御MPC}
制御対象の内部を表す状態をx、操作入力をu、外乱・予報をwとすると、離散モデルは`x[k+1] = f(x[k], u[k], w[k])`と書ける。ソーラーカーではSoC、電池温度、距離、速度などが状態候補、速度目標や駆動トルクが入力候補になる。


\[
\min_{u_0,\ldots,u_{N-1}} \sum_{k=0}^{N-1}\ell(x_k,u_k,w_k)+V_f(x_N)
\]

MPCは現在状態からNステップ先まで予測し、目的関数と制約を満たす入力系列を求める。ただし実際に適用するのは通常先頭入力だけで、次回は新しい実測状態から再び解く。これがreceding horizonである。


予測モデル、目的関数、制約、ホライズン、solver、初期値のどれかが変わると答えも変わる。「MPCを使う」だけでは仕様は決まらず、これらを単位付きで追う必要がある。

\paragraph{根拠資料}
\begin{itemize}[leftmargin=1.6em]
\item \href{https://docs.scipy.org/doc/scipy/reference/generated/scipy.optimize.minimize.html}{SciPy公式: scipy.optimize.minimize}
\end{itemize}

\subsection{上位MPCと下位MPCの役割}
上位層は長い距離または時間を見て、エネルギー、到着、停止、速度制限、終端SoCを考えた速度計画を出す。下位層は短い周期で実測速度を見て、上位速度へ追従する駆動・回生入力を求める。


\begin{lstlisting}
予報・route・SoC -> 上位MPC -> 将来速度列 -> 下位MPC -> throttle/回生/drive mode -> driverまたはvehicle -> 新しいtelemetry
\end{lstlisting}


上位解が遅い間も下位出力を止めないこと、古い計画を安全に保持すること、上位と下位で単位と時刻基準を一致させることが実装上重要である。



\subsection{CSV/YAMLのdata contractと検証}
data contractは列名、型、単位、timezone、欠損可否、並び順、重複、許容範囲、先頭行、encodingを事前に決めた仕様である。単にCSVとして読めることは、モデル入力として正しいことを意味しない。


同定用実測、map、route、forecast、stop、scheduleはそれぞれgrainが異なる。生成時にschema validation、物理範囲、時間単調性、route範囲、coverageを検査し、検査結果をartifactとして残す。


学習用データと独立検証データを分離し、RMSEだけでなくbias、時系列残差、energy積算誤差、終端SoC、温度・電圧制約、外挿領域を評価する。



\subsection{単体試験、SILS、replay、観測可能性}
単体試験は関数やモデル項の局所契約、SILSは複数Nodeと時間進行を含む閉ループ、historical replayは実測入力への再現性、本番preflightは実機通信と停止動作を確認する。目的が異なるため一つで代用しない。


再現可能な診断には、入力、出力、内部状態、solver成否、候補数、計算時間、fallback、source revision、profile hashを同じrun IDで記録する。


非同期不具合は最終値だけでは追えない。topic周期、message age、callback所要時間、queue遅延、Node生存、publisher数を同時に観測する。

\paragraph{根拠資料}
\begin{itemize}[leftmargin=1.6em]
\item \href{https://docs.ros.org/en/ros2_packages/humble/api/rqt_graph/index.html}{ROS 2 Humble公式: rqt\_graph}
\item \href{https://docs.ros.org/en/humble/Tutorials/Beginner-CLI-Tools/Recording-And-Playing-Back-Data/Recording-And-Playing-Back-Data.html}{ROS 2 Humble公式: Recording and playing back data}
\item \href{https://docs.ros.org/en/rolling/Concepts/Intermediate/About-Executors.html}{ROS 2公式: Executors}
\end{itemize}



        \section{関数・クラスを上から順に解説}
        \subsection{L63 関数 load\_yaml}
            \begin{itemize}[leftmargin=1.6em]
              \item 行範囲: L63--L70
              \item 所属: module直下
              \item 定義: \path|load_yaml(path)|
              \item docstring: 明示されていない。
              \item このブロックが直接呼ぶ主な関数・メソッド: open, safe\_load
              \item 戻り値の要点: \{\} / yaml.safe\_load(f) or \{\} / \{\}
              \item 主なローカル代入名: f
              \item 読み取る主なインスタンス属性: 特になし。
              \item 更新する主なインスタンス属性: 特になし。
              \item 明示的に送出する例外: 明示的raiseはない。
              \item 制御構造の規模: 条件分岐 1、ループ 0、try 1
            \end{itemize}
            \paragraph{この定義を読むためのPython構文}
            \begin{itemize}[leftmargin=1.6em]
            \item with文はcontext managerに開始・終了処理を任せ、ファイルなどの資源を確実に解放する。
\item try文は例外が起き得る処理と、異常時または終了時の経路を分ける。
            \end{itemize}
            \paragraph{上から順の処理}
            \begin{enumerate}[leftmargin=1.8em]
            \item 条件 not path を判定し、真なら内部処理を行う。
\item   \{\} を返す。
\item 例外処理を伴う try ブロックを実行する。
\item   with 文で open(path, 'r', encoding='utf-8') を管理しながら処理する。
\item     yaml.safe\_load(f) or \{\} を返す。
\item   Exceptionを捕捉した場合:
\item   \{\} を返す。
            \end{enumerate}
            \paragraph{代表コード断片}
            \begin{lstlisting}[language=Python]
def load_yaml(path):
    if not path:
        return {}
    try:
        with open(path, 'r', encoding='utf-8') as f:
            return yaml.safe_load(f) or {}
    except Exception:
        return {}
\end{lstlisting}

\subsection{L73 関数 sim\_log}
            \begin{itemize}[leftmargin=1.6em]
              \item 行範囲: L73--L74
              \item 所属: module直下
              \item 定義: \path|sim_log(message: str)|
              \item docstring: 明示されていない。
              \item このブロックが直接呼ぶ主な関数・メソッド: print
              \item 戻り値の要点: 戻り値が無いか、return 文が明示されていない。
              \item 主なローカル代入名: 特になし。
              \item 読み取る主なインスタンス属性: 特になし。
              \item 更新する主なインスタンス属性: 特になし。
              \item 明示的に送出する例外: 明示的raiseはない。
              \item 制御構造の規模: 条件分岐 0、ループ 0、try 0
            \end{itemize}
            \paragraph{この定義を読むためのPython構文}
            \begin{itemize}[leftmargin=1.6em]
            \item 特別な構文上の補足は少ない。
            \end{itemize}
            \paragraph{上から順の処理}
            \begin{enumerate}[leftmargin=1.8em]
            \item print(...) を実行する。
            \end{enumerate}
            \paragraph{代表コード断片}
            \begin{lstlisting}[language=Python]
def sim_log(message: str):
    print(f"[solar_sim] {message}", flush=True)
\end{lstlisting}

\subsection{L77 関数 select\_optimized\_vector}
            \begin{itemize}[leftmargin=1.6em]
              \item 行範囲: L77--L112
              \item 所属: module直下
              \item 定義: \path|select_optimized_vector(res, x0, *, label: str)|
              \item docstring: 明示されていない。
              \item このブロックが直接呼ぶ主な関数・メソッド: all, asarray, bool, getattr, isfinite, sim\_log, str, strip
              \item 戻り値の要点: x\_arr / x0\_arr / x0\_arr / x0\_arr
              \item 主なローカル代入名: message, raw\_x, status, success, x0\_arr, x\_arr
              \item 読み取る主なインスタンス属性: 特になし。
              \item 更新する主なインスタンス属性: 特になし。
              \item 明示的に送出する例外: 明示的raiseはない。
              \item 制御構造の規模: 条件分岐 4、ループ 0、try 1
            \end{itemize}
            \paragraph{この定義を読むためのPython構文}
            \begin{itemize}[leftmargin=1.6em]
            \item try文は例外が起き得る処理と、異常時または終了時の経路を分ける。
            \end{itemize}
            \paragraph{上から順の処理}
            \begin{enumerate}[leftmargin=1.8em]
            \item x0\_arr に np.asarray(x0, dtype=float) の結果を代入する。
\item 条件 res is None を判定し、真なら内部処理を行う。
\item   sim\_log(...) を実行する。
\item   x0\_arr を返す。
\item success に bool(getattr(res, 'success', False)) の結果を代入する。
\item status に getattr(res, 'status', None) の結果を代入する。
\item message に str(getattr(res, 'message', '') or '').strip() の結果を代入する。
\item raw\_x に getattr(res, 'x', None) の結果を代入する。
\item 条件 raw\_x is None を判定し、真なら内部処理を行う。
\item   sim\_log(...) を実行する。
\item   x0\_arr を返す。
\item 例外処理を伴う try ブロックを実行する。
\item   x\_arr に np.asarray(raw\_x, dtype=float) の結果を代入する。
\item   Exceptionを捕捉した場合:
\item   sim\_log(...) を実行する。
\item   x0\_arr を返す。
\item 条件 x\_arr.shape != x0\_arr.shape or not np.all(np.isfinite(x\_arr)) を判定し、真なら内部処理を行う。
\item   sim\_log(...) を実行する。
\item   x0\_arr を返す。
\item 条件 not success を判定し、真なら内部処理を行う。
\item   sim\_log(...) を実行する。
\item   x0\_arr を返す。
\item x\_arr を返す。
            \end{enumerate}
            \paragraph{代表コード断片}
            \begin{lstlisting}[language=Python]
def select_optimized_vector(res, x0, *, label: str):
    x0_arr = np.asarray(x0, dtype=float)
    if res is None:
        sim_log(f"{label} fallback: optimizer returned no result; using the initial guess.")
        return x0_arr
    success = bool(getattr(res, 'success', False))
    status = getattr(res, 'status', None)
    message = str(getattr(res, 'message', '') or '').strip()
    raw_x = getattr(res, 'x', None)
    if raw_x is None:
        sim_log(
            f"{label} fallback: optimizer produced no candidate "
            f"(success={success} status={status} message={message!r}); using the initial guess."
        )
        return x0_arr
    try:
        x_arr = np.asarray(raw_x, dtype=float)
    except Exception:
        sim_log(
            f"{label} fallback: optimizer candidate could not be converted to float "
            f"(success={success} status={status} message={message!r}); using the initial guess."
        )
        return x0_arr
    if x_arr.shape != x0_arr.shape or not np.all(np.isfinite(x_arr)):
        sim_log(
            f"{label} fallback: optimizer candidate was non-finite or mismatched "
            f"(success={success} status={status} message={message!r}); using the initial guess."
        )
        return x0_arr
    if not success:
        sim_log(
            f"{label} fallback: optimizer reported failure "
            f"(status={status} message={message!r}); using the initial guess."
        )
        return x0_arr
...
\end{lstlisting}

\subsection{L115 関数 limit\_step\_duration\_to\_distance}
            \begin{itemize}[leftmargin=1.6em]
              \item 行範囲: L115--L126
              \item 所属: module直下
              \item 定義: \path!limit_step_duration_to_distance(step_sec: float, speed_kmh: float, s0_km: float, race_km: float | None) -> float!
              \item docstring: 明示されていない。
              \item このブロックが直接呼ぶ主な関数・メソッド: float, max, min
              \item 戻り値の要点: min(step\_sec, max\_step\_sec) / step\_sec / 0.0 / step\_sec
              \item 主なローカル代入名: max\_step\_sec, remaining\_km, speed\_kmh, step\_sec
              \item 読み取る主なインスタンス属性: 特になし。
              \item 更新する主なインスタンス属性: 特になし。
              \item 明示的に送出する例外: 明示的raiseはない。
              \item 制御構造の規模: 条件分岐 3、ループ 0、try 0
            \end{itemize}
            \paragraph{この定義を読むためのPython構文}
            \begin{itemize}[leftmargin=1.6em]
            \item `->`以後は戻り値の型注釈であり、実行時の値を自動変換する命令ではない。
            \end{itemize}
            \paragraph{上から順の処理}
            \begin{enumerate}[leftmargin=1.8em]
            \item step\_sec に max(0.0, float(step\_sec)) の結果を代入する。
\item 条件 race\_km is None を判定し、真なら内部処理を行う。
\item   step\_sec を返す。
\item remaining\_km に max(0.0, float(race\_km) - float(s0\_km)) の結果を代入する。
\item 条件 remaining\_km <= DISTANCE\_EPS\_KM を判定し、真なら内部処理を行う。
\item   0.0 を返す。
\item speed\_kmh に max(0.0, float(speed\_kmh)) の結果を代入する。
\item 条件 speed\_kmh <= 1e-09 を判定し、真なら内部処理を行う。
\item   step\_sec を返す。
\item max\_step\_sec に remaining\_km * 3600.0 / speed\_kmh の結果を代入する。
\item min(step\_sec, max\_step\_sec) を返す。
            \end{enumerate}
            \paragraph{代表コード断片}
            \begin{lstlisting}[language=Python]
def limit_step_duration_to_distance(step_sec: float, speed_kmh: float, s0_km: float, race_km: float | None) -> float:
    step_sec = max(0.0, float(step_sec))
    if race_km is None:
        return step_sec
    remaining_km = max(0.0, float(race_km) - float(s0_km))
    if remaining_km <= DISTANCE_EPS_KM:
        return 0.0
    speed_kmh = max(0.0, float(speed_kmh))
    if speed_kmh <= 1.0e-9:
        return step_sec
    max_step_sec = remaining_km * 3600.0 / speed_kmh
    return min(step_sec, max_step_sec)
\end{lstlisting}

\subsection{L129 関数 terminal\_soc\_predictions}
            \begin{itemize}[leftmargin=1.6em]
              \item 行範囲: L129--L142
              \item 所属: module直下
              \item 定義: \path|terminal_soc_predictions(upper_solve_log: list[dict]) -> tuple[float, float]|
              \item docstring: Return initial full-horizon and latest nontrivial terminal-SoC predictions.
              \item このブロックが直接呼ぶ主な関数・メソッド: append, dict, float, get, int, isfinite
              \item 戻り値の要点: (float(initial\_soc), float(latest\_nontrivial\_soc)) / (math.nan, math.nan)
              \item 主なローカル代入名: candidates, initial\_soc, latest\_nontrivial\_soc, nontrivial, prediction, row, steps, terminal\_soc
              \item 読み取る主なインスタンス属性: 特になし。
              \item 更新する主なインスタンス属性: 特になし。
              \item 明示的に送出する例外: 明示的raiseはない。
              \item 制御構造の規模: 条件分岐 2、ループ 1、try 0
            \end{itemize}
            \paragraph{この定義を読むためのPython構文}
            \begin{itemize}[leftmargin=1.6em]
            \item `->`以後は戻り値の型注釈であり、実行時の値を自動変換する命令ではない。
\item 内包表記は、反復・条件・値生成を一つの式で記述してcollectionを作る。
            \end{itemize}
            \paragraph{上から順の処理}
            \begin{enumerate}[leftmargin=1.8em]
            \item candidates に [] の結果を代入する。
\item upper\_solve\_log を順に走査し、各要素を row に入れて処理する。
\item   prediction に dict(row.get('selected\_prediction', \{\}) or \{\}) の結果を代入する。
\item   terminal\_soc に float(prediction.get('terminal\_soc', math.nan)) の結果を代入する。
\item   条件 np.isfinite(terminal\_soc) を判定し、真なら内部処理を行う。
\item     candidates.append(...) を実行する。
\item 条件 not candidates を判定し、真なら内部処理を行う。
\item   (math.nan, math.nan) を返す。
\item initial\_soc に candidates[0][1] の結果を代入する。
\item nontrivial に [terminal\_soc for steps, terminal\_soc in candidates if steps > 1] の結果を代入する。
\item latest\_nontrivial\_soc に nontrivial[-1] if nontrivial else candidates[-1][1] の結果を代入する。
\item (float(initial\_soc), float(latest\_nontrivial\_soc)) を返す。
            \end{enumerate}
            \paragraph{代表コード断片}
            \begin{lstlisting}[language=Python]
def terminal_soc_predictions(upper_solve_log: list[dict]) -> tuple[float, float]:
    """Return initial full-horizon and latest nontrivial terminal-SoC predictions."""
    candidates = []
    for row in upper_solve_log:
        prediction = dict(row.get('selected_prediction', {}) or {})
        terminal_soc = float(prediction.get('terminal_soc', math.nan))
        if np.isfinite(terminal_soc):
            candidates.append((int(row.get('prediction_steps', 0) or 0), terminal_soc))
    if not candidates:
        return math.nan, math.nan
    initial_soc = candidates[0][1]
    nontrivial = [terminal_soc for steps, terminal_soc in candidates if steps > 1]
    latest_nontrivial_soc = nontrivial[-1] if nontrivial else candidates[-1][1]
    return float(initial_soc), float(latest_nontrivial_soc)
\end{lstlisting}

\subsection{L145 関数 advance\_rate\_limiter\_to\_distance\_boundary}
            \begin{itemize}[leftmargin=1.6em]
              \item 行範囲: L145--L192
              \item 所属: module直下
              \item 定義: \path!advance_rate_limiter_to_distance_boundary(limiter, target_kmh: float, *, start_time_sec: float, step_sec: float, s0_km: float, distance_limit_km: float | None) -> tuple[float, float]!
              \item docstring: Advance a limiter by the actual substep, including a short boundary remainder.
              \item このブロックが直接呼ぶ主な関数・メソッド: abs, float, limit\_step\_duration\_to\_distance, max, range, update
              \item 戻り値の要点: (speed\_kmh, actual\_dt) / (float(limiter.value), 0.0) / (speed\_kmh, next\_dt)
              \item 主なローカル代入名: \_, actual\_dt, initial\_value, limited\_dt, next\_dt, requested\_dt, speed\_kmh
              \item 読み取る主なインスタンス属性: 特になし。
              \item 更新する主なインスタンス属性: 特になし。
              \item 明示的に送出する例外: 明示的raiseはない。
              \item 制御構造の規模: 条件分岐 2、ループ 1、try 0
            \end{itemize}
            \paragraph{この定義を読むためのPython構文}
            \begin{itemize}[leftmargin=1.6em]
            \item `->`以後は戻り値の型注釈であり、実行時の値を自動変換する命令ではない。
            \end{itemize}
            \paragraph{上から順の処理}
            \begin{enumerate}[leftmargin=1.8em]
            \item requested\_dt に max(0.0, float(step\_sec)) の結果を代入する。
\item 条件 requested\_dt <= 0.0 を判定し、真なら内部処理を行う。
\item   (float(limiter.value), 0.0) を返す。
\item initial\_value に float(limiter.value) の結果を代入する。
\item actual\_dt に requested\_dt の結果を代入する。
\item speed\_kmh に initial\_value の結果を代入する。
\item range(8) を順に走査し、各要素を \_ に入れて処理する。
\item   limiter.value に initial\_value の結果を代入する。
\item   limiter.last\_time に float(start\_time\_sec) の結果を代入する。
\item   speed\_kmh に float(limiter.update(target\_kmh, now=float(start\_time\_sec) + actual\_dt)) の結果を代入する。
\item   limited\_dt に limit\_step\_duration\_to\_distance(requested\_dt, speed\_kmh, s0\_km, distance\_limit\_km) の結果を代入する。
\item   next\_dt に max(0.0, float(limited\_dt)) の結果を代入する。
\item   条件 abs(next\_dt - actual\_dt) <= 1e-09 を判定し、真なら内部処理を行う。
\item     (speed\_kmh, next\_dt) を返す。
\item   actual\_dt に next\_dt の結果を代入する。
\item limiter.value に initial\_value の結果を代入する。
\item limiter.last\_time に float(start\_time\_sec) の結果を代入する。
\item speed\_kmh に float(limiter.update(target\_kmh, now=float(start\_time\_sec) + actual\_dt)) の結果を代入する。
\item (speed\_kmh, actual\_dt) を返す。
            \end{enumerate}
            \paragraph{代表コード断片}
            \begin{lstlisting}[language=Python]
def advance_rate_limiter_to_distance_boundary(
    limiter,
    target_kmh: float,
    *,
    start_time_sec: float,
    step_sec: float,
    s0_km: float,
    distance_limit_km: float | None,
) -> tuple[float, float]:
    """Advance a limiter by the actual substep, including a short boundary remainder."""
    requested_dt = max(0.0, float(step_sec))
    if requested_dt <= 0.0:
        return float(limiter.value), 0.0

    initial_value = float(limiter.value)
    actual_dt = requested_dt
    speed_kmh = initial_value
    for _ in range(8):
        limiter.value = initial_value
        # The stored value is the state at the beginning of this substep.  Do
        # not let a stale limiter timestamp add time to a short boundary step.
        limiter.last_time = float(start_time_sec)
        speed_kmh = float(
            limiter.update(
                target_kmh,
                now=float(start_time_sec) + actual_dt,
            )
        )
        limited_dt = limit_step_duration_to_distance(
            requested_dt,
            speed_kmh,
            s0_km,
            distance_limit_km,
        )
        next_dt = max(0.0, float(limited_dt))
...
\end{lstlisting}

\subsection{L195 関数 snap\_execution\_stop\_speed\_kmh}
            \begin{itemize}[leftmargin=1.6em]
              \item 行範囲: L195--L209
              \item 所属: module直下
              \item 定義: \path|snap_execution_stop_speed_kmh(speed_kmh: float, target_kmh: float, *, stop_requested: bool, deadband_kmh: float, quantize_step_kmh: float) -> tuple[float, bool]|
              \item docstring: Remove the final quantized crawl once a stop command has nearly settled.
              \item このブロックが直接呼ぶ主な関数・メソッド: bool, float, max
              \item 戻り値の要点: (speed, False) / (0.0, True)
              \item 主なローカル代入名: speed, target, threshold
              \item 読み取る主なインスタンス属性: 特になし。
              \item 更新する主なインスタンス属性: 特になし。
              \item 明示的に送出する例外: 明示的raiseはない。
              \item 制御構造の規模: 条件分岐 1、ループ 0、try 0
            \end{itemize}
            \paragraph{この定義を読むためのPython構文}
            \begin{itemize}[leftmargin=1.6em]
            \item `->`以後は戻り値の型注釈であり、実行時の値を自動変換する命令ではない。
            \end{itemize}
            \paragraph{上から順の処理}
            \begin{enumerate}[leftmargin=1.8em]
            \item speed に max(0.0, float(speed\_kmh)) の結果を代入する。
\item target に max(0.0, float(target\_kmh)) の結果を代入する。
\item threshold に max(1e-06, float(deadband\_kmh), float(quantize\_step\_kmh)) の結果を代入する。
\item 条件 bool(stop\_requested) and target <= 1e-09 and (speed <= threshold + 1e-09) を判定し、真なら内部処理を行う。
\item   (0.0, True) を返す。
\item (speed, False) を返す。
            \end{enumerate}
            \paragraph{代表コード断片}
            \begin{lstlisting}[language=Python]
def snap_execution_stop_speed_kmh(
    speed_kmh: float,
    target_kmh: float,
    *,
    stop_requested: bool,
    deadband_kmh: float,
    quantize_step_kmh: float,
) -> tuple[float, bool]:
    """Remove the final quantized crawl once a stop command has nearly settled."""
    speed = max(0.0, float(speed_kmh))
    target = max(0.0, float(target_kmh))
    threshold = max(1.0e-6, float(deadband_kmh), float(quantize_step_kmh))
    if bool(stop_requested) and target <= 1.0e-9 and speed <= threshold + 1.0e-9:
        return 0.0, True
    return speed, False
\end{lstlisting}

\subsection{L212 関数 choose\_integration\_step\_seconds}
            \begin{itemize}[leftmargin=1.6em]
              \item 行範囲: L212--L215
              \item 所属: module直下
              \item 定義: \path|choose_integration_step_seconds(simulation_step_sec: float, weather_step_sec: float) -> float|
              \item docstring: 明示されていない。
              \item このブロックが直接呼ぶ主な関数・メソッド: float, max, min
              \item 戻り値の要点: max(60.0, min(simulation\_step, weather\_step))
              \item 主なローカル代入名: simulation\_step, weather\_step
              \item 読み取る主なインスタンス属性: 特になし。
              \item 更新する主なインスタンス属性: 特になし。
              \item 明示的に送出する例外: 明示的raiseはない。
              \item 制御構造の規模: 条件分岐 0、ループ 0、try 0
            \end{itemize}
            \paragraph{この定義を読むためのPython構文}
            \begin{itemize}[leftmargin=1.6em]
            \item `->`以後は戻り値の型注釈であり、実行時の値を自動変換する命令ではない。
            \end{itemize}
            \paragraph{上から順の処理}
            \begin{enumerate}[leftmargin=1.8em]
            \item simulation\_step に float(max(simulation\_step\_sec, 1.0)) の結果を代入する。
\item weather\_step に float(max(weather\_step\_sec, 1.0)) の結果を代入する。
\item max(60.0, min(simulation\_step, weather\_step)) を返す。
            \end{enumerate}
            \paragraph{代表コード断片}
            \begin{lstlisting}[language=Python]
def choose_integration_step_seconds(simulation_step_sec: float, weather_step_sec: float) -> float:
    simulation_step = float(max(simulation_step_sec, 1.0))
    weather_step = float(max(weather_step_sec, 1.0))
    return max(60.0, min(simulation_step, weather_step))
\end{lstlisting}

\subsection{L218 関数 choose\_stationary\_integration\_step\_seconds}
            \begin{itemize}[leftmargin=1.6em]
              \item 行範囲: L218--L227
              \item 所属: module直下
              \item 定義: \path|choose_stationary_integration_step_seconds(simulation_step_sec: float, weather_step_sec: float, requested_step_sec: float = 60.0) -> float|
              \item docstring: Use a fine state step after weather interpolation during long stops.
              \item このブロックが直接呼ぶ主な関数・メソッド: float, max, min
              \item 戻り値の要点: min(requested, simulation\_step, weather\_step)
              \item 主なローカル代入名: requested, simulation\_step, weather\_step
              \item 読み取る主なインスタンス属性: 特になし。
              \item 更新する主なインスタンス属性: 特になし。
              \item 明示的に送出する例外: 明示的raiseはない。
              \item 制御構造の規模: 条件分岐 0、ループ 0、try 0
            \end{itemize}
            \paragraph{この定義を読むためのPython構文}
            \begin{itemize}[leftmargin=1.6em]
            \item `->`以後は戻り値の型注釈であり、実行時の値を自動変換する命令ではない。
            \end{itemize}
            \paragraph{上から順の処理}
            \begin{enumerate}[leftmargin=1.8em]
            \item requested に max(1.0, float(requested\_step\_sec)) の結果を代入する。
\item simulation\_step に max(1.0, float(simulation\_step\_sec)) の結果を代入する。
\item weather\_step に max(1.0, float(weather\_step\_sec)) の結果を代入する。
\item min(requested, simulation\_step, weather\_step) を返す。
            \end{enumerate}
            \paragraph{代表コード断片}
            \begin{lstlisting}[language=Python]
def choose_stationary_integration_step_seconds(
    simulation_step_sec: float,
    weather_step_sec: float,
    requested_step_sec: float = 60.0,
) -> float:
    """Use a fine state step after weather interpolation during long stops."""
    requested = max(1.0, float(requested_step_sec))
    simulation_step = max(1.0, float(simulation_step_sec))
    weather_step = max(1.0, float(weather_step_sec))
    return min(requested, simulation_step, weather_step)
\end{lstlisting}

\subsection{L230 関数 write\_model\_snapshot}
            \begin{itemize}[leftmargin=1.6em]
              \item 行範囲: L230--L275
              \item 所属: module直下
              \item 定義: \path|write_model_snapshot(detail_csv: str, parameters: dict, map_paths: dict) -> tuple[str, str]|
              \item docstring: Write immutable model provenance once instead of repeating map paths at 1 Hz.
              \item このブロックが直接呼ぶ主な関数・メソッド: Path, bool, dict, dumps, encode, endswith, ensure\_parent\_dir, hexdigest, int, is\_file, items, iter
              \item 戻り値の要点: (snapshot\_id, output)
              \item 主なローカル代入名: block, canonical, digest, entry, key, maps, output, path, payload, raw, raw\_base, raw\_path, snapshot\_id, stream
              \item 読み取る主なインスタンス属性: 特になし。
              \item 更新する主なインスタンス属性: 特になし。
              \item 明示的に送出する例外: 明示的raiseはない。
              \item 制御構造の規模: 条件分岐 3、ループ 2、try 0
            \end{itemize}
            \paragraph{この定義を読むためのPython構文}
            \begin{itemize}[leftmargin=1.6em]
            \item `->`以後は戻り値の型注釈であり、実行時の値を自動変換する命令ではない。
\item lambdaは名前を付けずに短い関数オブジェクトを作る構文である。
\item with文はcontext managerに開始・終了処理を任せ、ファイルなどの資源を確実に解放する。
            \end{itemize}
            \paragraph{上から順の処理}
            \begin{enumerate}[leftmargin=1.8em]
            \item maps に \{\} の結果を代入する。
\item sorted((map\_paths or \{\}).items()) を順に走査し、各要素を (key, raw\_path) に入れて処理する。
\item   raw に str(raw\_path or '') の結果を代入する。
\item   path に Path(raw) の結果を代入する。
\item   entry に \{'path': raw, 'exists': bool(path.is\_file())\} の結果を代入する。
\item   条件 path.is\_file() を判定し、真なら内部処理を行う。
\item     digest に hashlib.sha256() の結果を代入する。
\item     with 文で path.open('rb') を管理しながら処理する。
\item       iter(lambda: stream.read(1024 * 1024), b'') を順に走査し、各要素を block に入れて処理する。
\item     entry['sha256'] に digest.hexdigest() の結果を代入する。
\item     entry['size\_bytes'] に int(path.stat().st\_size) の結果を代入する。
\item     上の条件が偽の場合:
\item     entry['sha256'] に '' の結果を代入する。
\item     entry['size\_bytes'] に 0 の結果を代入する。
\item   maps[str(key)] に entry の結果を代入する。
\item canonical に \{'parameters': dict(sorted((parameters or \{\}).items())), 'maps': maps\} の結果を代入する。
\item snapshot\_id に hashlib.sha256(json.dumps(canonical, sort\_keys=True, ensure\_ascii=False, default=str).encode('utf-8')).hexdigest()[:16] の結果を代入する。
\item raw\_base に str(detail\_csv) の結果を代入する。
\item 条件 raw\_base.lower().endswith('.gz') を判定し、真なら内部処理を行う。
\item   raw\_base に raw\_base[:-3] の結果を代入する。
\item 条件 raw\_base.lower().endswith('.csv') を判定し、真なら内部処理を行う。
\item   raw\_base に raw\_base[:-4] の結果を代入する。
\item output に raw\_base + '\_model\_snapshot.yaml' の結果を代入する。
\item payload に \{'schema\_version': 1, 'model\_snapshot\_id': snapshot\_id, 'detail\_csv': str(detail\_csv), 'parameters\_repeated\_in\_detail\_csv': True, 'maps\_referenced\_by\_snapshot\_id': True, **canonical\} の結果を代入する。
\item ensure\_parent\_dir(...) を実行する。
\item Path(output).write\_text(...) を実行する。
\item (snapshot\_id, output) を返す。
            \end{enumerate}
            \paragraph{代表コード断片}
            \begin{lstlisting}[language=Python]
def write_model_snapshot(detail_csv: str, parameters: dict, map_paths: dict) -> tuple[str, str]:
    """Write immutable model provenance once instead of repeating map paths at 1 Hz."""
    maps = {}
    for key, raw_path in sorted((map_paths or {}).items()):
        raw = str(raw_path or "")
        path = Path(raw)
        entry = {"path": raw, "exists": bool(path.is_file())}
        if path.is_file():
            digest = hashlib.sha256()
            with path.open("rb") as stream:
                for block in iter(lambda: stream.read(1024 * 1024), b""):
                    digest.update(block)
            entry["sha256"] = digest.hexdigest()
            entry["size_bytes"] = int(path.stat().st_size)
        else:
            entry["sha256"] = ""
            entry["size_bytes"] = 0
        maps[str(key)] = entry
    canonical = {
        "parameters": dict(sorted((parameters or {}).items())),
        "maps": maps,
    }
    snapshot_id = hashlib.sha256(
        json.dumps(canonical, sort_keys=True, ensure_ascii=False, default=str).encode("utf-8")
    ).hexdigest()[:16]
    raw_base = str(detail_csv)
    if raw_base.lower().endswith(".gz"):
        raw_base = raw_base[:-3]
    if raw_base.lower().endswith(".csv"):
        raw_base = raw_base[:-4]
    output = raw_base + "_model_snapshot.yaml"
    payload = {
        "schema_version": 1,
        "model_snapshot_id": snapshot_id,
        "detail_csv": str(detail_csv),
...
\end{lstlisting}

\subsection{L278 関数 get\_workspace\_revision}
            \begin{itemize}[leftmargin=1.6em]
              \item 行範囲: L278--L328
              \item 所属: module直下
              \item 定義: \path|get_workspace_revision(root_dir: Path) -> dict|
              \item docstring: 明示されていない。
              \item このブロックが直接呼ぶ主な関数・メソッド: bool, lower, run, strip
              \item 戻り値の要点: info / info / info / info
              \item 主なローカル代入名: branch, dirty, head, info, inside
              \item 読み取る主なインスタンス属性: 特になし。
              \item 更新する主なインスタンス属性: 特になし。
              \item 明示的に送出する例外: 明示的raiseはない。
              \item 制御構造の規模: 条件分岐 4、ループ 0、try 2
            \end{itemize}
            \paragraph{この定義を読むためのPython構文}
            \begin{itemize}[leftmargin=1.6em]
            \item `->`以後は戻り値の型注釈であり、実行時の値を自動変換する命令ではない。
\item try文は例外が起き得る処理と、異常時または終了時の経路を分ける。
            \end{itemize}
            \paragraph{上から順の処理}
            \begin{enumerate}[leftmargin=1.8em]
            \item info に \{'git\_available': False, 'git\_head': '', 'git\_branch': '', 'git\_dirty': None\} の結果を代入する。
\item 例外処理を伴う try ブロックを実行する。
\item   inside に subprocess.run(['git', 'rev-parse', '--is-inside-work-tree'], cwd=root\_dir, capture\_output=True, text=True, check=False) の結果を代入する。
\item   Exceptionを捕捉した場合:
\item   info を返す。
\item 条件 inside.returncode != 0 or inside.stdout.strip().lower() != 'true' を判定し、真なら内部処理を行う。
\item   info を返す。
\item info['git\_available'] に True の結果を代入する。
\item 例外処理を伴う try ブロックを実行する。
\item   head に subprocess.run(['git', 'rev-parse', 'HEAD'], cwd=root\_dir, capture\_output=True, text=True, check=False) の結果を代入する。
\item   条件 head.returncode == 0 を判定し、真なら内部処理を行う。
\item     info['git\_head'] に head.stdout.strip() の結果を代入する。
\item   branch に subprocess.run(['git', 'rev-parse', '--abbrev-ref', 'HEAD'], cwd=root\_dir, capture\_output=True, text=True, check=False) の結果を代入する。
\item   条件 branch.returncode == 0 を判定し、真なら内部処理を行う。
\item     info['git\_branch'] に branch.stdout.strip() の結果を代入する。
\item   dirty に subprocess.run(['git', 'status', '--short'], cwd=root\_dir, capture\_output=True, text=True, check=False) の結果を代入する。
\item   条件 dirty.returncode == 0 を判定し、真なら内部処理を行う。
\item     info['git\_dirty'] に bool(dirty.stdout.strip()) の結果を代入する。
\item   Exceptionを捕捉した場合:
\item   info を返す。
\item info を返す。
            \end{enumerate}
            \paragraph{代表コード断片}
            \begin{lstlisting}[language=Python]
def get_workspace_revision(root_dir: Path) -> dict:
    info = {
        'git_available': False,
        'git_head': '',
        'git_branch': '',
        'git_dirty': None,
    }
    try:
        inside = subprocess.run(
            ['git', 'rev-parse', '--is-inside-work-tree'],
            cwd=root_dir,
            capture_output=True,
            text=True,
            check=False,
        )
    except Exception:
        return info
    if inside.returncode != 0 or inside.stdout.strip().lower() != 'true':
        return info
    info['git_available'] = True
    try:
        head = subprocess.run(
            ['git', 'rev-parse', 'HEAD'],
            cwd=root_dir,
            capture_output=True,
            text=True,
            check=False,
        )
        if head.returncode == 0:
            info['git_head'] = head.stdout.strip()
        branch = subprocess.run(
            ['git', 'rev-parse', '--abbrev-ref', 'HEAD'],
            cwd=root_dir,
            capture_output=True,
            text=True,
...
\end{lstlisting}

\subsection{L331 関数 timestamp\_ns}
            \begin{itemize}[leftmargin=1.6em]
              \item 行範囲: L331--L332
              \item 所属: module直下
              \item 定義: \path|timestamp_ns(value) -> int|
              \item docstring: 明示されていない。
              \item このブロックが直接呼ぶ主な関数・メソッド: Timestamp, int
              \item 戻り値の要点: int(pd.Timestamp(value).value)
              \item 主なローカル代入名: 特になし。
              \item 読み取る主なインスタンス属性: 特になし。
              \item 更新する主なインスタンス属性: 特になし。
              \item 明示的に送出する例外: 明示的raiseはない。
              \item 制御構造の規模: 条件分岐 0、ループ 0、try 0
            \end{itemize}
            \paragraph{この定義を読むためのPython構文}
            \begin{itemize}[leftmargin=1.6em]
            \item `->`以後は戻り値の型注釈であり、実行時の値を自動変換する命令ではない。
            \end{itemize}
            \paragraph{上から順の処理}
            \begin{enumerate}[leftmargin=1.8em]
            \item int(pd.Timestamp(value).value) を返す。
            \end{enumerate}
            \paragraph{代表コード断片}
            \begin{lstlisting}[language=Python]
def timestamp_ns(value) -> int:
    return int(pd.Timestamp(value).value)
\end{lstlisting}

\subsection{L335 関数 load\_stops}
            \begin{itemize}[leftmargin=1.6em]
              \item 行範囲: L335--L362
              \item 所属: module直下
              \item 定義: \path|load_stops(path)|
              \item docstring: 明示されていない。
              \item このブロックが直接呼ぶ主な関数・メソッド: append, float, get, isinstance, load\_yaml, max, str
              \item 戻り値の要点: stops
              \item 主なローカル代入名: cfg, dwell\_sec, item, raw\_stops, s\_km, stops
              \item 読み取る主なインスタンス属性: 特になし。
              \item 更新する主なインスタンス属性: 特になし。
              \item 明示的に送出する例外: 明示的raiseはない。
              \item 制御構造の規模: 条件分岐 1、ループ 1、try 3
            \end{itemize}
            \paragraph{この定義を読むためのPython構文}
            \begin{itemize}[leftmargin=1.6em]
            \item try文は例外が起き得る処理と、異常時または終了時の経路を分ける。
            \end{itemize}
            \paragraph{上から順の処理}
            \begin{enumerate}[leftmargin=1.8em]
            \item cfg に load\_yaml(path) の結果を代入する。
\item raw\_stops に cfg.get('stops', []) if isinstance(cfg, dict) else [] の結果を代入する。
\item stops に [] の結果を代入する。
\item raw\_stops or [] を順に走査し、各要素を item に入れて処理する。
\item   条件 isinstance(item, dict) を判定し、真なら内部処理を行う。
\item     例外処理を伴う try ブロックを実行する。
\item       s\_km に float(item.get('s\_km', item.get('dist\_km', 0.0))) の結果を代入する。
\item       Exceptionを捕捉した場合:
\item       Continue 文を実行する。
\item     例外処理を伴う try ブロックを実行する。
\item       dwell\_sec に float(item.get('dwell\_sec', item.get('duration\_sec', 1800.0))) の結果を代入する。
\item       Exceptionを捕捉した場合:
\item       dwell\_sec に 1800.0 の結果を代入する。
\item     stops.append(...) を実行する。
\item     上の条件が偽の場合:
\item     例外処理を伴う try ブロックを実行する。
\item       s\_km に float(item) の結果を代入する。
\item       Exceptionを捕捉した場合:
\item       Continue 文を実行する。
\item     stops.append(...) を実行する。
\item stops を返す。
            \end{enumerate}
            \paragraph{代表コード断片}
            \begin{lstlisting}[language=Python]
def load_stops(path):
    cfg = load_yaml(path)
    raw_stops = cfg.get('stops', []) if isinstance(cfg, dict) else []
    stops = []
    for item in raw_stops or []:
        if isinstance(item, dict):
            try:
                s_km = float(item.get('s_km', item.get('dist_km', 0.0)))
            except Exception:
                continue
            try:
                dwell_sec = float(item.get('dwell_sec', item.get('duration_sec', 1800.0)))
            except Exception:
                dwell_sec = 1800.0
            stops.append({
                's_km': s_km,
                'dwell_sec': max(0.0, dwell_sec),
                'label': str(item.get('label', item.get('name', '')) or ''),
                'window_open_utc': str(item.get('window_open_utc', '') or ''),
                'window_close_utc': str(item.get('window_close_utc', '') or ''),
            })
        else:
            try:
                s_km = float(item)
            except Exception:
                continue
            stops.append({'s_km': s_km, 'dwell_sec': 1800.0, 'label': ''})
    return stops
\end{lstlisting}

\subsection{L365 関数 load\_profile}
            \begin{itemize}[leftmargin=1.6em]
              \item 行範囲: L365--L371
              \item 所属: module直下
              \item 定義: \path|load_profile(path)|
              \item docstring: 明示されていない。
              \item このブロックが直接呼ぶ主な関数・メソッド: read\_csv
              \item 戻り値の要点: None / pd.read\_csv(path) / None
              \item 主なローカル代入名: 特になし。
              \item 読み取る主なインスタンス属性: 特になし。
              \item 更新する主なインスタンス属性: 特になし。
              \item 明示的に送出する例外: 明示的raiseはない。
              \item 制御構造の規模: 条件分岐 1、ループ 0、try 1
            \end{itemize}
            \paragraph{この定義を読むためのPython構文}
            \begin{itemize}[leftmargin=1.6em]
            \item try文は例外が起き得る処理と、異常時または終了時の経路を分ける。
            \end{itemize}
            \paragraph{上から順の処理}
            \begin{enumerate}[leftmargin=1.8em]
            \item 条件 not path を判定し、真なら内部処理を行う。
\item   None を返す。
\item 例外処理を伴う try ブロックを実行する。
\item   pd.read\_csv(path) を返す。
\item   Exceptionを捕捉した場合:
\item   None を返す。
            \end{enumerate}
            \paragraph{代表コード断片}
            \begin{lstlisting}[language=Python]
def load_profile(path):
    if not path:
        return None
    try:
        return pd.read_csv(path)
    except Exception:
        return None
\end{lstlisting}

\subsection{L373 関数 forecast\_distance\_column}
            \begin{itemize}[leftmargin=1.6em]
              \item 行範囲: L373--L380
              \item 所属: module直下
              \item 定義: \path|forecast_distance_column(df: pd.DataFrame) -> str|
              \item docstring: 明示されていない。
              \item このブロックが直接呼ぶ主な関数・メソッド: isinstance
              \item 戻り値の要点: '' / '' / 's\_km' / 'route\_progress\_km'
              \item 主なローカル代入名: 特になし。
              \item 読み取る主なインスタンス属性: 特になし。
              \item 更新する主なインスタンス属性: 特になし。
              \item 明示的に送出する例外: 明示的raiseはない。
              \item 制御構造の規模: 条件分岐 3、ループ 0、try 0
            \end{itemize}
            \paragraph{この定義を読むためのPython構文}
            \begin{itemize}[leftmargin=1.6em]
            \item `->`以後は戻り値の型注釈であり、実行時の値を自動変換する命令ではない。
            \end{itemize}
            \paragraph{上から順の処理}
            \begin{enumerate}[leftmargin=1.8em]
            \item 条件 not isinstance(df, pd.DataFrame) を判定し、真なら内部処理を行う。
\item   '' を返す。
\item 条件 's\_km' in df.columns を判定し、真なら内部処理を行う。
\item   's\_km' を返す。
\item 条件 'route\_progress\_km' in df.columns を判定し、真なら内部処理を行う。
\item   'route\_progress\_km' を返す。
\item '' を返す。
            \end{enumerate}
            \paragraph{代表コード断片}
            \begin{lstlisting}[language=Python]
def forecast_distance_column(df: pd.DataFrame) -> str:
    if not isinstance(df, pd.DataFrame):
        return ''
    if 's_km' in df.columns:
        return 's_km'
    if 'route_progress_km' in df.columns:
        return 'route_progress_km'
    return ''
\end{lstlisting}

\subsection{L383 関数 load\_forecast\_dataframe}
            \begin{itemize}[leftmargin=1.6em]
              \item 行範囲: L383--L407
              \item 所属: module直下
              \item 定義: \path|load_forecast_dataframe(path: str, tzname: str) -> pd.DataFrame|
              \item docstring: 明示されていない。
              \item このブロックが直接呼ぶ主な関数・メソッド: ZoneInfo, copy, drop\_duplicates, forecast\_distance\_column, notna, read\_csv, reset\_index, sort\_values, str, to\_datetime, tz\_convert, tz\_localize
              \item 戻り値の要点: df / df
              \item 主なローカル代入名: dedup\_cols, df, dist\_col, sort\_cols, t, tzname
              \item 読み取る主なインスタンス属性: 特になし。
              \item 更新する主なインスタンス属性: 特になし。
              \item 明示的に送出する例外: 明示的raiseはない。
              \item 制御構造の規模: 条件分岐 3、ループ 0、try 0
            \end{itemize}
            \paragraph{この定義を読むためのPython構文}
            \begin{itemize}[leftmargin=1.6em]
            \item `->`以後は戻り値の型注釈であり、実行時の値を自動変換する命令ではない。
            \end{itemize}
            \paragraph{上から順の処理}
            \begin{enumerate}[leftmargin=1.8em]
            \item df に pd.read\_csv(path) の結果を代入する。
\item 条件 'time' not in df.columns を判定し、真なら内部処理を行う。
\item   df を返す。
\item t に pd.to\_datetime(df['time'], format='mixed', errors='coerce') の結果を代入する。
\item tzname に str(tzname or 'UTC') の結果を代入する。
\item 条件 t.dt.tz is None を判定し、真なら内部処理を行う。
\item   条件 tzname.upper() == 'UTC' を判定し、真なら内部処理を行う。
\item     t に t.dt.tz\_localize('UTC') の結果を代入する。
\item     上の条件が偽の場合:
\item     t に t.dt.tz\_localize(ZoneInfo(tzname), ambiguous='NaT', nonexistent='NaT').dt.tz\_convert('UTC') の結果を代入する。
\item   上の条件が偽の場合:
\item   t に t.dt.tz\_convert('UTC') の結果を代入する。
\item df に df.copy() の結果を代入する。
\item df['time'] に t の結果を代入する。
\item dist\_col に forecast\_distance\_column(df) の結果を代入する。
\item sort\_cols に ['time'] + ([dist\_col] if dist\_col else []) の結果を代入する。
\item dedup\_cols に ['time'] + ([dist\_col] if dist\_col else []) の結果を代入する。
\item df に df.loc[df['time'].notna()].sort\_values(sort\_cols).drop\_duplicates(subset=dedup\_cols, keep='last').reset\_index(drop=True) の結果を代入する。
\item df を返す。
            \end{enumerate}
            \paragraph{代表コード断片}
            \begin{lstlisting}[language=Python]
def load_forecast_dataframe(path: str, tzname: str) -> pd.DataFrame:
    df = pd.read_csv(path)
    if 'time' not in df.columns:
        return df
    t = pd.to_datetime(df['time'], format='mixed', errors='coerce')
    tzname = str(tzname or 'UTC')
    if t.dt.tz is None:
        if tzname.upper() == 'UTC':
            t = t.dt.tz_localize('UTC')
        else:
            t = t.dt.tz_localize(ZoneInfo(tzname), ambiguous='NaT', nonexistent='NaT').dt.tz_convert('UTC')
    else:
        t = t.dt.tz_convert('UTC')
    df = df.copy()
    df['time'] = t
    dist_col = forecast_distance_column(df)
    sort_cols = ['time'] + ([dist_col] if dist_col else [])
    dedup_cols = ['time'] + ([dist_col] if dist_col else [])
    df = (
        df.loc[df['time'].notna()]
        .sort_values(sort_cols)
        .drop_duplicates(subset=dedup_cols, keep='last')
        .reset_index(drop=True)
    )
    return df
\end{lstlisting}

\subsection{L410 関数 merge\_forecast\_dataframes}
            \begin{itemize}[leftmargin=1.6em]
              \item 行範囲: L410--L420
              \item 所属: module直下
              \item 定義: \path|merge_forecast_dataframes(primary: pd.DataFrame, fallback: pd.DataFrame) -> pd.DataFrame|
              \item docstring: 明示されていない。
              \item このブロックが直接呼ぶ主な関数・メソッド: DataFrame, combine\_first, copy, len, reset\_index, set\_index, sort\_index
              \item 戻り値の要点: merged / fallback.copy() if fallback is not None else pd.DataFrame() / primary.copy() / primary.copy()
              \item 主なローカル代入名: fallback\_idx, merged, primary\_idx
              \item 読み取る主なインスタンス属性: 特になし。
              \item 更新する主なインスタンス属性: 特になし。
              \item 明示的に送出する例外: 明示的raiseはない。
              \item 制御構造の規模: 条件分岐 3、ループ 0、try 0
            \end{itemize}
            \paragraph{この定義を読むためのPython構文}
            \begin{itemize}[leftmargin=1.6em]
            \item `->`以後は戻り値の型注釈であり、実行時の値を自動変換する命令ではない。
            \end{itemize}
            \paragraph{上から順の処理}
            \begin{enumerate}[leftmargin=1.8em]
            \item 条件 primary is None or len(primary) == 0 を判定し、真なら内部処理を行う。
\item   fallback.copy() if fallback is not None else pd.DataFrame() を返す。
\item 条件 fallback is None or len(fallback) == 0 を判定し、真なら内部処理を行う。
\item   primary.copy() を返す。
\item 条件 'time' not in primary.columns or 'time' not in fallback.columns を判定し、真なら内部処理を行う。
\item   primary.copy() を返す。
\item primary\_idx に primary.set\_index('time') の結果を代入する。
\item fallback\_idx に fallback.set\_index('time') の結果を代入する。
\item merged に primary\_idx.combine\_first(fallback\_idx).sort\_index().reset\_index() の結果を代入する。
\item merged を返す。
            \end{enumerate}
            \paragraph{代表コード断片}
            \begin{lstlisting}[language=Python]
def merge_forecast_dataframes(primary: pd.DataFrame, fallback: pd.DataFrame) -> pd.DataFrame:
    if primary is None or len(primary) == 0:
        return fallback.copy() if fallback is not None else pd.DataFrame()
    if fallback is None or len(fallback) == 0:
        return primary.copy()
    if 'time' not in primary.columns or 'time' not in fallback.columns:
        return primary.copy()
    primary_idx = primary.set_index('time')
    fallback_idx = fallback.set_index('time')
    merged = primary_idx.combine_first(fallback_idx).sort_index().reset_index()
    return merged
\end{lstlisting}

\subsection{L423 関数 build\_forecast\_grid\_payload}
            \begin{itemize}[leftmargin=1.6em]
              \item 行範囲: L423--L473
              \item 所属: module直下
              \item 定義: \path!build_forecast_grid_payload(df: pd.DataFrame) -> dict | None!
              \item docstring: 明示されていない。
              \item このブロックが直接呼ぶ主な関数・メソッド: Index, apply, array, bfill, copy, drop\_duplicates, dropna, ffill, forecast\_distance\_column, interpolate, len, max
              \item 戻り値の要点: \{'dist\_col': dist\_col, 'time\_ns': np.array([timestamp\_ns(value) for value in time\_index], dtype=np.int64), 's\_grid': s\_grid, 'matrices': matrices\} / None / None / None
              \item 主なローカル代入名: col, dist\_col, matrices, pivot, s\_grid, time\_index, value, work
              \item 読み取る主なインスタンス属性: 特になし。
              \item 更新する主なインスタンス属性: 特になし。
              \item 明示的に送出する例外: 明示的raiseはない。
              \item 制御構造の規模: 条件分岐 6、ループ 1、try 0
            \end{itemize}
            \paragraph{この定義を読むためのPython構文}
            \begin{itemize}[leftmargin=1.6em]
            \item `->`以後は戻り値の型注釈であり、実行時の値を自動変換する命令ではない。
\item 内包表記は、反復・条件・値生成を一つの式で記述してcollectionを作る。
            \end{itemize}
            \paragraph{上から順の処理}
            \begin{enumerate}[leftmargin=1.8em]
            \item dist\_col に forecast\_distance\_column(df) の結果を代入する。
\item 条件 not dist\_col or 'time' not in df.columns を判定し、真なら内部処理を行う。
\item   None を返す。
\item work に df.copy() の結果を代入する。
\item work[dist\_col] に pd.to\_numeric(work[dist\_col], errors='coerce') の結果を代入する。
\item work に work.dropna(subset=['time', dist\_col]).sort\_values(['time', dist\_col]) の結果を代入する。
\item 条件 work.empty を判定し、真なら内部処理を行う。
\item   None を返す。
\item time\_index に pd.Index(work['time'].drop\_duplicates().sort\_values()) の結果を代入する。
\item s\_grid に np.array(sorted(work[dist\_col].dropna().unique()), dtype=float) の結果を代入する。
\item 条件 len(time\_index) < 2 or len(s\_grid) < 2 を判定し、真なら内部処理を行う。
\item   None を返す。
\item 条件 len(work) <= max(len(time\_index), len(s\_grid)) を判定し、真なら内部処理を行う。
\item   None を返す。
\item work に work.drop\_duplicates(subset=['time', dist\_col], keep='last') の結果を代入する。
\item matrices に \{\} の結果を代入する。
\item ('GHI', 'DNI', 'DHI', 'POA\_drive', 'POA\_stop\_ideal', 'Tamb\_C', 'Tcell\_C', 'Tcell\_drive\_C', 'Tcell\_stop\_ideal\_C', 'headwind\_ms') を順に走査し、各要素を col に入れて処理する。
\item   条件 col not in work.columns を判定し、真なら内部処理を行う。
\item     Continue 文を実行する。
\item   pivot に work.pivot(index='time', columns=dist\_col, values=col).reindex(index=time\_index, columns=s\_grid).apply(pd.to\_numeric, errors='coerce').interpolate(axis=0, limit\_direction='both').interpolate(axis=1, limit\_direction='both').ffill().bfill().ffill(axis=1).bfill(axis=1) の結果を代入する。
\item   matrices[col] に pivot.to\_numpy(dtype=float) の結果を代入する。
\item 条件 not matrices を判定し、真なら内部処理を行う。
\item   None を返す。
\item \{'dist\_col': dist\_col, 'time\_ns': np.array([timestamp\_ns(value) for value in time\_index], dtype=np.int64), 's\_grid': s\_grid, 'matrices': matrices\} を返す。
            \end{enumerate}
            \paragraph{代表コード断片}
            \begin{lstlisting}[language=Python]
def build_forecast_grid_payload(df: pd.DataFrame) -> dict | None:
    dist_col = forecast_distance_column(df)
    if not dist_col or 'time' not in df.columns:
        return None
    work = df.copy()
    work[dist_col] = pd.to_numeric(work[dist_col], errors='coerce')
    work = work.dropna(subset=['time', dist_col]).sort_values(['time', dist_col])
    if work.empty:
        return None
    time_index = pd.Index(work['time'].drop_duplicates().sort_values())
    s_grid = np.array(sorted(work[dist_col].dropna().unique()), dtype=float)
    if len(time_index) < 2 or len(s_grid) < 2:
        return None
    if len(work) <= max(len(time_index), len(s_grid)):
        return None
    work = work.drop_duplicates(subset=['time', dist_col], keep='last')
    matrices = {}
    for col in (
        'GHI',
        'DNI',
        'DHI',
        'POA_drive',
        'POA_stop_ideal',
        'Tamb_C',
        'Tcell_C',
        'Tcell_drive_C',
        'Tcell_stop_ideal_C',
        'headwind_ms',
    ):
        if col not in work.columns:
            continue
        pivot = (
            work.pivot(index='time', columns=dist_col, values=col)
            .reindex(index=time_index, columns=s_grid)
            .apply(pd.to_numeric, errors='coerce')
...
\end{lstlisting}

\subsection{L476 関数 interp\_forecast\_grid}
            \begin{itemize}[leftmargin=1.6em]
              \item 行範囲: L476--L520
              \item 所属: module直下
              \item 定義: \path!interp_forecast_grid(payload: dict | None, col: str, t_utc: datetime, s_km: float | None, default: float) -> float!
              \item docstring: 明示されていない。
              \item このブロックが直接呼ぶ主な関数・メソッド: clip, float, get, int, len, max, searchsorted, timestamp\_ns
              \item 戻り値の要点: float((1.0 - wt) * v0 + wt * v1) / float(default) / float(default)
              \item 主なローカル代入名: denom\_s, denom\_t, i0, i1, i\_hi, j0, j1, j\_hi, matrix, s\_val, sg, t\_ns, tg, v0, v00, v01, v1, v10, v11, ws
              \item 読み取る主なインスタンス属性: 特になし。
              \item 更新する主なインスタンス属性: 特になし。
              \item 明示的に送出する例外: 明示的raiseはない。
              \item 制御構造の規模: 条件分岐 6、ループ 0、try 0
            \end{itemize}
            \paragraph{この定義を読むためのPython構文}
            \begin{itemize}[leftmargin=1.6em]
            \item `->`以後は戻り値の型注釈であり、実行時の値を自動変換する命令ではない。
            \end{itemize}
            \paragraph{上から順の処理}
            \begin{enumerate}[leftmargin=1.8em]
            \item 条件 not payload を判定し、真なら内部処理を行う。
\item   float(default) を返す。
\item matrix に payload.get('matrices', \{\}).get(col) の結果を代入する。
\item tg に payload.get('time\_ns') の結果を代入する。
\item sg に payload.get('s\_grid') の結果を代入する。
\item 条件 matrix is None or tg is None or sg is None or (len(tg) == 0) or (len(sg) == 0) を判定し、真なら内部処理を行う。
\item   float(default) を返す。
\item t\_ns に int(np.clip(timestamp\_ns(t\_utc), int(tg[0]), int(tg[-1]))) の結果を代入する。
\item s\_val に float(s\_km if s\_km is not None else sg[0]) の結果を代入する。
\item s\_val に float(np.clip(s\_val, float(sg[0]), float(sg[-1]))) の結果を代入する。
\item i\_hi に int(np.searchsorted(tg, t\_ns, side='left')) の結果を代入する。
\item 条件 i\_hi <= 0 を判定し、真なら内部処理を行う。
\item   i0, i1 に 0 の結果を代入する。
\item   wt に 0.0 の結果を代入する。
\item   上の条件が偽の場合:
\item   条件 i\_hi >= len(tg) を判定し、真なら内部処理を行う。
\item     i0, i1 に len(tg) - 1 の結果を代入する。
\item     wt に 0.0 の結果を代入する。
\item     上の条件が偽の場合:
\item     i0 に i\_hi - 1 の結果を代入する。
\item     i1 に i\_hi の結果を代入する。
\item     denom\_t に max(int(tg[i1]) - int(tg[i0]), 1) の結果を代入する。
\item     wt に float((t\_ns - int(tg[i0])) / denom\_t) の結果を代入する。
\item j\_hi に int(np.searchsorted(sg, s\_val, side='left')) の結果を代入する。
\item 条件 j\_hi <= 0 を判定し、真なら内部処理を行う。
\item   j0, j1 に 0 の結果を代入する。
\item   ws に 0.0 の結果を代入する。
\item   上の条件が偽の場合:
\item   条件 j\_hi >= len(sg) を判定し、真なら内部処理を行う。
\item     j0, j1 に len(sg) - 1 の結果を代入する。
\item     ws に 0.0 の結果を代入する。
\item     上の条件が偽の場合:
\item     j0 に j\_hi - 1 の結果を代入する。
\item     j1 に j\_hi の結果を代入する。
\item     denom\_s に max(float(sg[j1]) - float(sg[j0]), 1e-09) の結果を代入する。
\item     ws に float((s\_val - float(sg[j0])) / denom\_s) の結果を代入する。
\item v00 に float(matrix[i0, j0]) の結果を代入する。
\item v01 に float(matrix[i0, j1]) の結果を代入する。
\item v10 に float(matrix[i1, j0]) の結果を代入する。
\item v11 に float(matrix[i1, j1]) の結果を代入する。
\item v0 に (1.0 - ws) * v00 + ws * v01 の結果を代入する。
\item v1 に (1.0 - ws) * v10 + ws * v11 の結果を代入する。
\item float((1.0 - wt) * v0 + wt * v1) を返す。
            \end{enumerate}
            \paragraph{代表コード断片}
            \begin{lstlisting}[language=Python]
def interp_forecast_grid(payload: dict | None, col: str, t_utc: datetime, s_km: float | None, default: float) -> float:
    if not payload:
        return float(default)
    matrix = payload.get('matrices', {}).get(col)
    tg = payload.get('time_ns')
    sg = payload.get('s_grid')
    if matrix is None or tg is None or sg is None or len(tg) == 0 or len(sg) == 0:
        return float(default)
    t_ns = int(np.clip(timestamp_ns(t_utc), int(tg[0]), int(tg[-1])))
    s_val = float(s_km if s_km is not None else sg[0])
    s_val = float(np.clip(s_val, float(sg[0]), float(sg[-1])))

    i_hi = int(np.searchsorted(tg, t_ns, side='left'))
    if i_hi <= 0:
        i0 = i1 = 0
        wt = 0.0
    elif i_hi >= len(tg):
        i0 = i1 = len(tg) - 1
        wt = 0.0
    else:
        i0 = i_hi - 1
        i1 = i_hi
        denom_t = max(int(tg[i1]) - int(tg[i0]), 1)
        wt = float((t_ns - int(tg[i0])) / denom_t)

    j_hi = int(np.searchsorted(sg, s_val, side='left'))
    if j_hi <= 0:
        j0 = j1 = 0
        ws = 0.0
    elif j_hi >= len(sg):
        j0 = j1 = len(sg) - 1
        ws = 0.0
    else:
        j0 = j_hi - 1
        j1 = j_hi
...
\end{lstlisting}

\subsection{L523 関数 load\_progress\_reference\_dataframe}
            \begin{itemize}[leftmargin=1.6em]
              \item 行範囲: L523--L540
              \item 所属: module直下
              \item 定義: \path!load_progress_reference_dataframe(path: str) -> pd.DataFrame | None!
              \item docstring: 明示されていない。
              \item このブロックが直接呼ぶ主な関数・メソッド: DataFrame, drop\_duplicates, dropna, read\_csv, reset\_index, sort\_values, to\_datetime, to\_numeric
              \item 戻り値の要点: out.reset\_index(drop=True) / None / None / None
              \item 主なローカル代入名: df, out, t, time\_col
              \item 読み取る主なインスタンス属性: 特になし。
              \item 更新する主なインスタンス属性: 特になし。
              \item 明示的に送出する例外: 明示的raiseはない。
              \item 制御構造の規模: 条件分岐 4、ループ 0、try 1
            \end{itemize}
            \paragraph{この定義を読むためのPython構文}
            \begin{itemize}[leftmargin=1.6em]
            \item `->`以後は戻り値の型注釈であり、実行時の値を自動変換する命令ではない。
\item try文は例外が起き得る処理と、異常時または終了時の経路を分ける。
            \end{itemize}
            \paragraph{上から順の処理}
            \begin{enumerate}[leftmargin=1.8em]
            \item 条件 not path を判定し、真なら内部処理を行う。
\item   None を返す。
\item 例外処理を伴う try ブロックを実行する。
\item   df に pd.read\_csv(path, low\_memory=False) の結果を代入する。
\item   Exceptionを捕捉した場合:
\item   None を返す。
\item time\_col に 'time\_utc' if 'time\_utc' in df.columns else 'time' if 'time' in df.columns else '' の結果を代入する。
\item 条件 not time\_col or 's\_km' not in df.columns を判定し、真なら内部処理を行う。
\item   None を返す。
\item t に pd.to\_datetime(df[time\_col], format='mixed', utc=True, errors='coerce') の結果を代入する。
\item out に pd.DataFrame(\{'time\_utc': t, 's\_km': pd.to\_numeric(df['s\_km'], errors='coerce')\}) の結果を代入する。
\item 条件 'speed\_kmh' in df.columns を判定し、真なら内部処理を行う。
\item   out['speed\_kmh'] に pd.to\_numeric(df['speed\_kmh'], errors='coerce') の結果を代入する。
\item out に out.dropna(subset=['time\_utc', 's\_km']).sort\_values('time\_utc').drop\_duplicates(subset=['time\_utc'], keep='last') の結果を代入する。
\item 条件 out.empty を判定し、真なら内部処理を行う。
\item   None を返す。
\item out.reset\_index(drop=True) を返す。
            \end{enumerate}
            \paragraph{代表コード断片}
            \begin{lstlisting}[language=Python]
def load_progress_reference_dataframe(path: str) -> pd.DataFrame | None:
    if not path:
        return None
    try:
        df = pd.read_csv(path, low_memory=False)
    except Exception:
        return None
    time_col = 'time_utc' if 'time_utc' in df.columns else ('time' if 'time' in df.columns else '')
    if not time_col or 's_km' not in df.columns:
        return None
    t = pd.to_datetime(df[time_col], format='mixed', utc=True, errors='coerce')
    out = pd.DataFrame({'time_utc': t, 's_km': pd.to_numeric(df['s_km'], errors='coerce')})
    if 'speed_kmh' in df.columns:
        out['speed_kmh'] = pd.to_numeric(df['speed_kmh'], errors='coerce')
    out = out.dropna(subset=['time_utc', 's_km']).sort_values('time_utc').drop_duplicates(subset=['time_utc'], keep='last')
    if out.empty:
        return None
    return out.reset_index(drop=True)
\end{lstlisting}

\subsection{L543 関数 ensure\_parent\_dir}
            \begin{itemize}[leftmargin=1.6em]
              \item 行範囲: L543--L546
              \item 所属: module直下
              \item 定義: \path|ensure_parent_dir(path)|
              \item docstring: 明示されていない。
              \item このブロックが直接呼ぶ主な関数・メソッド: dirname, makedirs
              \item 戻り値の要点: 戻り値が無いか、return 文が明示されていない。
              \item 主なローカル代入名: parent
              \item 読み取る主なインスタンス属性: 特になし。
              \item 更新する主なインスタンス属性: 特になし。
              \item 明示的に送出する例外: 明示的raiseはない。
              \item 制御構造の規模: 条件分岐 1、ループ 0、try 0
            \end{itemize}
            \paragraph{この定義を読むためのPython構文}
            \begin{itemize}[leftmargin=1.6em]
            \item 特別な構文上の補足は少ない。
            \end{itemize}
            \paragraph{上から順の処理}
            \begin{enumerate}[leftmargin=1.8em]
            \item parent に os.path.dirname(path) の結果を代入する。
\item 条件 parent を判定し、真なら内部処理を行う。
\item   os.makedirs(...) を実行する。
            \end{enumerate}
            \paragraph{代表コード断片}
            \begin{lstlisting}[language=Python]
def ensure_parent_dir(path):
    parent = os.path.dirname(path)
    if parent:
        os.makedirs(parent, exist_ok=True)
\end{lstlisting}

\subsection{L549 クラス DetailCsvStream}
            \begin{itemize}[leftmargin=1.6em]
              \item 行範囲: L549--L580
              \item 所属: module直下
              \item 定義: \path|DetailCsvStream(bases=none)|
              \item docstring: Write the 1 Hz execution trace without retaining the full race in RAM.
              \item このブロックが直接呼ぶ主な関数・メソッド: 特に明示されていない。
              \item 戻り値の要点: 戻り値が無いか、return 文が明示されていない。
              \item 主なローカル代入名: 特になし。
              \item 読み取る主なインスタンス属性: 特になし。
              \item 更新する主なインスタンス属性: 特になし。
              \item 明示的に送出する例外: 明示的raiseはない。
              \item 制御構造の規模: 条件分岐 0、ループ 0、try 0
            \end{itemize}
            \paragraph{この定義を読むためのPython構文}
            \begin{itemize}[leftmargin=1.6em]
            \item class文は新しい型を定義する。丸括弧内は継承する基底クラスである。
            \end{itemize}
            \paragraph{上から順の処理}
            \begin{enumerate}[leftmargin=1.8em]
            \item docstring または説明文字列を置く。
\item 関数 \_\_init\_\_ を定義する。
\item 関数 write を定義する。
\item 関数 close を定義する。
\item 関数 \_\_enter\_\_ を定義する。
\item 関数 \_\_exit\_\_ を定義する。
            \end{enumerate}
            \paragraph{代表コード断片}
            \begin{lstlisting}[language=Python]
class DetailCsvStream:
    """Write the 1 Hz execution trace without retaining the full race in RAM."""

    def __init__(self, path):
        self.path = str(path)
        self._file = None
        self._writer = None
        self.fieldnames = []
        self.row_count = 0

    def write(self, row):
        if self._writer is None:
            ensure_parent_dir(self.path)
            self.fieldnames = list(row.keys())
            opener = gzip.open if self.path.lower().endswith('.gz') else open
            self._file = opener(self.path, 'wt' if opener is gzip.open else 'w', encoding='utf-8', newline='')
            self._writer = csv.DictWriter(self._file, fieldnames=self.fieldnames, extrasaction='ignore')
            self._writer.writeheader()
        self._writer.writerow(row)
        self.row_count += 1

    def close(self):
        if self._file is not None:
            self._file.flush()
            self._file.close()
            self._file = None

    def __enter__(self):
        return self

    def __exit__(self, exc_type, exc, tb):
        self.close()
\end{lstlisting}

\subsection{L552 関数 DetailCsvStream.\_\_init\_\_}
            \begin{itemize}[leftmargin=1.6em]
              \item 行範囲: L552--L557
              \item 所属: \path|DetailCsvStream|
              \item 定義: \path|__init__(self, path)|
              \item docstring: 明示されていない。
              \item このブロックが直接呼ぶ主な関数・メソッド: str
              \item 戻り値の要点: 戻り値が無いか、return 文が明示されていない。
              \item 主なローカル代入名: 特になし。
              \item 読み取る主なインスタンス属性: 特になし。
              \item 更新する主なインスタンス属性: self.\_file, self.\_writer, self.fieldnames, self.path, self.row\_count
              \item 明示的に送出する例外: 明示的raiseはない。
              \item 制御構造の規模: 条件分岐 0、ループ 0、try 0
            \end{itemize}
            \paragraph{この定義を読むためのPython構文}
            \begin{itemize}[leftmargin=1.6em]
            \item 第1引数selfは、このメソッドを呼んだインスタンス自身を受け取る慣習名である。
            \end{itemize}
            \paragraph{上から順の処理}
            \begin{enumerate}[leftmargin=1.8em]
            \item self.path に str(path) の結果を代入する。
\item self.\_file に None の結果を代入する。
\item self.\_writer に None の結果を代入する。
\item self.fieldnames に [] の結果を代入する。
\item self.row\_count に 0 の結果を代入する。
            \end{enumerate}
            \paragraph{代表コード断片}
            \begin{lstlisting}[language=Python]
    def __init__(self, path):
        self.path = str(path)
        self._file = None
        self._writer = None
        self.fieldnames = []
        self.row_count = 0
\end{lstlisting}

\subsection{L559 関数 DetailCsvStream.write}
            \begin{itemize}[leftmargin=1.6em]
              \item 行範囲: L559--L568
              \item 所属: \path|DetailCsvStream|
              \item 定義: \path|write(self, row)|
              \item docstring: 明示されていない。
              \item このブロックが直接呼ぶ主な関数・メソッド: DictWriter, endswith, ensure\_parent\_dir, keys, list, lower, opener, writeheader, writerow
              \item 戻り値の要点: 戻り値が無いか、return 文が明示されていない。
              \item 主なローカル代入名: opener
              \item 読み取る主なインスタンス属性: self.\_file, self.\_writer, self.fieldnames, self.path
              \item 更新する主なインスタンス属性: self.\_file, self.\_writer, self.fieldnames, self.row\_count
              \item 明示的に送出する例外: 明示的raiseはない。
              \item 制御構造の規模: 条件分岐 1、ループ 0、try 0
            \end{itemize}
            \paragraph{この定義を読むためのPython構文}
            \begin{itemize}[leftmargin=1.6em]
            \item 第1引数selfは、このメソッドを呼んだインスタンス自身を受け取る慣習名である。
            \end{itemize}
            \paragraph{上から順の処理}
            \begin{enumerate}[leftmargin=1.8em]
            \item 条件 self.\_writer is None を判定し、真なら内部処理を行う。
\item   ensure\_parent\_dir(...) を実行する。
\item   self.fieldnames に list(row.keys()) の結果を代入する。
\item   opener に gzip.open if self.path.lower().endswith('.gz') else open の結果を代入する。
\item   self.\_file に opener(self.path, 'wt' if opener is gzip.open else 'w', encoding='utf-8', newline='') の結果を代入する。
\item   self.\_writer に csv.DictWriter(self.\_file, fieldnames=self.fieldnames, extrasaction='ignore') の結果を代入する。
\item   self.\_writer.writeheader(...) を実行する。
\item self.\_writer.writerow(...) を実行する。
\item self.row\_count を Add で更新する。
            \end{enumerate}
            \paragraph{代表コード断片}
            \begin{lstlisting}[language=Python]
    def write(self, row):
        if self._writer is None:
            ensure_parent_dir(self.path)
            self.fieldnames = list(row.keys())
            opener = gzip.open if self.path.lower().endswith('.gz') else open
            self._file = opener(self.path, 'wt' if opener is gzip.open else 'w', encoding='utf-8', newline='')
            self._writer = csv.DictWriter(self._file, fieldnames=self.fieldnames, extrasaction='ignore')
            self._writer.writeheader()
        self._writer.writerow(row)
        self.row_count += 1
\end{lstlisting}

\subsection{L570 関数 DetailCsvStream.close}
            \begin{itemize}[leftmargin=1.6em]
              \item 行範囲: L570--L574
              \item 所属: \path|DetailCsvStream|
              \item 定義: \path|close(self)|
              \item docstring: 明示されていない。
              \item このブロックが直接呼ぶ主な関数・メソッド: close, flush
              \item 戻り値の要点: 戻り値が無いか、return 文が明示されていない。
              \item 主なローカル代入名: 特になし。
              \item 読み取る主なインスタンス属性: self.\_file
              \item 更新する主なインスタンス属性: self.\_file
              \item 明示的に送出する例外: 明示的raiseはない。
              \item 制御構造の規模: 条件分岐 1、ループ 0、try 0
            \end{itemize}
            \paragraph{この定義を読むためのPython構文}
            \begin{itemize}[leftmargin=1.6em]
            \item 第1引数selfは、このメソッドを呼んだインスタンス自身を受け取る慣習名である。
            \end{itemize}
            \paragraph{上から順の処理}
            \begin{enumerate}[leftmargin=1.8em]
            \item 条件 self.\_file is not None を判定し、真なら内部処理を行う。
\item   self.\_file.flush(...) を実行する。
\item   self.\_file.close(...) を実行する。
\item   self.\_file に None の結果を代入する。
            \end{enumerate}
            \paragraph{代表コード断片}
            \begin{lstlisting}[language=Python]
    def close(self):
        if self._file is not None:
            self._file.flush()
            self._file.close()
            self._file = None
\end{lstlisting}

\subsection{L576 関数 DetailCsvStream.\_\_enter\_\_}
            \begin{itemize}[leftmargin=1.6em]
              \item 行範囲: L576--L577
              \item 所属: \path|DetailCsvStream|
              \item 定義: \path|__enter__(self)|
              \item docstring: 明示されていない。
              \item このブロックが直接呼ぶ主な関数・メソッド: 特に明示されていない。
              \item 戻り値の要点: self
              \item 主なローカル代入名: 特になし。
              \item 読み取る主なインスタンス属性: 特になし。
              \item 更新する主なインスタンス属性: 特になし。
              \item 明示的に送出する例外: 明示的raiseはない。
              \item 制御構造の規模: 条件分岐 0、ループ 0、try 0
            \end{itemize}
            \paragraph{この定義を読むためのPython構文}
            \begin{itemize}[leftmargin=1.6em]
            \item 第1引数selfは、このメソッドを呼んだインスタンス自身を受け取る慣習名である。
            \end{itemize}
            \paragraph{上から順の処理}
            \begin{enumerate}[leftmargin=1.8em]
            \item self を返す。
            \end{enumerate}
            \paragraph{代表コード断片}
            \begin{lstlisting}[language=Python]
    def __enter__(self):
        return self
\end{lstlisting}

\subsection{L579 関数 DetailCsvStream.\_\_exit\_\_}
            \begin{itemize}[leftmargin=1.6em]
              \item 行範囲: L579--L580
              \item 所属: \path|DetailCsvStream|
              \item 定義: \path|__exit__(self, exc_type, exc, tb)|
              \item docstring: 明示されていない。
              \item このブロックが直接呼ぶ主な関数・メソッド: close
              \item 戻り値の要点: 戻り値が無いか、return 文が明示されていない。
              \item 主なローカル代入名: 特になし。
              \item 読み取る主なインスタンス属性: self.close
              \item 更新する主なインスタンス属性: 特になし。
              \item 明示的に送出する例外: 明示的raiseはない。
              \item 制御構造の規模: 条件分岐 0、ループ 0、try 0
            \end{itemize}
            \paragraph{この定義を読むためのPython構文}
            \begin{itemize}[leftmargin=1.6em]
            \item 第1引数selfは、このメソッドを呼んだインスタンス自身を受け取る慣習名である。
            \end{itemize}
            \paragraph{上から順の処理}
            \begin{enumerate}[leftmargin=1.8em]
            \item self.close(...) を実行する。
            \end{enumerate}
            \paragraph{代表コード断片}
            \begin{lstlisting}[language=Python]
    def __exit__(self, exc_type, exc, tb):
        self.close()
\end{lstlisting}

\subsection{L583 関数 \_deep\_copy\_cfg}
            \begin{itemize}[leftmargin=1.6em]
              \item 行範囲: L583--L584
              \item 所属: module直下
              \item 定義: \path|_deep_copy_cfg(cfg)|
              \item docstring: 明示されていない。
              \item このブロックが直接呼ぶ主な関数・メソッド: deepcopy, isinstance
              \item 戻り値の要点: copy.deepcopy(cfg) if isinstance(cfg, dict) else \{\}
              \item 主なローカル代入名: 特になし。
              \item 読み取る主なインスタンス属性: 特になし。
              \item 更新する主なインスタンス属性: 特になし。
              \item 明示的に送出する例外: 明示的raiseはない。
              \item 制御構造の規模: 条件分岐 0、ループ 0、try 0
            \end{itemize}
            \paragraph{この定義を読むためのPython構文}
            \begin{itemize}[leftmargin=1.6em]
            \item 特別な構文上の補足は少ない。
            \end{itemize}
            \paragraph{上から順の処理}
            \begin{enumerate}[leftmargin=1.8em]
            \item copy.deepcopy(cfg) if isinstance(cfg, dict) else \{\} を返す。
            \end{enumerate}
            \paragraph{代表コード断片}
            \begin{lstlisting}[language=Python]
def _deep_copy_cfg(cfg):
    return copy.deepcopy(cfg) if isinstance(cfg, dict) else {}
\end{lstlisting}

\subsection{L587 関数 \_set\_nested}
            \begin{itemize}[leftmargin=1.6em]
              \item 行範囲: L587--L598
              \item 所属: module直下
              \item 定義: \path|_set_nested(cfg, dotted_key: str, value)|
              \item docstring: 明示されていない。
              \item このブロックが直接呼ぶ主な関数・メソッド: ValueError, get, isinstance, split, str
              \item 戻り値の要点: 戻り値が無いか、return 文が明示されていない。
              \item 主なローカル代入名: cur, next\_val, part, parts
              \item 読み取る主なインスタンス属性: 特になし。
              \item 更新する主なインスタンス属性: 特になし。
              \item 明示的に送出する例外: ValueError('override key is empty')
              \item 制御構造の規模: 条件分岐 2、ループ 1、try 0
            \end{itemize}
            \paragraph{この定義を読むためのPython構文}
            \begin{itemize}[leftmargin=1.6em]
            \item 内包表記は、反復・条件・値生成を一つの式で記述してcollectionを作る。
            \end{itemize}
            \paragraph{上から順の処理}
            \begin{enumerate}[leftmargin=1.8em]
            \item parts に [part for part in str(dotted\_key).split('.') if part] の結果を代入する。
\item 条件 not parts を判定し、真なら内部処理を行う。
\item   ValueError('override key is empty') を送出する。
\item cur に cfg の結果を代入する。
\item parts[:-1] を順に走査し、各要素を part に入れて処理する。
\item   next\_val に cur.get(part) の結果を代入する。
\item   条件 not isinstance(next\_val, dict) を判定し、真なら内部処理を行う。
\item     next\_val に \{\} の結果を代入する。
\item     cur[part] に next\_val の結果を代入する。
\item   cur に next\_val の結果を代入する。
\item cur[parts[-1]] に value の結果を代入する。
            \end{enumerate}
            \paragraph{代表コード断片}
            \begin{lstlisting}[language=Python]
def _set_nested(cfg, dotted_key: str, value):
    parts = [part for part in str(dotted_key).split('.') if part]
    if not parts:
        raise ValueError('override key is empty')
    cur = cfg
    for part in parts[:-1]:
        next_val = cur.get(part)
        if not isinstance(next_val, dict):
            next_val = {}
            cur[part] = next_val
        cur = next_val
    cur[parts[-1]] = value
\end{lstlisting}

\subsection{L601 関数 apply\_overrides}
            \begin{itemize}[leftmargin=1.6em]
              \item 行範囲: L601--L615
              \item 所属: module直下
              \item 定義: \path|apply_overrides(cfg, overrides)|
              \item docstring: 明示されていない。
              \item このブロックが直接呼ぶ主な関数・メソッド: ValueError, \_deep\_copy\_cfg, \_set\_nested, append, safe\_load, split, str, strip
              \item 戻り値の要点: (cfg, applied)
              \item 主なローカル代入名: applied, cfg, item, key, raw, raw\_value, value
              \item 読み取る主なインスタンス属性: 特になし。
              \item 更新する主なインスタンス属性: 特になし。
              \item 明示的に送出する例外: ValueError(f'Invalid override: \{item\}')
              \item 制御構造の規模: 条件分岐 2、ループ 1、try 0
            \end{itemize}
            \paragraph{この定義を読むためのPython構文}
            \begin{itemize}[leftmargin=1.6em]
            \item 特別な構文上の補足は少ない。
            \end{itemize}
            \paragraph{上から順の処理}
            \begin{enumerate}[leftmargin=1.8em]
            \item cfg に \_deep\_copy\_cfg(cfg) の結果を代入する。
\item applied に [] の結果を代入する。
\item overrides or [] を順に走査し、各要素を raw に入れて処理する。
\item   item に str(raw).strip() の結果を代入する。
\item   条件 not item を判定し、真なら内部処理を行う。
\item     Continue 文を実行する。
\item   条件 '=' not in item を判定し、真なら内部処理を行う。
\item     ValueError(f'Invalid override: \{item\}') を送出する。
\item   (key, raw\_value) に item.split('=', 1) の結果を代入する。
\item   key に key.strip() の結果を代入する。
\item   value に yaml.safe\_load(raw\_value) の結果を代入する。
\item   \_set\_nested(...) を実行する。
\item   applied.append(...) を実行する。
\item (cfg, applied) を返す。
            \end{enumerate}
            \paragraph{代表コード断片}
            \begin{lstlisting}[language=Python]
def apply_overrides(cfg, overrides):
    cfg = _deep_copy_cfg(cfg)
    applied = []
    for raw in overrides or []:
        item = str(raw).strip()
        if not item:
            continue
        if '=' not in item:
            raise ValueError(f'Invalid override: {item}')
        key, raw_value = item.split('=', 1)
        key = key.strip()
        value = yaml.safe_load(raw_value)
        _set_nested(cfg, key, value)
        applied.append({'key': key, 'value': value})
    return cfg, applied
\end{lstlisting}

\subsection{L618 関数 build\_config\_tag}
            \begin{itemize}[leftmargin=1.6em]
              \item 行範囲: L618--L630
              \item 所属: module直下
              \item 定義: \path|build_config_tag(profile_path: str, profile_cfg: dict, overrides = None) -> str|
              \item docstring: 明示されていない。
              \item このブロックが直接呼ぶ主な関数・メソッド: encode, fromtimestamp, getmtime, hexdigest, isinstance, now, safe\_dump, sha1, strftime
              \item 戻り値の要点: f'\{stamp\}\_\{digest\}'
              \item 主なローカル代入名: canonical, digest, mtime, payload, stamp
              \item 読み取る主なインスタンス属性: 特になし。
              \item 更新する主なインスタンス属性: 特になし。
              \item 明示的に送出する例外: 明示的raiseはない。
              \item 制御構造の規模: 条件分岐 0、ループ 0、try 1
            \end{itemize}
            \paragraph{この定義を読むためのPython構文}
            \begin{itemize}[leftmargin=1.6em]
            \item `->`以後は戻り値の型注釈であり、実行時の値を自動変換する命令ではない。
\item try文は例外が起き得る処理と、異常時または終了時の経路を分ける。
            \end{itemize}
            \paragraph{上から順の処理}
            \begin{enumerate}[leftmargin=1.8em]
            \item payload に \{'profile\_cfg': profile\_cfg if isinstance(profile\_cfg, dict) else \{\}, 'overrides': overrides or []\} の結果を代入する。
\item canonical に yaml.safe\_dump(payload, sort\_keys=True, allow\_unicode=True) の結果を代入する。
\item digest に hashlib.sha1(canonical.encode('utf-8')).hexdigest()[:8] の結果を代入する。
\item 例外処理を伴う try ブロックを実行する。
\item   mtime に datetime.fromtimestamp(os.path.getmtime(profile\_path), timezone.utc) の結果を代入する。
\item   Exceptionを捕捉した場合:
\item   mtime に datetime.now(timezone.utc) の結果を代入する。
\item stamp に mtime.strftime('\%Y\%m\%d\_\%H\%M\%S') の結果を代入する。
\item f'\{stamp\}\_\{digest\}' を返す。
            \end{enumerate}
            \paragraph{代表コード断片}
            \begin{lstlisting}[language=Python]
def build_config_tag(profile_path: str, profile_cfg: dict, overrides=None) -> str:
    payload = {
        'profile_cfg': profile_cfg if isinstance(profile_cfg, dict) else {},
        'overrides': overrides or [],
    }
    canonical = yaml.safe_dump(payload, sort_keys=True, allow_unicode=True)
    digest = hashlib.sha1(canonical.encode('utf-8')).hexdigest()[:8]
    try:
        mtime = datetime.fromtimestamp(os.path.getmtime(profile_path), timezone.utc)
    except Exception:
        mtime = datetime.now(timezone.utc)
    stamp = mtime.strftime('%Y%m%d_%H%M%S')
    return f'{stamp}_{digest}'
\end{lstlisting}

\subsection{L633 関数 tag\_output\_path}
            \begin{itemize}[leftmargin=1.6em]
              \item 行範囲: L633--L642
              \item 所属: module直下
              \item 定義: \path|tag_output_path(path_value: str, tag: str, default_ext: str = '') -> str|
              \item docstring: 明示されていない。
              \item このブロックが直接呼ぶ主な関数・メソッド: replace, splitext, str, strip
              \item 戻り値の要点: f'\{stem\}\_\{tag\}\{ext\}' / raw / raw.replace('\{tag\}', tag)
              \item 主なローカル代入名: ext, raw, stem
              \item 読み取る主なインスタンス属性: 特になし。
              \item 更新する主なインスタンス属性: 特になし。
              \item 明示的に送出する例外: 明示的raiseはない。
              \item 制御構造の規模: 条件分岐 3、ループ 0、try 0
            \end{itemize}
            \paragraph{この定義を読むためのPython構文}
            \begin{itemize}[leftmargin=1.6em]
            \item `->`以後は戻り値の型注釈であり、実行時の値を自動変換する命令ではない。
            \end{itemize}
            \paragraph{上から順の処理}
            \begin{enumerate}[leftmargin=1.8em]
            \item raw に str(path\_value or '').strip() の結果を代入する。
\item 条件 not raw or not tag を判定し、真なら内部処理を行う。
\item   raw を返す。
\item 条件 '\{tag\}' in raw を判定し、真なら内部処理を行う。
\item   raw.replace('\{tag\}', tag) を返す。
\item (stem, ext) に os.path.splitext(raw) の結果を代入する。
\item 条件 not ext and default\_ext を判定し、真なら内部処理を行う。
\item   ext に default\_ext の結果を代入する。
\item f'\{stem\}\_\{tag\}\{ext\}' を返す。
            \end{enumerate}
            \paragraph{代表コード断片}
            \begin{lstlisting}[language=Python]
def tag_output_path(path_value: str, tag: str, default_ext: str = '') -> str:
    raw = str(path_value or '').strip()
    if not raw or not tag:
        return raw
    if '{tag}' in raw:
        return raw.replace('{tag}', tag)
    stem, ext = os.path.splitext(raw)
    if not ext and default_ext:
        ext = default_ext
    return f'{stem}_{tag}{ext}'
\end{lstlisting}

\subsection{L645 関数 apply\_profile\_cfg\_to\_args}
            \begin{itemize}[leftmargin=1.6em]
              \item 行範囲: L645--L795
              \item 所属: module直下
              \item 定義: \path|apply_profile_cfg_to_args(profile_path, profile_cfg, args, *, force_output_defaults = False)|
              \item docstring: 明示されていない。
              \item このブロックが直接呼ぶ主な関数・メソッド: bool, build\_config\_tag, clip, endswith, float, get, get\_path, get\_section, getattr, int, isinstance, join
              \item 戻り値の要点: 戻り値が無いか、return 文が明示されていない。
              \item 主なローカル代入名: auto\_version\_outputs, command\_bridge\_cfg, detail\_compression, detail\_csv, exec\_cfg, latest\_manifest\_json, live\_cfg, logging\_cfg, mpc\_cfg, output\_dir, output\_prefix, output\_tag, plan\_csv, report\_html, resolved\_yaml, runtime\_cfg, sim\_cfg, summary\_json
              \item 読み取る主なインスタンス属性: 特になし。
              \item 更新する主なインスタンス属性: 特になし。
              \item 明示的に送出する例外: 明示的raiseはない。
              \item 制御構造の規模: 条件分岐 9、ループ 0、try 0
            \end{itemize}
            \paragraph{この定義を読むためのPython構文}
            \begin{itemize}[leftmargin=1.6em]
            \item 特別な構文上の補足は少ない。
            \end{itemize}
            \paragraph{上から順の処理}
            \begin{enumerate}[leftmargin=1.8em]
            \item 条件 not profile\_cfg を判定し、真なら内部処理を行う。
\item    を返す。
\item sim\_cfg に get\_section(profile\_cfg, 'simulation') の結果を代入する。
\item runtime\_cfg に get\_section(profile\_cfg, 'runtime') の結果を代入する。
\item logging\_cfg に get\_section(profile\_cfg, 'logging') の結果を代入する。
\item auto\_version\_outputs に bool(sim\_cfg.get('auto\_version\_outputs', True)) の結果を代入する。
\item output\_tag に build\_config\_tag(profile\_path, profile\_cfg, getattr(args, 'override', [])) if auto\_version\_outputs else '' の結果を代入する。
\item args.params\_yaml に profile\_path の結果を代入する。
\item args.profile\_path\_resolved に profile\_path の結果を代入する。
\item args.auto\_version\_outputs に auto\_version\_outputs の結果を代入する。
\item args.output\_tag に output\_tag の結果を代入する。
\item args.forecast\_csv に get\_path(profile\_cfg, profile\_path, 'forecast\_csv', args.forecast\_csv) の結果を代入する。
\item args.forecast\_fill\_csv に get\_path(profile\_cfg, profile\_path, 'forecast\_fill\_csv', getattr(args, 'forecast\_fill\_csv', '')) の結果を代入する。
\item args.progress\_reference\_csv に get\_path(profile\_cfg, profile\_path, 'progress\_reference\_csv', getattr(args, 'progress\_reference\_csv', '')) の結果を代入する。
\item args.route\_profile\_csv に get\_path(profile\_cfg, profile\_path, 'route\_profile\_csv', args.route\_profile\_csv) の結果を代入する。
\item args.speed\_profile\_csv に get\_path(profile\_cfg, profile\_path, 'speed\_profile\_csv', args.speed\_profile\_csv) の結果を代入する。
\item args.stop\_yaml に get\_path(profile\_cfg, profile\_path, 'stop\_yaml', args.stop\_yaml) の結果を代入する。
\item args.drive\_schedule\_yaml に get\_path(profile\_cfg, profile\_path, 'drive\_schedule\_yaml', args.drive\_schedule\_yaml) の結果を代入する。
\item args.drive\_eff\_map に get\_path(profile\_cfg, profile\_path, 'drive\_eff\_map', args.drive\_eff\_map) の結果を代入する。
\item args.regen\_eff\_map に get\_path(profile\_cfg, profile\_path, 'regen\_eff\_map', args.regen\_eff\_map) の結果を代入する。
\item args.rint\_map に get\_path(profile\_cfg, profile\_path, 'rint\_map', args.rint\_map) の結果を代入する。
\item args.panel\_eff\_map に get\_path(profile\_cfg, profile\_path, 'panel\_eff\_map', args.panel\_eff\_map) の結果を代入する。
\item args.mppt\_eff\_map に get\_path(profile\_cfg, profile\_path, 'mppt\_eff\_map', args.mppt\_eff\_map) の結果を代入する。
\item args.ocv\_soc\_map に get\_path(profile\_cfg, profile\_path, 'ocv\_soc\_map', args.ocv\_soc\_map) の結果を代入する。
\item args.drive\_map\_eco に get\_path(profile\_cfg, profile\_path, 'drive\_map\_eco', args.drive\_map\_eco) の結果を代入する。
\item args.drive\_map\_power に get\_path(profile\_cfg, profile\_path, 'drive\_map\_power', args.drive\_map\_power) の結果を代入する。
\item args.regen\_map\_eco に get\_path(profile\_cfg, profile\_path, 'regen\_map\_eco', args.regen\_map\_eco) の結果を代入する。
\item args.regen\_map\_power に get\_path(profile\_cfg, profile\_path, 'regen\_map\_power', args.regen\_map\_power) の結果を代入する。
\item args.forecast\_time\_mode に str(sim\_cfg.get('forecast\_time\_mode', runtime\_cfg.get('forecast\_time\_mode', args.forecast\_time\_mode))) の結果を代入する。
\item args.forecast\_time\_tz に str(runtime\_cfg.get('forecast\_time\_tz', args.forecast\_time\_tz)) の結果を代入する。
\item args.dt に float(get\_section(profile\_cfg, 'mpc').get('dt', args.dt)) の結果を代入する。
\item args.horizon\_steps に int(get\_section(profile\_cfg, 'mpc').get('horizon\_steps', args.horizon\_steps)) の結果を代入する。
\item args.soc0 に float(sim\_cfg.get('soc0', args.soc0)) の結果を代入する。
\item args.Tb0 に float(sim\_cfg.get('Tb0', args.Tb0)) の結果を代入する。
\item args.v0\_kmh に float(sim\_cfg.get('v0\_kmh', args.v0\_kmh)) の結果を代入する。
\item args.start\_utc に str(sim\_cfg.get('start\_utc', args.start\_utc or '') or '') の結果を代入する。
\item args.forecast\_start\_time\_utc に str(sim\_cfg.get('forecast\_start\_time\_utc', args.forecast\_start\_time\_utc or '') or '') の結果を代入する。
\item args.start\_s\_km に float(sim\_cfg.get('start\_s\_km', args.start\_s\_km)) の結果を代入する。
\item args.energy\_budget に bool(sim\_cfg.get('energy\_budget', args.energy\_budget)) の結果を代入する。
\item args.solar\_gain に float(get\_section(profile\_cfg, 'mpc').get('solar\_gain', args.solar\_gain)) の結果を代入する。
\item args.poa\_gain\_drive に float(get\_section(profile\_cfg, 'mpc').get('poa\_gain\_drive', args.poa\_gain\_drive)) の結果を代入する。
\item args.poa\_gain\_stop に float(get\_section(profile\_cfg, 'mpc').get('poa\_gain\_stop', args.poa\_gain\_stop)) の結果を代入する。
\item args.stop\_tilt\_fraction に float(np.clip(get\_section(profile\_cfg, 'mpc').get('stop\_tilt\_fraction', args.stop\_tilt\_fraction), 0.0, 1.0)) の結果を代入する。
\item args.control\_stop\_tilt\_fraction に float(np.clip(get\_section(profile\_cfg, 'mpc').get('control\_stop\_tilt\_fraction', args.control\_stop\_tilt\_fraction), 0.0, 1.0)) の結果を代入する。
\item mpc\_cfg に get\_section(profile\_cfg, 'mpc') の結果を代入する。
\item live\_cfg に get\_section(profile\_cfg, 'live') の結果を代入する。
\item command\_bridge\_cfg に live\_cfg.get('command\_bridge', \{\}) if isinstance(live\_cfg.get('command\_bridge', \{\}), dict) else \{\} の結果を代入する。
\item exec\_cfg に sim\_cfg.get('execution\_model', \{\}) if isinstance(sim\_cfg.get('execution\_model', \{\}), dict) else \{\} の結果を代入する。
\item args.exec\_model\_enabled に bool(exec\_cfg.get('enabled', args.exec\_model\_enabled)) の結果を代入する。
\item args.exec\_inner\_dt\_sec に float(exec\_cfg.get('inner\_dt\_sec', args.exec\_inner\_dt\_sec)) の結果を代入する。
\item args.detail\_rate\_hz に float(sim\_cfg.get('detail\_rate\_hz', args.detail\_rate\_hz)) の結果を代入する。
\item args.exec\_tau\_sec に float(exec\_cfg.get('tau\_sec', command\_bridge\_cfg.get('filter\_tau\_sec', mpc\_cfg.get('speed\_meas\_filter\_tau\_sec', args.exec\_tau\_sec)))) の結果を代入する。
\item args.exec\_accel\_limit\_kmhps に float(exec\_cfg.get('accel\_limit\_kmhps', command\_bridge\_cfg.get('accel\_limit\_kmhps', mpc\_cfg.get('lower\_ref\_accel\_limit\_kmhps', args.exec\_accel\_limit\_kmhps)))) の結果を代入する。
\item args.exec\_decel\_limit\_kmhps に float(exec\_cfg.get('decel\_limit\_kmhps', command\_bridge\_cfg.get('decel\_limit\_kmhps', mpc\_cfg.get('lower\_ref\_decel\_limit\_kmhps', args.exec\_decel\_limit\_kmhps)))) の結果を代入する。
\item args.exec\_deadband\_kmh に float(exec\_cfg.get('deadband\_kmh', command\_bridge\_cfg.get('speed\_deadband\_kmh', mpc\_cfg.get('lower\_ref\_deadband\_kmh', args.exec\_deadband\_kmh)))) の結果を代入する。
\item args.exec\_quantize\_step\_kmh に float(exec\_cfg.get('quantize\_step\_kmh', command\_bridge\_cfg.get('speed\_quantize\_step\_kmh', args.exec\_quantize\_step\_kmh))) の結果を代入する。
\item args.exec\_reaction\_delay\_sec に float(exec\_cfg.get('reaction\_delay\_sec', args.exec\_reaction\_delay\_sec)) の結果を代入する。
\item output\_dir に sim\_cfg.get('output\_dir', os.path.join('outputs', 'prerace')) の結果を代入する。
\item output\_prefix に str(sim\_cfg.get('output\_prefix', logging\_cfg.get('file\_prefix', 'solar\_prerace'))) の結果を代入する。
\item 条件 force\_output\_defaults or not args.out\_csv or args.out\_csv == 'solar\_sim\_log.csv' を判定し、真なら内部処理を行う。
\item   args.out\_csv に tag\_output\_path(os.path.join(output\_dir, f'\{output\_prefix\}.csv'), output\_tag, '.csv') の結果を代入する。
\item 条件 force\_output\_defaults or not args.out\_detail\_csv を判定し、真なら内部処理を行う。
\item   detail\_csv に str(sim\_cfg.get('out\_detail\_csv', '') or '') の結果を代入する。
\item   args.out\_detail\_csv に tag\_output\_path(detail\_csv, output\_tag, '.csv') if detail\_csv else tag\_output\_path(os.path.join(output\_dir, f'\{output\_prefix\}\_detail.csv'), output\_tag, '.csv') の結果を代入する。
\item detail\_compression に str(sim\_cfg.get('detail\_compression', '') or '').strip().lower() の結果を代入する。
\item 条件 detail\_compression in ('gzip', 'gz') and (not args.out\_detail\_csv.lower().endswith('.gz')) を判定し、真なら内部処理を行う。
\item   args.out\_detail\_csv を Add で更新する。
\item 条件 force\_output\_defaults or not args.out\_plan\_csv を判定し、真なら内部処理を行う。
\item   plan\_csv に str(sim\_cfg.get('out\_plan\_csv', '') or '') の結果を代入する。
\item   args.out\_plan\_csv に tag\_output\_path(plan\_csv, output\_tag, '.csv') if plan\_csv else tag\_output\_path(os.path.join(output\_dir, f'\{output\_prefix\}\_upper\_plan.csv'), output\_tag, '.csv') の結果を代入する。
\item 条件 force\_output\_defaults or not args.report\_html を判定し、真なら内部処理を行う。
\item   report\_html に str(sim\_cfg.get('report\_html', '') or '') の結果を代入する。
\item   args.report\_html に tag\_output\_path(report\_html, output\_tag, '.html') if report\_html else tag\_output\_path(os.path.join(output\_dir, f'\{output\_prefix\}\_report.html'), output\_tag, '.html') の結果を代入する。
\item 条件 force\_output\_defaults or not args.summary\_json を判定し、真なら内部処理を行う。
\item   summary\_json に str(sim\_cfg.get('summary\_json', '') or '') の結果を代入する。
\item   args.summary\_json に tag\_output\_path(summary\_json, output\_tag, '.json') if summary\_json else tag\_output\_path(os.path.join(output\_dir, f'\{output\_prefix\}\_summary.json'), output\_tag, '.json') の結果を代入する。
\item 条件 force\_output\_defaults or not args.resolved\_yaml を判定し、真なら内部処理を行う。
\item   resolved\_yaml に str(sim\_cfg.get('resolved\_yaml', '') or '') の結果を代入する。
\item   args.resolved\_yaml に tag\_output\_path(resolved\_yaml, output\_tag, '.yaml') if resolved\_yaml else tag\_output\_path(os.path.join(output\_dir, f'\{output\_prefix\}\_resolved.yaml'), output\_tag, '.yaml') の結果を代入する。
\item latest\_manifest\_json に str(sim\_cfg.get('latest\_manifest\_json', '') or '') の結果を代入する。
            \end{enumerate}
            \paragraph{代表コード断片}
            \begin{lstlisting}[language=Python]
def apply_profile_cfg_to_args(profile_path, profile_cfg, args, *, force_output_defaults=False):
    if not profile_cfg:
        return

    sim_cfg = get_section(profile_cfg, 'simulation')
    runtime_cfg = get_section(profile_cfg, 'runtime')
    logging_cfg = get_section(profile_cfg, 'logging')
    auto_version_outputs = bool(sim_cfg.get('auto_version_outputs', True))
    output_tag = build_config_tag(profile_path, profile_cfg, getattr(args, 'override', [])) if auto_version_outputs else ''

    args.params_yaml = profile_path
    args.profile_path_resolved = profile_path
    args.auto_version_outputs = auto_version_outputs
    args.output_tag = output_tag
    args.forecast_csv = get_path(profile_cfg, profile_path, 'forecast_csv', args.forecast_csv)
    args.forecast_fill_csv = get_path(profile_cfg, profile_path, 'forecast_fill_csv', getattr(args, 'forecast_fill_csv', ''))
    args.progress_reference_csv = get_path(profile_cfg, profile_path, 'progress_reference_csv', getattr(args, 'progress_reference_csv', ''))
    args.route_profile_csv = get_path(profile_cfg, profile_path, 'route_profile_csv', args.route_profile_csv)
    args.speed_profile_csv = get_path(profile_cfg, profile_path, 'speed_profile_csv', args.speed_profile_csv)
    args.stop_yaml = get_path(profile_cfg, profile_path, 'stop_yaml', args.stop_yaml)
    args.drive_schedule_yaml = get_path(profile_cfg, profile_path, 'drive_schedule_yaml', args.drive_schedule_yaml)
    args.drive_eff_map = get_path(profile_cfg, profile_path, 'drive_eff_map', args.drive_eff_map)
    args.regen_eff_map = get_path(profile_cfg, profile_path, 'regen_eff_map', args.regen_eff_map)
    args.rint_map = get_path(profile_cfg, profile_path, 'rint_map', args.rint_map)
    args.panel_eff_map = get_path(profile_cfg, profile_path, 'panel_eff_map', args.panel_eff_map)
    args.mppt_eff_map = get_path(profile_cfg, profile_path, 'mppt_eff_map', args.mppt_eff_map)
    args.ocv_soc_map = get_path(profile_cfg, profile_path, 'ocv_soc_map', args.ocv_soc_map)
    args.drive_map_eco = get_path(profile_cfg, profile_path, 'drive_map_eco', args.drive_map_eco)
    args.drive_map_power = get_path(profile_cfg, profile_path, 'drive_map_power', args.drive_map_power)
    args.regen_map_eco = get_path(profile_cfg, profile_path, 'regen_map_eco', args.regen_map_eco)
    args.regen_map_power = get_path(profile_cfg, profile_path, 'regen_map_power', args.regen_map_power)

    args.forecast_time_mode = str(sim_cfg.get('forecast_time_mode', runtime_cfg.get('forecast_time_mode', args.forecast_time_mode)))
    args.forecast_time_tz = str(runtime_cfg.get('forecast_time_tz', args.forecast_time_tz))
    args.dt = float(get_section(profile_cfg, 'mpc').get('dt', args.dt))
...
\end{lstlisting}

\subsection{L798 関数 apply\_profile\_defaults}
            \begin{itemize}[leftmargin=1.6em]
              \item 行範囲: L798--L803
              \item 所属: module直下
              \item 定義: \path|apply_profile_defaults(args)|
              \item docstring: 明示されていない。
              \item このブロックが直接呼ぶ主な関数・メソッド: apply\_profile\_cfg\_to\_args, load\_workflow\_profile
              \item 戻り値の要点: profile\_cfg / \{\}
              \item 主なローカル代入名: profile\_cfg, profile\_path
              \item 読み取る主なインスタンス属性: 特になし。
              \item 更新する主なインスタンス属性: 特になし。
              \item 明示的に送出する例外: 明示的raiseはない。
              \item 制御構造の規模: 条件分岐 1、ループ 0、try 0
            \end{itemize}
            \paragraph{この定義を読むためのPython構文}
            \begin{itemize}[leftmargin=1.6em]
            \item 特別な構文上の補足は少ない。
            \end{itemize}
            \paragraph{上から順の処理}
            \begin{enumerate}[leftmargin=1.8em]
            \item 条件 not args.profile\_yaml を判定し、真なら内部処理を行う。
\item   \{\} を返す。
\item (profile\_path, profile\_cfg) に load\_workflow\_profile(args.profile\_yaml) の結果を代入する。
\item apply\_profile\_cfg\_to\_args(...) を実行する。
\item profile\_cfg を返す。
            \end{enumerate}
            \paragraph{代表コード断片}
            \begin{lstlisting}[language=Python]
def apply_profile_defaults(args):
    if not args.profile_yaml:
        return {}
    profile_path, profile_cfg = load_workflow_profile(args.profile_yaml)
    apply_profile_cfg_to_args(profile_path, profile_cfg, args)
    return profile_cfg
\end{lstlisting}

\subsection{L806 関数 resolve\_config\_path}
            \begin{itemize}[leftmargin=1.6em]
              \item 行範囲: L806--L817
              \item 所属: module直下
              \item 定義: \path|resolve_config_path(profile_path, value)|
              \item docstring: 明示されていない。
              \item このブロックが直接呼ぶ主な関数・メソッド: Path, exists, is\_absolute, resolve, str, strip
              \item 戻り値の要点: str((ROOT / path).resolve()) / '' / str(path) / str(candidate)
              \item 主なローカル代入名: candidate, path, profile\_dir, value
              \item 読み取る主なインスタンス属性: 特になし。
              \item 更新する主なインスタンス属性: 特になし。
              \item 明示的に送出する例外: 明示的raiseはない。
              \item 制御構造の規模: 条件分岐 3、ループ 0、try 0
            \end{itemize}
            \paragraph{この定義を読むためのPython構文}
            \begin{itemize}[leftmargin=1.6em]
            \item 特別な構文上の補足は少ない。
            \end{itemize}
            \paragraph{上から順の処理}
            \begin{enumerate}[leftmargin=1.8em]
            \item value に str(value or '').strip() の結果を代入する。
\item 条件 not value を判定し、真なら内部処理を行う。
\item   '' を返す。
\item path に Path(value) の結果を代入する。
\item 条件 path.is\_absolute() を判定し、真なら内部処理を行う。
\item   str(path) を返す。
\item profile\_dir に Path(profile\_path).resolve().parent if profile\_path else ROOT の結果を代入する。
\item candidate に (profile\_dir / path).resolve() の結果を代入する。
\item 条件 candidate.exists() を判定し、真なら内部処理を行う。
\item   str(candidate) を返す。
\item str((ROOT / path).resolve()) を返す。
            \end{enumerate}
            \paragraph{代表コード断片}
            \begin{lstlisting}[language=Python]
def resolve_config_path(profile_path, value):
    value = str(value or '').strip()
    if not value:
        return ''
    path = Path(value)
    if path.is_absolute():
        return str(path)
    profile_dir = Path(profile_path).resolve().parent if profile_path else ROOT
    candidate = (profile_dir / path).resolve()
    if candidate.exists():
        return str(candidate)
    return str((ROOT / path).resolve())
\end{lstlisting}

\subsection{L820 関数 evaluate\_model\_validation\_gate}
            \begin{itemize}[leftmargin=1.6em]
              \item 行範囲: L820--L1029
              \item 所属: module直下
              \item 定義: \path|evaluate_model_validation_gate(cfg, profile_path)|
              \item docstring: 明示されていない。
              \item このブロックが直接呼ぶ主な関数・メソッド: Path, abs, all, bool, dict, exists, float, get, int, isfinite, isinstance, load\_yaml
              \item 戻り値の要点: result / result
              \item 主なローカル代入名: acceleration\_fit, checks, fit\_summary\_path, gate\_cfg, grade\_fit, identification, label, metrics, observed, predicted, prefix, result, spread, summary, terminal, terminal\_anchor, terminal\_anchor\_role, terminal\_path, terminal\_voltage\_observed, terminal\_voltage\_predicted
              \item 読み取る主なインスタンス属性: 特になし。
              \item 更新する主なインスタンス属性: 特になし。
              \item 明示的に送出する例外: 明示的raiseはない。
              \item 制御構造の規模: 条件分岐 5、ループ 1、try 0
            \end{itemize}
            \paragraph{この定義を読むためのPython構文}
            \begin{itemize}[leftmargin=1.6em]
            \item 特別な構文上の補足は少ない。
            \end{itemize}
            \paragraph{上から順の処理}
            \begin{enumerate}[leftmargin=1.8em]
            \item identification に cfg.get('identification', \{\}) if isinstance(cfg, dict) else \{\} の結果を代入する。
\item gate\_cfg に identification.get('validation\_gate', \{\}) if isinstance(identification, dict) else \{\} の結果を代入する。
\item 条件 not isinstance(gate\_cfg, dict) を判定し、真なら内部処理を行う。
\item   gate\_cfg に \{\} の結果を代入する。
\item fit\_summary\_path に resolve\_config\_path(profile\_path, identification.get('fit\_summary\_yaml', '')) の結果を代入する。
\item terminal\_path に resolve\_config\_path(profile\_path, identification.get('terminal\_consistency\_yaml', '')) の結果を代入する。
\item result に \{'available': False, 'fit\_summary\_yaml': fit\_summary\_path, 'terminal\_consistency\_yaml': terminal\_path, 'gate\_pass': False, 'checks': \{\}, 'thresholds': \{\}\} の結果を代入する。
\item 条件 not fit\_summary\_path or not Path(fit\_summary\_path).exists() を判定し、真なら内部処理を行う。
\item   result['reason'] に 'fit summary is not configured or does not exist' の結果を代入する。
\item   result を返す。
\item summary に load\_yaml(fit\_summary\_path) の結果を代入する。
\item metrics に summary.get('validation\_metrics', \{\}) if isinstance(summary, dict) else \{\} の結果を代入する。
\item thresholds に \{'vehicle\_power\_rmse\_max\_w': float(gate\_cfg.get('vehicle\_power\_rmse\_max\_w', 150.0)), 'vehicle\_voltage\_rmse\_max\_v': float(gate\_cfg.get('vehicle\_voltage\_rmse\_max\_v', 1.0)), 'conditional\_power\_rmse\_max\_w': float(gate\_cfg.get('conditional\_power\_rmse\_max\_w', 150.0)), 'conditional\_voltage\_rmse\_max\_v': float(gate\_cfg.get('conditional\_voltage\_rmse\_max\_v', 1.0)), 'end\_to\_end\_power\_rmse\_max\_w': float(gate\_cfg.get('end\_to\_end\_power\_rmse\_max\_w', 200.0)), 'end\_to\_end\_voltage\_rmse\_max\_v': float(gate\_cfg.get('end\_to\_end\_voltage\_rmse\_max\_v', 2.0)), 'moving\_pv\_rmse\_max\_w': float(gate\_cfg.get('moving\_pv\_rmse\_max\_w', 150.0)), 'pv\_lodo\_moving\_rmse\_max\_w': float(gate\_cfg.get('pv\_lodo\_moving\_rmse\_max\_w', 150.0)), 'pv\_lodo\_deployed\_stop\_rmse\_max\_w': float(gate\_cfg.get('pv\_lodo\_deployed\_stop\_rmse\_max\_w', 200.0)), 'power\_residual\_mean\_120s\_rmse\_max\_w': float(gate\_cfg.get('power\_residual\_mean\_120s\_rmse\_max\_w', 150.0)), 'energy\_error\_25km\_rmse\_max\_wh': float(gate\_cfg.get('energy\_error\_25km\_rmse\_max\_wh', 35.0)), 'terminal\_soc\_evidence\_spread\_max': float(gate\_cfg.get('terminal\_soc\_evidence\_spread\_max', 0.05)), 'terminal\_replay\_soc\_error\_max': float(gate\_cfg.get('terminal\_replay\_soc\_error\_max', 0.02)), 'terminal\_replay\_voltage\_error\_max\_v': float(gate\_cfg.get('terminal\_replay\_voltage\_error\_max\_v', 0.5)), 'vehicle\_terminal\_soc\_error\_max': float(gate\_cfg.get('vehicle\_terminal\_soc\_error\_max', 0.02)), 'vehicle\_terminal\_voltage\_error\_max\_v': float(gate\_cfg.get('vehicle\_terminal\_voltage\_error\_max\_v', 0.5)), 'end\_to\_end\_terminal\_soc\_error\_max': float(gate\_cfg.get('end\_to\_end\_terminal\_soc\_error\_max', 0.03)), 'end\_to\_end\_terminal\_voltage\_error\_max\_v': float(gate\_cfg.get('end\_to\_end\_terminal\_voltage\_error\_max\_v', 1.0)), 'acceleration\_validation\_rmse\_ratio\_max': float(gate\_cfg.get('acceleration\_validation\_rmse\_ratio\_max', 1.02)), 'acceleration\_validation\_min\_samples': int(gate\_cfg.get('acceleration\_validation\_min\_samples', 100)), 'grade\_validation\_rmse\_ratio\_max': float(gate\_cfg.get('grade\_validation\_rmse\_ratio\_max', 1.02)), 'grade\_validation\_min\_samples': int(gate\_cfg.get('grade\_validation\_min\_samples', 100))\} の結果を代入する。
\item values に \{'vehicle\_power\_rmse\_w': float(metrics.get('power\_rmse\_clean\_w', math.inf)), 'vehicle\_voltage\_rmse\_v': float(metrics.get('voltage\_rmse\_clean\_v', math.inf)), 'conditional\_power\_rmse\_w': float(metrics.get('battery\_conditional\_power\_rmse\_clean\_w', math.inf)), 'conditional\_voltage\_rmse\_v': float(metrics.get('battery\_conditional\_voltage\_rmse\_clean\_v', math.inf)), 'end\_to\_end\_power\_rmse\_w': float(metrics.get('end\_to\_end\_power\_rmse\_clean\_w', math.inf)), 'end\_to\_end\_voltage\_rmse\_v': float(metrics.get('end\_to\_end\_voltage\_rmse\_clean\_v', math.inf)), 'moving\_pv\_rmse\_w': float(metrics.get('end\_to\_end\_moving\_pv\_rmse\_w', math.inf)), 'pv\_lodo\_moving\_rmse\_w': float(metrics.get('pv\_lodo\_moving\_rmse\_w', math.inf)), 'pv\_lodo\_deployed\_stop\_rmse\_w': float(metrics.get('pv\_lodo\_deployed\_stop\_rmse\_w', math.inf)), 'power\_residual\_mean\_120s\_rmse\_w': float(metrics.get('end\_to\_end\_power\_residual\_mean\_120s\_rmse\_w', math.inf)), 'energy\_error\_25km\_rmse\_wh': float(metrics.get('end\_to\_end\_energy\_error\_25km\_rmse\_wh', math.inf))\} の結果を代入する。
\item terminal に load\_yaml(terminal\_path) if terminal\_path and Path(terminal\_path).exists() else \{\} の結果を代入する。
\item spread に float(terminal.get('evidence\_interval\_max', math.inf)) - float(terminal.get('evidence\_interval\_min', -math.inf)) の結果を代入する。
\item values['terminal\_soc\_evidence\_spread'] に spread の結果を代入する。
\item terminal\_anchor\_role に str((summary.get('fit\_plan', \{\}) or \{\}).get('terminal\_anchor\_role', 'unknown')).strip().lower() の結果を代入する。
\item values['terminal\_anchor\_role'] に terminal\_anchor\_role の結果を代入する。
\item values['terminal\_replay\_soc\_error'] に abs(float(metrics.get('battery\_conditional\_retire\_anchor\_soc\_error', math.inf))) の結果を代入する。
\item terminal\_voltage\_observed に float(metrics.get('battery\_conditional\_retire\_anchor\_voltage\_obs\_v', math.nan)) の結果を代入する。
\item terminal\_voltage\_predicted に float(metrics.get('battery\_conditional\_retire\_anchor\_voltage\_pred\_v', math.nan)) の結果を代入する。
\item values['terminal\_replay\_voltage\_error\_v'] に abs(terminal\_voltage\_predicted - terminal\_voltage\_observed) の結果を代入する。
\item (('vehicle', ''), ('end\_to\_end', 'end\_to\_end\_')) を順に走査し、各要素を (label, prefix) に入れて処理する。
\item   values[f'\{label\}\_terminal\_soc\_error'] に abs(float(metrics.get(f'\{prefix\}retire\_anchor\_soc\_error', math.inf))) の結果を代入する。
\item   observed に float(metrics.get(f'\{prefix\}retire\_anchor\_voltage\_obs\_v', math.nan)) の結果を代入する。
\item   predicted に float(metrics.get(f'\{prefix\}retire\_anchor\_voltage\_pred\_v', math.nan)) の結果を代入する。
\item   values[f'\{label\}\_terminal\_voltage\_error\_v'] に abs(predicted - observed) の結果を代入する。
\item checks に \{'vehicle\_power\_rmse': values['vehicle\_power\_rmse\_w'] <= thresholds['vehicle\_power\_rmse\_max\_w'], 'vehicle\_voltage\_rmse': values['vehicle\_voltage\_rmse\_v'] <= thresholds['vehicle\_voltage\_rmse\_max\_v'], 'conditional\_power\_rmse': values['conditional\_power\_rmse\_w'] <= thresholds['conditional\_power\_rmse\_max\_w'], 'conditional\_voltage\_rmse': values['conditional\_voltage\_rmse\_v'] <= thresholds['conditional\_voltage\_rmse\_max\_v'], 'end\_to\_end\_power\_rmse': values['end\_to\_end\_power\_rmse\_w'] <= thresholds['end\_to\_end\_power\_rmse\_max\_w'], 'end\_to\_end\_voltage\_rmse': values['end\_to\_end\_voltage\_rmse\_v'] <= thresholds['end\_to\_end\_voltage\_rmse\_max\_v'], 'moving\_pv\_rmse': values['moving\_pv\_rmse\_w'] <= thresholds['moving\_pv\_rmse\_max\_w'], 'pv\_lodo\_moving\_rmse': values['pv\_lodo\_moving\_rmse\_w'] <= thresholds['pv\_lodo\_moving\_rmse\_max\_w'], 'pv\_lodo\_deployed\_stop\_rmse': values['pv\_lodo\_deployed\_stop\_rmse\_w'] <= thresholds['pv\_lodo\_deployed\_stop\_rmse\_max\_w'], 'power\_residual\_mean\_120s\_rmse': values['power\_residual\_mean\_120s\_rmse\_w'] <= thresholds['power\_residual\_mean\_120s\_rmse\_max\_w'], 'energy\_error\_25km\_rmse': values['energy\_error\_25km\_rmse\_wh'] <= thresholds['energy\_error\_25km\_rmse\_max\_wh'], 'terminal\_soc\_evidence\_spread': spread <= thresholds['terminal\_soc\_evidence\_spread\_max'], 'terminal\_replay\_soc': math.isfinite(values['terminal\_replay\_soc\_error']) and values['terminal\_replay\_soc\_error'] <= thresholds['terminal\_replay\_soc\_error\_max'], 'terminal\_replay\_voltage': math.isfinite(values['terminal\_replay\_voltage\_error\_v']) and values['terminal\_replay\_voltage\_error\_v'] <= thresholds['terminal\_replay\_voltage\_error\_max\_v'], 'vehicle\_terminal\_soc': math.isfinite(values['vehicle\_terminal\_soc\_error']) and values['vehicle\_terminal\_soc\_error'] <= thresholds['vehicle\_terminal\_soc\_error\_max'], 'vehicle\_terminal\_voltage': math.isfinite(values['vehicle\_terminal\_voltage\_error\_v']) and values['vehicle\_terminal\_voltage\_error\_v'] <= thresholds['vehicle\_terminal\_voltage\_error\_max\_v'], 'end\_to\_end\_terminal\_soc': math.isfinite(values['end\_to\_end\_terminal\_soc\_error']) and values['end\_to\_end\_terminal\_soc\_error'] <= thresholds['end\_to\_end\_terminal\_soc\_error\_max'], 'end\_to\_end\_terminal\_voltage': math.isfinite(values['end\_to\_end\_terminal\_voltage\_error\_v']) and values['end\_to\_end\_terminal\_voltage\_error\_v'] <= thresholds['end\_to\_end\_terminal\_voltage\_error\_max\_v'], 'terminal\_high\_precision\_evidence': bool(terminal.get('high\_precision\_gate\_pass', False)), 'terminal\_anchor\_role\_operational': terminal\_anchor\_role in \{'independent\_consensus', 'independent\_measurement'\}\} の結果を代入する。
\item terminal\_anchor に dict(summary.get('terminal\_anchor', \{\}) or \{\}) の結果を代入する。
\item 条件 terminal\_anchor を判定し、真なら内部処理を行う。
\item   checks['terminal\_local\_anchor\_quality'] に bool(terminal\_anchor.get('quality\_gate\_pass', False)) の結果を代入する。
\item   checks['terminal\_independent\_cross\_channel\_consistency'] に bool(terminal\_anchor.get('weak\_channel\_cross\_consistency\_gate\_pass', False)) の結果を代入する。
\item acceleration\_fit に dict((summary.get('fit\_plan', \{\}) or \{\}).get('acceleration\_observation\_fit', \{\}) or \{\}) の結果を代入する。
\item 条件 bool(acceleration\_fit.get('enabled', False)) を判定し、真なら内部処理を行う。
\item   validation\_count に int(acceleration\_fit.get('validation\_sample\_count', 0) or 0) の結果を代入する。
\item   validation\_ratio に float(acceleration\_fit.get('validation\_rmse\_ratio', math.inf)) の結果を代入する。
\item   values['acceleration\_validation\_sample\_count'] に validation\_count の結果を代入する。
\item   values['acceleration\_validation\_rmse\_ratio'] に validation\_ratio の結果を代入する。
\item   checks['acceleration\_timestamp\_holdout'] に bool(validation\_count >= thresholds['acceleration\_validation\_min\_samples'] and math.isfinite(validation\_ratio) and (validation\_ratio <= thresholds['acceleration\_validation\_rmse\_ratio\_max'])) の結果を代入する。
\item grade\_fit に dict((summary.get('fit\_plan', \{\}) or \{\}).get('grade\_observation\_fit', \{\}) or \{\}) の結果を代入する。
\item 条件 bool(grade\_fit.get('enabled', False)) を判定し、真なら内部処理を行う。
\item   validation\_count に int(grade\_fit.get('validation\_sample\_count', 0) or 0) の結果を代入する。
\item   validation\_ratio に float(grade\_fit.get('validation\_rmse\_ratio', math.inf)) の結果を代入する。
\item   values['grade\_validation\_sample\_count'] に validation\_count の結果を代入する。
\item   values['grade\_validation\_rmse\_ratio'] に validation\_ratio の結果を代入する。
\item   values['grade\_observation\_adopted'] に bool(grade\_fit.get('adopted', False)) の結果を代入する。
\item   checks['grade\_observation\_holdout'] に bool(validation\_count >= thresholds['grade\_validation\_min\_samples'] and math.isfinite(validation\_ratio) and (validation\_ratio <= thresholds['grade\_validation\_rmse\_ratio\_max'])) の結果を代入する。
\item result.update(...) を実行する。
\item result を返す。
            \end{enumerate}
            \paragraph{代表コード断片}
            \begin{lstlisting}[language=Python]
def evaluate_model_validation_gate(cfg, profile_path):
    identification = cfg.get('identification', {}) if isinstance(cfg, dict) else {}
    gate_cfg = identification.get('validation_gate', {}) if isinstance(identification, dict) else {}
    if not isinstance(gate_cfg, dict):
        gate_cfg = {}
    fit_summary_path = resolve_config_path(profile_path, identification.get('fit_summary_yaml', ''))
    terminal_path = resolve_config_path(profile_path, identification.get('terminal_consistency_yaml', ''))
    result = {
        'available': False,
        'fit_summary_yaml': fit_summary_path,
        'terminal_consistency_yaml': terminal_path,
        'gate_pass': False,
        'checks': {},
        'thresholds': {},
    }
    if not fit_summary_path or not Path(fit_summary_path).exists():
        result['reason'] = 'fit summary is not configured or does not exist'
        return result
    summary = load_yaml(fit_summary_path)
    metrics = summary.get('validation_metrics', {}) if isinstance(summary, dict) else {}
    thresholds = {
        'vehicle_power_rmse_max_w': float(gate_cfg.get('vehicle_power_rmse_max_w', 150.0)),
        'vehicle_voltage_rmse_max_v': float(gate_cfg.get('vehicle_voltage_rmse_max_v', 1.0)),
        'conditional_power_rmse_max_w': float(gate_cfg.get('conditional_power_rmse_max_w', 150.0)),
        'conditional_voltage_rmse_max_v': float(gate_cfg.get('conditional_voltage_rmse_max_v', 1.0)),
        'end_to_end_power_rmse_max_w': float(gate_cfg.get('end_to_end_power_rmse_max_w', 200.0)),
        'end_to_end_voltage_rmse_max_v': float(gate_cfg.get('end_to_end_voltage_rmse_max_v', 2.0)),
        'moving_pv_rmse_max_w': float(gate_cfg.get('moving_pv_rmse_max_w', 150.0)),
        'pv_lodo_moving_rmse_max_w': float(gate_cfg.get('pv_lodo_moving_rmse_max_w', 150.0)),
        'pv_lodo_deployed_stop_rmse_max_w': float(
            gate_cfg.get('pv_lodo_deployed_stop_rmse_max_w', 200.0)
        ),
        'power_residual_mean_120s_rmse_max_w': float(
            gate_cfg.get('power_residual_mean_120s_rmse_max_w', 150.0)
        ),
...
\end{lstlisting}

\subsection{L1032 関数 get\_profile\_val}
            \begin{itemize}[leftmargin=1.6em]
              \item 行範囲: L1032--L1038
              \item 所属: module直下
              \item 定義: \path|get_profile_val(df, s_km, field, default = 0.0)|
              \item docstring: 明示されていない。
              \item このブロックが直接呼ぶ主な関数・メソッド: float, interpolate\_profile, isfinite
              \item 戻り値の要点: float(val) / float(default) / float(default)
              \item 主なローカル代入名: val
              \item 読み取る主なインスタンス属性: 特になし。
              \item 更新する主なインスタンス属性: 特になし。
              \item 明示的に送出する例外: 明示的raiseはない。
              \item 制御構造の規模: 条件分岐 2、ループ 0、try 0
            \end{itemize}
            \paragraph{この定義を読むためのPython構文}
            \begin{itemize}[leftmargin=1.6em]
            \item 特別な構文上の補足は少ない。
            \end{itemize}
            \paragraph{上から順の処理}
            \begin{enumerate}[leftmargin=1.8em]
            \item 条件 df is None or field not in df.columns を判定し、真なら内部処理を行う。
\item   float(default) を返す。
\item val に float(interpolate\_profile(df, s\_km, field, default)) の結果を代入する。
\item 条件 not math.isfinite(val) を判定し、真なら内部処理を行う。
\item   float(default) を返す。
\item float(val) を返す。
            \end{enumerate}
            \paragraph{代表コード断片}
            \begin{lstlisting}[language=Python]
def get_profile_val(df, s_km, field, default=0.0):
    if df is None or field not in df.columns:
        return float(default)
    val = float(interpolate_profile(df, s_km, field, default))
    if not math.isfinite(val):
        return float(default)
    return float(val)
\end{lstlisting}

\subsection{L1040 関数 parse\_utc\_arg}
            \begin{itemize}[leftmargin=1.6em]
              \item 行範囲: L1040--L1049
              \item 所属: module直下
              \item 定義: \path|parse_utc_arg(raw_value: str)|
              \item docstring: 明示されていない。
              \item このブロックが直接呼ぶ主な関数・メソッド: astimezone, endswith, fromisoformat, replace, str, strip
              \item 戻り値の要点: dt.astimezone(timezone.utc) / None
              \item 主なローカル代入名: dt, text
              \item 読み取る主なインスタンス属性: 特になし。
              \item 更新する主なインスタンス属性: 特になし。
              \item 明示的に送出する例外: 明示的raiseはない。
              \item 制御構造の規模: 条件分岐 3、ループ 0、try 0
            \end{itemize}
            \paragraph{この定義を読むためのPython構文}
            \begin{itemize}[leftmargin=1.6em]
            \item 特別な構文上の補足は少ない。
            \end{itemize}
            \paragraph{上から順の処理}
            \begin{enumerate}[leftmargin=1.8em]
            \item text に str(raw\_value or '').strip() の結果を代入する。
\item 条件 not text を判定し、真なら内部処理を行う。
\item   None を返す。
\item 条件 text.endswith('Z') を判定し、真なら内部処理を行う。
\item   text に text[:-1] + '+00:00' の結果を代入する。
\item dt に datetime.fromisoformat(text) の結果を代入する。
\item 条件 dt.tzinfo is None を判定し、真なら内部処理を行う。
\item   dt に dt.replace(tzinfo=timezone.utc) の結果を代入する。
\item dt.astimezone(timezone.utc) を返す。
            \end{enumerate}
            \paragraph{代表コード断片}
            \begin{lstlisting}[language=Python]
def parse_utc_arg(raw_value: str):
    text = str(raw_value or '').strip()
    if not text:
        return None
    if text.endswith('Z'):
        text = text[:-1] + '+00:00'
    dt = datetime.fromisoformat(text)
    if dt.tzinfo is None:
        dt = dt.replace(tzinfo=timezone.utc)
    return dt.astimezone(timezone.utc)
\end{lstlisting}

\subsection{L1052 関数 resolve\_forecast\_mode}
            \begin{itemize}[leftmargin=1.6em]
              \item 行範囲: L1052--L1061
              \item 所属: module直下
              \item 定義: \path|resolve_forecast_mode(mode: str, df: pd.DataFrame) -> str|
              \item docstring: 明示されていない。
              \item このブロックが直接呼ぶ主な関数・メソッド: all, isna, lower, str, strip
              \item 戻り値の要点: raw / 'absolute' if has\_time else 'relative' / 'relative' / 'relative'
              \item 主なローカル代入名: has\_time, raw
              \item 読み取る主なインスタンス属性: 特になし。
              \item 更新する主なインスタンス属性: 特になし。
              \item 明示的に送出する例外: 明示的raiseはない。
              \item 制御構造の規模: 条件分岐 3、ループ 0、try 0
            \end{itemize}
            \paragraph{この定義を読むためのPython構文}
            \begin{itemize}[leftmargin=1.6em]
            \item `->`以後は戻り値の型注釈であり、実行時の値を自動変換する命令ではない。
            \end{itemize}
            \paragraph{上から順の処理}
            \begin{enumerate}[leftmargin=1.8em]
            \item raw に str(mode or 'auto').strip().lower() の結果を代入する。
\item has\_time に 'time' in df.columns and (not df['time'].isna().all()) の結果を代入する。
\item 条件 raw == 'auto' を判定し、真なら内部処理を行う。
\item   'absolute' if has\_time else 'relative' を返す。
\item 条件 raw not in ('absolute', 'relative', 'loop') を判定し、真なら内部処理を行う。
\item   'relative' を返す。
\item 条件 raw == 'absolute' and (not has\_time) を判定し、真なら内部処理を行う。
\item   'relative' を返す。
\item raw を返す。
            \end{enumerate}
            \paragraph{代表コード断片}
            \begin{lstlisting}[language=Python]
def resolve_forecast_mode(mode: str, df: pd.DataFrame) -> str:
    raw = str(mode or 'auto').strip().lower()
    has_time = ('time' in df.columns) and (not df['time'].isna().all())
    if raw == 'auto':
        return 'absolute' if has_time else 'relative'
    if raw not in ('absolute', 'relative', 'loop'):
        return 'relative'
    if raw == 'absolute' and not has_time:
        return 'relative'
    return raw
\end{lstlisting}

\subsection{L1064 関数 forecast\_row\_index}
            \begin{itemize}[leftmargin=1.6em]
              \item 行範囲: L1064--L1079
              \item 所属: module直下
              \item 定義: \path|forecast_row_index(df: pd.DataFrame, sim_t: datetime, *, dt_sec: float, mode: str, forecast_start_time: datetime) -> int|
              \item docstring: 明示されていない。
              \item このブロックが直接呼ぶ主な関数・メソッド: Timestamp, all, clip, int, isna, len, max, searchsorted, total\_seconds
              \item 戻り値の要点: int(np.clip(idx, 0, len(df) - 1)) / 0 / int(idx \% len(df)) / int(np.clip(idx, 0, len(df) - 1))
              \item 主なローカル代入名: elapsed, idx, mode, t\_max, t\_min
              \item 読み取る主なインスタンス属性: 特になし。
              \item 更新する主なインスタンス属性: 特になし。
              \item 明示的に送出する例外: 明示的raiseはない。
              \item 制御構造の規模: 条件分岐 4、ループ 0、try 0
            \end{itemize}
            \paragraph{この定義を読むためのPython構文}
            \begin{itemize}[leftmargin=1.6em]
            \item `->`以後は戻り値の型注釈であり、実行時の値を自動変換する命令ではない。
            \end{itemize}
            \paragraph{上から順の処理}
            \begin{enumerate}[leftmargin=1.8em]
            \item 条件 len(df) == 0 を判定し、真なら内部処理を行う。
\item   0 を返す。
\item 条件 mode == 'absolute' and 'time' in df.columns and (not df['time'].isna().all()) を判定し、真なら内部処理を行う。
\item   t\_min に df['time'].iloc[0] の結果を代入する。
\item   t\_max に df['time'].iloc[-1] の結果を代入する。
\item   条件 t\_min <= sim\_t <= t\_max を判定し、真なら内部処理を行う。
\item     idx に int(df['time'].searchsorted(pd.Timestamp(sim\_t), side='right') - 1) の結果を代入する。
\item     int(np.clip(idx, 0, len(df) - 1)) を返す。
\item   mode に 'relative' の結果を代入する。
\item elapsed に max(0.0, (sim\_t - forecast\_start\_time).total\_seconds()) の結果を代入する。
\item idx に int(elapsed / max(dt\_sec, 0.001)) の結果を代入する。
\item 条件 mode == 'loop' を判定し、真なら内部処理を行う。
\item   int(idx \% len(df)) を返す。
\item int(np.clip(idx, 0, len(df) - 1)) を返す。
            \end{enumerate}
            \paragraph{代表コード断片}
            \begin{lstlisting}[language=Python]
def forecast_row_index(df: pd.DataFrame, sim_t: datetime, *, dt_sec: float, mode: str, forecast_start_time: datetime) -> int:
    if len(df) == 0:
        return 0
    if mode == 'absolute' and 'time' in df.columns and not df['time'].isna().all():
        t_min = df['time'].iloc[0]
        t_max = df['time'].iloc[-1]
        if t_min <= sim_t <= t_max:
            idx = int(df['time'].searchsorted(pd.Timestamp(sim_t), side='right') - 1)
            return int(np.clip(idx, 0, len(df) - 1))
        mode = 'relative'

    elapsed = max(0.0, (sim_t - forecast_start_time).total_seconds())
    idx = int(elapsed / max(dt_sec, 1.0e-3))
    if mode == 'loop':
        return int(idx % len(df))
    return int(np.clip(idx, 0, len(df) - 1))
\end{lstlisting}

\subsection{L1082 関数 format\_metric}
            \begin{itemize}[leftmargin=1.6em]
              \item 行範囲: L1082--L1091
              \item 所属: module直下
              \item 定義: \path|format_metric(value, digits = 3, default = '--')|
              \item docstring: 明示されていない。
              \item このブロックが直接呼ぶ主な関数・メソッド: float, isfinite
              \item 戻り値の要点: default / default / f'\{fval:.\{digits\}f\}'
              \item 主なローカル代入名: fval
              \item 読み取る主なインスタンス属性: 特になし。
              \item 更新する主なインスタンス属性: 特になし。
              \item 明示的に送出する例外: 明示的raiseはない。
              \item 制御構造の規模: 条件分岐 2、ループ 0、try 1
            \end{itemize}
            \paragraph{この定義を読むためのPython構文}
            \begin{itemize}[leftmargin=1.6em]
            \item try文は例外が起き得る処理と、異常時または終了時の経路を分ける。
            \end{itemize}
            \paragraph{上から順の処理}
            \begin{enumerate}[leftmargin=1.8em]
            \item 条件 value is None を判定し、真なら内部処理を行う。
\item   default を返す。
\item 例外処理を伴う try ブロックを実行する。
\item   fval に float(value) の結果を代入する。
\item   条件 math.isfinite(fval) を判定し、真なら内部処理を行う。
\item     f'\{fval:.\{digits\}f\}' を返す。
\item   Exceptionを捕捉した場合:
\item   Pass 文を実行する。
\item default を返す。
            \end{enumerate}
            \paragraph{代表コード断片}
            \begin{lstlisting}[language=Python]
def format_metric(value, digits=3, default='--'):
    if value is None:
        return default
    try:
        fval = float(value)
        if math.isfinite(fval):
            return f'{fval:.{digits}f}'
    except Exception:
        pass
    return default
\end{lstlisting}

\subsection{L1094 関数 decimate\_xy}
            \begin{itemize}[leftmargin=1.6em]
              \item 行範囲: L1094--L1102
              \item 所属: module直下
              \item 定義: \path|decimate_xy(xs, ys, max_points = 600)|
              \item docstring: 明示されていない。
              \item このブロックが直接呼ぶ主な関数・メソッド: append, ceil, float, int, isfinite, len, max, zip
              \item 戻り値の要点: reduced / pts
              \item 主なローカル代入名: pts, reduced, step, x, y
              \item 読み取る主なインスタンス属性: 特になし。
              \item 更新する主なインスタンス属性: 特になし。
              \item 明示的に送出する例外: 明示的raiseはない。
              \item 制御構造の規模: 条件分岐 2、ループ 0、try 0
            \end{itemize}
            \paragraph{この定義を読むためのPython構文}
            \begin{itemize}[leftmargin=1.6em]
            \item 内包表記は、反復・条件・値生成を一つの式で記述してcollectionを作る。
            \end{itemize}
            \paragraph{上から順の処理}
            \begin{enumerate}[leftmargin=1.8em]
            \item pts に [(float(x), float(y)) for x, y in zip(xs, ys) if math.isfinite(float(y))] の結果を代入する。
\item 条件 len(pts) <= max\_points を判定し、真なら内部処理を行う。
\item   pts を返す。
\item step に max(1, int(math.ceil(len(pts) / max\_points))) の結果を代入する。
\item reduced に pts[::step] の結果を代入する。
\item 条件 reduced[-1] != pts[-1] を判定し、真なら内部処理を行う。
\item   reduced.append(...) を実行する。
\item reduced を返す。
            \end{enumerate}
            \paragraph{代表コード断片}
            \begin{lstlisting}[language=Python]
def decimate_xy(xs, ys, max_points=600):
    pts = [(float(x), float(y)) for x, y in zip(xs, ys) if math.isfinite(float(y))]
    if len(pts) <= max_points:
        return pts
    step = max(1, int(math.ceil(len(pts) / max_points)))
    reduced = pts[::step]
    if reduced[-1] != pts[-1]:
        reduced.append(pts[-1])
    return reduced
\end{lstlisting}

\subsection{L1105 関数 build\_svg\_chart}
            \begin{itemize}[leftmargin=1.6em]
              \item 行範囲: L1105--L1143
              \item 所属: module直下
              \item 定義: \path|build_svg_chart(xs, ys, *, color = '#135d66', width = 920, height = 220, pad = 26, label = '')|
              \item docstring: 明示されていない。
              \item このブロックが直接呼ぶ主な関数・メソッド: abs, append, decimate\_xy, escape, format\_metric, join, max, min
              \item 戻り値の要点: f'<svg viewBox="0 0 \{width\} \{height\}" class="chart-svg" role="img" aria-label="\{html.escape(label)\}"><rect x="0" y="0" width="\{width\}" height="\{height\}" rx="18" ry="18" fill="\#fffdf8" stroke="\#d8d1c4" /><line x1="\{pad\}" y1="\{pad\}" x2="\{pad\}" y2="\{height - pad\}" stroke="\#a9a293" stroke-width="1" /><line x1="\{pad\}" y1="\{height - pad\}" x2="\{width - pad\}" y2="\{height - pad\}" stroke="\#a9a293" stroke-width="1" /><polyline fill="none" stroke="\{color\}" stroke-width="2.6" points="\{polyline\}" /><text x="\{pad\}" y="18" font-size="12" fill="\#50483f">\{html.escape(label)\}</text><text x="\{width - pad\}" y="18" text-anchor="end" font-size="11" fill="\#6f665b">\{format\_metric(y\_max, 2)\}</text><text x="\{width - pad\}" y="\{height - 8\}" text-anchor="end" font-size="11" fill="\#6f665b">\{format\_metric(x\_max, 1)\}</text><text x="\{pad\}" y="\{height - 8\}" font-size="11" fill="\#6f665b">\{format\_metric(x\_min, 1)\}</text><text x="\{width - pad\}" y="\{height / 2:.1f\}" text-anchor="end" font-size="11" fill="\#6f665b">\{format\_metric(y\_mid, 2)\}</text><text x="\{width - pad\}" y="\{height - pad + 16\}" text-anchor="end" font-size="11" fill="\#6f665b">x</text></svg>' / '<div class="chart-empty">no data</div>'
              \item 主なローカル代入名: coords, inner\_h, inner\_w, p, polyline, pts, span, sx, sy, x\_max, x\_min, x\_val, x\_vals, y\_max, y\_mid, y\_min, y\_val, y\_vals
              \item 読み取る主なインスタンス属性: 特になし。
              \item 更新する主なインスタンス属性: 特になし。
              \item 明示的に送出する例外: 明示的raiseはない。
              \item 制御構造の規模: 条件分岐 3、ループ 1、try 0
            \end{itemize}
            \paragraph{この定義を読むためのPython構文}
            \begin{itemize}[leftmargin=1.6em]
            \item 内包表記は、反復・条件・値生成を一つの式で記述してcollectionを作る。
            \end{itemize}
            \paragraph{上から順の処理}
            \begin{enumerate}[leftmargin=1.8em]
            \item pts に decimate\_xy(xs, ys, max\_points=700) の結果を代入する。
\item 条件 not pts を判定し、真なら内部処理を行う。
\item   '<div class="chart-empty">no data</div>' を返す。
\item x\_vals に [p[0] for p in pts] の結果を代入する。
\item y\_vals に [p[1] for p in pts] の結果を代入する。
\item x\_min に min(x\_vals) の結果を代入する。
\item x\_max に max(x\_vals) の結果を代入する。
\item y\_min に min(y\_vals) の結果を代入する。
\item y\_max に max(y\_vals) の結果を代入する。
\item 条件 x\_max <= x\_min を判定し、真なら内部処理を行う。
\item   x\_max に x\_min + 1.0 の結果を代入する。
\item 条件 y\_max <= y\_min を判定し、真なら内部処理を行う。
\item   span に 1.0 if abs(y\_max) < 1.0 else abs(y\_max) * 0.1 の結果を代入する。
\item   y\_min を Sub で更新する。
\item   y\_max を Add で更新する。
\item inner\_w に max(10.0, width - 2 * pad) の結果を代入する。
\item inner\_h に max(10.0, height - 2 * pad) の結果を代入する。
\item coords に [] の結果を代入する。
\item pts を順に走査し、各要素を (x\_val, y\_val) に入れて処理する。
\item   sx に pad + (x\_val - x\_min) / (x\_max - x\_min) * inner\_w の結果を代入する。
\item   sy に height - pad - (y\_val - y\_min) / (y\_max - y\_min) * inner\_h の結果を代入する。
\item   coords.append(...) を実行する。
\item polyline に ' '.join(coords) の結果を代入する。
\item y\_mid に 0.5 * (y\_min + y\_max) の結果を代入する。
\item f'<svg viewBox="0 0 \{width\} \{height\}" class="chart-svg" role="img" aria-label="\{html.escape(label)\}"><rect x="0" y="0" width="\{width\}" height="\{height\}" rx="18" ry="18" fill="\#fffdf8" stroke="\#d8d1c4" /><line x1="\{pad\}" y1="\{pad\}" x2="\{pad\}" y2="\{height - pad\}" stroke="\#a9a293" stroke-width="1" /><line x1="\{pad\}" y1="\{height - pad\}" x2="\{width - pad\}" y2="\{height - pad\}" stroke="\#a9a293" stroke-width="1" /><polyline fill="none" stroke="\{color\}" stroke-width="2.6" points="\{polyline\}" /><text x="\{pad\}" y="18" font-size="12" fill="\#50483f">\{html.escape(label)\}</text><text x="\{width - pad\}" y="18" text-anchor="end" font-size="11" fill="\#6f665b">\{format\_metric(y\_max, 2)\}</text><text x="\{width - pad\}" y="\{height - 8\}" text-anchor="end" font-size="11" fill="\#6f665b">\{format\_metric(x\_max, 1)\}</text><text x="\{pad\}" y="\{height - 8\}" font-size="11" fill="\#6f665b">\{format\_metric(x\_min, 1)\}</text><text x="\{width - pad\}" y="\{height / 2:.1f\}" text-anchor="end" font-size="11" fill="\#6f665b">\{format\_metric(y\_mid, 2)\}</text><text x="\{width - pad\}" y="\{height - pad + 16\}" text-anchor="end" font-size="11" fill="\#6f665b">x</text></svg>' を返す。
            \end{enumerate}
            \paragraph{代表コード断片}
            \begin{lstlisting}[language=Python]
def build_svg_chart(xs, ys, *, color='#135d66', width=920, height=220, pad=26, label=''):
    pts = decimate_xy(xs, ys, max_points=700)
    if not pts:
        return '<div class="chart-empty">no data</div>'
    x_vals = [p[0] for p in pts]
    y_vals = [p[1] for p in pts]
    x_min = min(x_vals)
    x_max = max(x_vals)
    y_min = min(y_vals)
    y_max = max(y_vals)
    if x_max <= x_min:
        x_max = x_min + 1.0
    if y_max <= y_min:
        span = 1.0 if abs(y_max) < 1.0 else abs(y_max) * 0.1
        y_min -= span
        y_max += span
    inner_w = max(10.0, width - 2 * pad)
    inner_h = max(10.0, height - 2 * pad)
    coords = []
    for x_val, y_val in pts:
        sx = pad + (x_val - x_min) / (x_max - x_min) * inner_w
        sy = height - pad - (y_val - y_min) / (y_max - y_min) * inner_h
        coords.append(f'{sx:.2f},{sy:.2f}')
    polyline = ' '.join(coords)
    y_mid = 0.5 * (y_min + y_max)
    return (
        f'<svg viewBox="0 0 {width} {height}" class="chart-svg" role="img" aria-label="{html.escape(label)}">'
        f'<rect x="0" y="0" width="{width}" height="{height}" rx="18" ry="18" fill="#fffdf8" stroke="#d8d1c4" />'
        f'<line x1="{pad}" y1="{pad}" x2="{pad}" y2="{height - pad}" stroke="#a9a293" stroke-width="1" />'
        f'<line x1="{pad}" y1="{height - pad}" x2="{width - pad}" y2="{height - pad}" stroke="#a9a293" stroke-width="1" />'
        f'<polyline fill="none" stroke="{color}" stroke-width="2.6" points="{polyline}" />'
        f'<text x="{pad}" y="18" font-size="12" fill="#50483f">{html.escape(label)}</text>'
        f'<text x="{width - pad}" y="18" text-anchor="end" font-size="11" fill="#6f665b">{format_metric(y_max, 2)}</text>'
        f'<text x="{width - pad}" y="{height - 8}" text-anchor="end" font-size="11" fill="#6f665b">{format_metric(x_max, 1)}</text>'
        f'<text x="{pad}" y="{height - 8}" font-size="11" fill="#6f665b">{format_metric(x_min, 1)}</text>'
...
\end{lstlisting}

\subsection{L1146 関数 flatten\_params\_for\_report}
            \begin{itemize}[leftmargin=1.6em]
              \item 行範囲: L1146--L1163
              \item 所属: module直下
              \item 定義: \path|flatten_params_for_report(cfg)|
              \item docstring: 明示されていない。
              \item このブロックが直接呼ぶ主な関数・メソッド: append, get, isinstance, items, str, visit
              \item 戻り値の要点: rows / rows
              \item 主なローカル代入名: child, child\_prefix, key, rows, section\_name
              \item 読み取る主なインスタンス属性: 特になし。
              \item 更新する主なインスタンス属性: 特になし。
              \item 明示的に送出する例外: 明示的raiseはない。
              \item 制御構造の規模: 条件分岐 3、ループ 2、try 0
            \end{itemize}
            \paragraph{この定義を読むためのPython構文}
            \begin{itemize}[leftmargin=1.6em]
            \item 特別な構文上の補足は少ない。
            \end{itemize}
            \paragraph{上から順の処理}
            \begin{enumerate}[leftmargin=1.8em]
            \item rows に [] の結果を代入する。
\item 関数 visit を定義する。
\item 条件 not isinstance(cfg, dict) を判定し、真なら内部処理を行う。
\item   rows を返す。
\item ('simulation', 'model', 'mpc', 'runtime') を順に走査し、各要素を section\_name に入れて処理する。
\item   visit(...) を実行する。
\item rows を返す。
            \end{enumerate}
            \paragraph{代表コード断片}
            \begin{lstlisting}[language=Python]
def flatten_params_for_report(cfg):
    rows = []

    def visit(prefix, value):
        if isinstance(value, dict):
            for key, child in value.items():
                child_prefix = f'{prefix}.{key}' if prefix else str(key)
                visit(child_prefix, child)
            return
        if isinstance(value, list):
            return
        rows.append((prefix, value))

    if not isinstance(cfg, dict):
        return rows
    for section_name in ('simulation', 'model', 'mpc', 'runtime'):
        visit(section_name, cfg.get(section_name, {}))
    return rows
\end{lstlisting}

\subsection{L1149 関数 flatten\_params\_for\_report.visit}
            \begin{itemize}[leftmargin=1.6em]
              \item 行範囲: L1149--L1157
              \item 所属: \path|flatten_params_for_report|
              \item 定義: \path|visit(prefix, value)|
              \item docstring: 明示されていない。
              \item このブロックが直接呼ぶ主な関数・メソッド: append, isinstance, items, str, visit
              \item 戻り値の要点: 戻り値が無いか、return 文が明示されていない。
              \item 主なローカル代入名: child, child\_prefix, key
              \item 読み取る主なインスタンス属性: 特になし。
              \item 更新する主なインスタンス属性: 特になし。
              \item 明示的に送出する例外: 明示的raiseはない。
              \item 制御構造の規模: 条件分岐 2、ループ 1、try 0
            \end{itemize}
            \paragraph{この定義を読むためのPython構文}
            \begin{itemize}[leftmargin=1.6em]
            \item 特別な構文上の補足は少ない。
            \end{itemize}
            \paragraph{上から順の処理}
            \begin{enumerate}[leftmargin=1.8em]
            \item 条件 isinstance(value, dict) を判定し、真なら内部処理を行う。
\item   value.items() を順に走査し、各要素を (key, child) に入れて処理する。
\item     child\_prefix に f'\{prefix\}.\{key\}' if prefix else str(key) の結果を代入する。
\item     visit(...) を実行する。
\item    を返す。
\item 条件 isinstance(value, list) を判定し、真なら内部処理を行う。
\item    を返す。
\item rows.append(...) を実行する。
            \end{enumerate}
            \paragraph{代表コード断片}
            \begin{lstlisting}[language=Python]
    def visit(prefix, value):
        if isinstance(value, dict):
            for key, child in value.items():
                child_prefix = f'{prefix}.{key}' if prefix else str(key)
                visit(child_prefix, child)
            return
        if isinstance(value, list):
            return
        rows.append((prefix, value))
\end{lstlisting}

\subsection{L1166 関数 write\_simulation\_report}
            \begin{itemize}[leftmargin=1.6em]
              \item 行範囲: L1166--L1286
              \item 所属: module直下
              \item 定義: \path|write_simulation_report(path, summary, detail_df, params_rows)|
              \item docstring: 明示されていない。
              \item このブロックが直接呼ぶ主な関数・メソッド: Series, append, build\_svg\_chart, dumps, ensure\_parent\_dir, escape, extend, format\_metric, get, join, len, list
              \item 戻り値の要点: 戻り値が無いか、return 文が明示されていない。
              \item 主なローカル代入名: cards\_html, charts, charts\_html, f, html\_text, item, key, label, overrides\_html, pack\_svg, params\_html, slope\_svg, soc\_svg, solar\_svg, speed\_cmd\_svg, speed\_exec\_present, speed\_series, speed\_svg, summary\_cards, value
              \item 読み取る主なインスタンス属性: 特になし。
              \item 更新する主なインスタンス属性: 特になし。
              \item 明示的に送出する例外: 明示的raiseはない。
              \item 制御構造の規模: 条件分岐 2、ループ 0、try 0
            \end{itemize}
            \paragraph{この定義を読むためのPython構文}
            \begin{itemize}[leftmargin=1.6em]
            \item 内包表記は、反復・条件・値生成を一つの式で記述してcollectionを作る。
\item with文はcontext managerに開始・終了処理を任せ、ファイルなどの資源を確実に解放する。
            \end{itemize}
            \paragraph{上から順の処理}
            \begin{enumerate}[leftmargin=1.8em]
            \item ensure\_parent\_dir(...) を実行する。
\item x\_index に list(range(len(detail\_df))) の結果を代入する。
\item speed\_exec\_present に 'v\_exec\_kmh' in detail\_df.columns の結果を代入する。
\item speed\_series に detail\_df.get('v\_exec\_kmh', detail\_df.get('v\_cmd\_kmh', pd.Series(dtype=float))) の結果を代入する。
\item speed\_svg に build\_svg\_chart(x\_index, speed\_series, color='\#0f766e', label='speed exec [km/h]' if speed\_exec\_present else 'speed cmd [km/h]') の結果を代入する。
\item speed\_cmd\_svg に '' の結果を代入する。
\item 条件 speed\_exec\_present を判定し、真なら内部処理を行う。
\item   speed\_cmd\_svg に build\_svg\_chart(x\_index, detail\_df.get('v\_cmd\_kmh', pd.Series(dtype=float)), color='\#475569', label='speed cmd [km/h]') の結果を代入する。
\item soc\_svg に build\_svg\_chart(x\_index, detail\_df.get('soc', pd.Series(dtype=float)), color='\#b45309', label='soc [-]') の結果を代入する。
\item pack\_svg に build\_svg\_chart(x\_index, detail\_df.get('P\_pack', pd.Series(dtype=float)), color='\#b91c1c', label='pack power [W]') の結果を代入する。
\item solar\_svg に build\_svg\_chart(x\_index, detail\_df.get('P\_pv', pd.Series(dtype=float)), color='\#2563eb', label='pv power [W]') の結果を代入する。
\item wind\_svg に build\_svg\_chart(x\_index, detail\_df.get('headwind\_ms', pd.Series(dtype=float)), color='\#7c3aed', label='headwind [m/s]') の結果を代入する。
\item slope\_svg に build\_svg\_chart(x\_index, detail\_df.get('slope\_pct', pd.Series(dtype=float)), color='\#57534e', label='slope [\%]') の結果を代入する。
\item summary\_cards に [('Finish reached', 'yes' if summary.get('finish\_reached') else 'no'), ('Final distance [km]', format\_metric(summary.get('final\_distance\_km'), 1)), ('Race progress [\%]', format\_metric(summary.get('race\_progress\_pct'), 1)), ('Final SoC [-]', format\_metric(summary.get('final\_soc'), 3)), ('Min SoC [-]', format\_metric(summary.get('min\_soc'), 3)), ('Avg speed [km/h]', format\_metric(summary.get('avg\_speed\_kmh'), 1)), ('Elapsed [h]', format\_metric(summary.get('elapsed\_hours'), 2)), ('Exec model', 'on' if summary.get('execution\_model\_enabled') else 'off'), ('Mean |v\_cmd-v\_exec| [km/h]', format\_metric(summary.get('mean\_tracking\_error\_kmh'), 2)), ('P95 |v\_cmd-v\_exec| [km/h]', format\_metric(summary.get('p95\_tracking\_error\_kmh'), 2)), ('Forecast mode', str(summary.get('forecast\_mode', '--'))), ('Overrides', str(summary.get('override\_count', 0)))] の結果を代入する。
\item overrides\_html に ''.join((f"<tr><td>\{html.escape(str(item.get('key')))\}</td><td><code>\{html.escape(json.dumps(item.get('value'), ensure\_ascii=False))\}</code></td></tr>" for item in summary.get('overrides', []))) or '<tr><td colspan="2">none</td></tr>' の結果を代入する。
\item params\_html に ''.join((f'<tr><td>\{html.escape(str(key))\}</td><td>\{html.escape(str(value))\}</td></tr>' for key, value in params\_rows)) の結果を代入する。
\item cards\_html に ''.join((f'<div class="metric"><div class="metric-label">\{html.escape(label)\}</div><div class="metric-value">\{html.escape(str(value))\}</div></div>' for label, value in summary\_cards)) の結果を代入する。
\item charts に [speed\_svg] の結果を代入する。
\item 条件 speed\_cmd\_svg を判定し、真なら内部処理を行う。
\item   charts.append(...) を実行する。
\item charts.extend(...) を実行する。
\item charts\_html に ''.join(charts) の結果を代入する。
\item warning\_text に html.escape(str(summary.get('warning', ''))) の結果を代入する。
\item html\_text に f"""<!doctype html>\textbackslash{}n<html lang="en">\textbackslash{}n<head>\textbackslash{}n  <meta charset="utf-8">\textbackslash{}n  <title>Solar Simulation Report</title>\textbackslash{}n  <style>\textbackslash{}n    body \{\{ margin: 0; font-family: "Segoe UI", sans-serif; background: \#f5efe3; color: \#1f2933; \}\}\textbackslash{}n    .wrap \{\{ max-width: 1180px; margin: 0 auto; padding: 24px; \}\}\textbackslash{}n    .hero \{\{ background: linear-gradient(135deg, \#fffaf0, \#efe6d2); border: 1px solid \#d9cfbd; border-radius: 24px; padding: 22px; \}\}\textbackslash{}n    h1, h2 \{\{ margin: 0 0 12px; \}\}\textbackslash{}n    p \{\{ margin: 8px 0; \}\}\textbackslash{}n    .muted \{\{ color: \#625a4e; \}\}\textbackslash{}n    .grid \{\{ display: grid; grid-template-columns: repeat(auto-fit, minmax(220px, 1fr)); gap: 12px; margin-top: 18px; \}\}\textbackslash{}n    .metric \{\{ background: \#fffdf8; border: 1px solid \#ddd2be; border-radius: 18px; padding: 14px; \}\}\textbackslash{}n    .metric-label \{\{ font-size: 12px; color: \#6f665b; text-transform: uppercase; letter-spacing: 0.06em; \}\}\textbackslash{}n    .metric-value \{\{ margin-top: 6px; font-size: 24px; font-weight: 700; \}\}\textbackslash{}n    .section \{\{ margin-top: 20px; background: \#fffaf2; border: 1px solid \#ddd2be; border-radius: 24px; padding: 18px; \}\}\textbackslash{}n    .charts \{\{ display: grid; grid-template-columns: 1fr; gap: 16px; \}\}\textbackslash{}n    .chart-svg \{\{ width: 100\%; height: auto; display: block; \}\}\textbackslash{}n    table \{\{ width: 100\%; border-collapse: collapse; \}\}\textbackslash{}n    th, td \{\{ border-bottom: 1px solid \#e6ddcf; padding: 8px 10px; text-align: left; vertical-align: top; \}\}\textbackslash{}n    th \{\{ font-size: 12px; color: \#6f665b; text-transform: uppercase; letter-spacing: 0.06em; \}\}\textbackslash{}n    code \{\{ font-family: Consolas, monospace; font-size: 12px; \}\}\textbackslash{}n    .warning \{\{ margin-top: 10px; padding: 12px; border-radius: 14px; background: \#fff0e1; color: \#8a3b12; \}\}\textbackslash{}n  </style>\textbackslash{}n</head>\textbackslash{}n<body>\textbackslash{}n  <div class="wrap">\textbackslash{}n    <div class="hero">\textbackslash{}n      <h1>Solar Simulation Report</h1>\textbackslash{}n      <p class="muted">profile: \{html.escape(str(summary.get('profile\_name', '--')))\} | generated: \{html.escape(str(summary.get('generated\_at\_utc', '--')))\}</p>\textbackslash{}n      <p class="muted">csv: \{html.escape(str(summary.get('out\_csv', '--')))\} | detail: \{html.escape(str(summary.get('detail\_csv', '--')))\}</p>\textbackslash{}n      \{(f'<div class="warning">\{warning\_text\}</div>' if warning\_text else '')\}\textbackslash{}n      <div class="grid">\{cards\_html\}</div>\textbackslash{}n    </div>\textbackslash{}n\textbackslash{}n    <div class="section">\textbackslash{}n      <h2>Key Charts</h2>\textbackslash{}n      <div class="charts">\textbackslash{}n        \{charts\_html\}\textbackslash{}n      </div>\textbackslash{}n    </div>\textbackslash{}n\textbackslash{}n    <div class="section">\textbackslash{}n      <h2>Overrides</h2>\textbackslash{}n      <table>\textbackslash{}n        <thead><tr><th>key</th><th>value</th></tr></thead>\textbackslash{}n        <tbody>\{overrides\_html\}</tbody>\textbackslash{}n      </table>\textbackslash{}n    </div>\textbackslash{}n\textbackslash{}n    <div class="section">\textbackslash{}n      <h2>Resolved Parameters</h2>\textbackslash{}n      <table>\textbackslash{}n        <thead><tr><th>parameter</th><th>value</th></tr></thead>\textbackslash{}n        <tbody>\{params\_html\}</tbody>\textbackslash{}n      </table>\textbackslash{}n    </div>\textbackslash{}n  </div>\textbackslash{}n</body>\textbackslash{}n</html>\textbackslash{}n""" の結果を代入する。
\item with 文で open(path, 'w', encoding='utf-8') を管理しながら処理する。
\item   f.write(...) を実行する。
            \end{enumerate}
            \paragraph{代表コード断片}
            \begin{lstlisting}[language=Python]
def write_simulation_report(path, summary, detail_df, params_rows):
    ensure_parent_dir(path)
    x_index = list(range(len(detail_df)))
    speed_exec_present = 'v_exec_kmh' in detail_df.columns
    speed_series = detail_df.get('v_exec_kmh', detail_df.get('v_cmd_kmh', pd.Series(dtype=float)))
    speed_svg = build_svg_chart(
        x_index,
        speed_series,
        color='#0f766e',
        label='speed exec [km/h]' if speed_exec_present else 'speed cmd [km/h]',
    )
    speed_cmd_svg = ''
    if speed_exec_present:
        speed_cmd_svg = build_svg_chart(
            x_index,
            detail_df.get('v_cmd_kmh', pd.Series(dtype=float)),
            color='#475569',
            label='speed cmd [km/h]',
        )
    soc_svg = build_svg_chart(x_index, detail_df.get('soc', pd.Series(dtype=float)), color='#b45309', label='soc [-]')
    pack_svg = build_svg_chart(x_index, detail_df.get('P_pack', pd.Series(dtype=float)), color='#b91c1c', label='pack power [W]')
    solar_svg = build_svg_chart(x_index, detail_df.get('P_pv', pd.Series(dtype=float)), color='#2563eb', label='pv power [W]')
    wind_svg = build_svg_chart(x_index, detail_df.get('headwind_ms', pd.Series(dtype=float)), color='#7c3aed', label='headwind [m/s]')
    slope_svg = build_svg_chart(x_index, detail_df.get('slope_pct', pd.Series(dtype=float)), color='#57534e', label='slope [%]')

    summary_cards = [
        ('Finish reached', 'yes' if summary.get('finish_reached') else 'no'),
        ('Final distance [km]', format_metric(summary.get('final_distance_km'), 1)),
        ('Race progress [%]', format_metric(summary.get('race_progress_pct'), 1)),
        ('Final SoC [-]', format_metric(summary.get('final_soc'), 3)),
        ('Min SoC [-]', format_metric(summary.get('min_soc'), 3)),
        ('Avg speed [km/h]', format_metric(summary.get('avg_speed_kmh'), 1)),
        ('Elapsed [h]', format_metric(summary.get('elapsed_hours'), 2)),
        ('Exec model', 'on' if summary.get('execution_model_enabled') else 'off'),
        ('Mean |v_cmd-v_exec| [km/h]', format_metric(summary.get('mean_tracking_error_kmh'), 2)),
...
\end{lstlisting}

\subsection{L1288 関数 mpc\_solve}
            \begin{itemize}[leftmargin=1.6em]
              \item 行範囲: L1288--L1399
              \item 所属: module直下
              \item 定義: \path|mpc_solve(model, data, z0, Tb0, s0_km, v0_kmh, v_max_kmh, term_soc_min, w_dv, w_dv_limit, dv_max_kmhps, w_T, w_speed_limit, w_current, speed_profile, soc_target, soc_band, w_soc_target, w_soc_band, schedule, soc_day_end_max, w_soc_day_max, soc_finish_target, soc_finish_tol, w_soc_progress, w_soc_terminal, race_km, soc_day_end_target, soc_day_end_tol, w_soc_day_track)|
              \item docstring: 明示されていない。
              \item このブロックが直接呼ぶ主な関数・メソッド: abs, array, clip, current\_drive\_window, dict, electrical\_balance, float, get, get\_profile\_val, is\_drive\_time, len, list
              \item 戻り値の要点: float(np.clip(v0\_kmh, 0.0, v\_max\_kmh)) / v0\_kmh / x * x / J
              \item 主なローカル代入名: I, J, Np, P\_pack, Tb, Tb\_next, V, bounds, d, dv, dv\_max\_msps, inertial\_power\_w, k, lb, loss\_int, out, p, prog, res, s\_km
              \item 読み取る主なインスタンス属性: 特になし。
              \item 更新する主なインスタンス属性: 特になし。
              \item 明示的に送出する例外: 明示的raiseはない。
              \item 制御構造の規模: 条件分岐 14、ループ 1、try 0
            \end{itemize}
            \paragraph{この定義を読むためのPython構文}
            \begin{itemize}[leftmargin=1.6em]
            \item 特別な構文上の補足は少ない。
            \end{itemize}
            \paragraph{上から順の処理}
            \begin{enumerate}[leftmargin=1.8em]
            \item p に model.p の結果を代入する。
\item Np に len(data) の結果を代入する。
\item 条件 Np <= 0 を判定し、真なら内部処理を行う。
\item   v0\_kmh を返す。
\item v0\_ms に v0\_kmh / 3.6 の結果を代入する。
\item x0 に np.array([v0\_ms] * Np, dtype=float) の結果を代入する。
\item lb に np.zeros(Np, dtype=float) の結果を代入する。
\item ub に np.ones(Np, dtype=float) * (v\_max\_kmh / 3.6) の結果を代入する。
\item dv\_max\_msps に dv\_max\_kmhps / 3.6 の結果を代入する。
\item 関数 quad\_penalty を定義する。
\item 関数 cost を定義する。
\item ImportFrom 文を実行する。
\item bounds に list(zip(lb, ub)) の結果を代入する。
\item res に minimize(cost, x0, method='L-BFGS-B', bounds=bounds, options=dict(maxiter=150)) の結果を代入する。
\item x\_best に select\_optimized\_vector(res, x0, label='lower solve') の結果を代入する。
\item v0\_kmh に float(x\_best[0]) * 3.6 の結果を代入する。
\item float(np.clip(v0\_kmh, 0.0, v\_max\_kmh)) を返す。
            \end{enumerate}
            \paragraph{代表コード断片}
            \begin{lstlisting}[language=Python]
def mpc_solve(model, data, z0, Tb0, s0_km, v0_kmh, v_max_kmh, term_soc_min,
              w_dv, w_dv_limit, dv_max_kmhps, w_T, w_speed_limit, w_current, speed_profile,
              soc_target, soc_band, w_soc_target, w_soc_band, schedule, soc_day_end_max, w_soc_day_max,
              soc_finish_target, soc_finish_tol, w_soc_progress, w_soc_terminal, race_km,
              soc_day_end_target, soc_day_end_tol, w_soc_day_track):
    p = model.p
    Np = len(data)
    if Np <= 0:
        return v0_kmh

    v0_ms = v0_kmh / 3.6
    x0 = np.array([v0_ms] * Np, dtype=float)
    lb = np.zeros(Np, dtype=float)
    ub = np.ones(Np, dtype=float) * (v_max_kmh / 3.6)

    dv_max_msps = dv_max_kmhps / 3.6

    def quad_penalty(x, cap=1.0e3):
        if x <= 0.0:
            return 0.0
        if x > cap:
            x = cap
        return x * x

    def cost(v):
        z = float(z0)
        Tb = float(Tb0)
        s_km = float(s0_km)
        v_prev = float(v0_ms)
        J = 0.0
        for k in range(Np):
            d = data[k]
            v_k = float(v[k])
            inertial_power_w = 0.5 * p.m * (v_k * v_k - v_prev * v_prev) / max(p.dt, 1.0e-9)
            out = model.electrical_balance(
...
\end{lstlisting}

\subsection{L1305 関数 mpc\_solve.quad\_penalty}
            \begin{itemize}[leftmargin=1.6em]
              \item 行範囲: L1305--L1310
              \item 所属: \path|mpc_solve|
              \item 定義: \path|quad_penalty(x, cap = 1000.0)|
              \item docstring: 明示されていない。
              \item このブロックが直接呼ぶ主な関数・メソッド: 特に明示されていない。
              \item 戻り値の要点: x * x / 0.0
              \item 主なローカル代入名: x
              \item 読み取る主なインスタンス属性: 特になし。
              \item 更新する主なインスタンス属性: 特になし。
              \item 明示的に送出する例外: 明示的raiseはない。
              \item 制御構造の規模: 条件分岐 2、ループ 0、try 0
            \end{itemize}
            \paragraph{この定義を読むためのPython構文}
            \begin{itemize}[leftmargin=1.6em]
            \item 特別な構文上の補足は少ない。
            \end{itemize}
            \paragraph{上から順の処理}
            \begin{enumerate}[leftmargin=1.8em]
            \item 条件 x <= 0.0 を判定し、真なら内部処理を行う。
\item   0.0 を返す。
\item 条件 x > cap を判定し、真なら内部処理を行う。
\item   x に cap の結果を代入する。
\item x * x を返す。
            \end{enumerate}
            \paragraph{代表コード断片}
            \begin{lstlisting}[language=Python]
    def quad_penalty(x, cap=1.0e3):
        if x <= 0.0:
            return 0.0
        if x > cap:
            x = cap
        return x * x
\end{lstlisting}

\subsection{L1312 関数 mpc\_solve.cost}
            \begin{itemize}[leftmargin=1.6em]
              \item 行範囲: L1312--L1392
              \item 所属: \path|mpc_solve|
              \item 定義: \path|cost(v)|
              \item docstring: 明示されていない。
              \item このブロックが直接呼ぶ主な関数・メソッド: abs, current\_drive\_window, electrical\_balance, float, get, get\_profile\_val, is\_drive\_time, max, min, quad\_penalty, range, soc\_step
              \item 戻り値の要点: J
              \item 主なローカル代入名: I, J, P\_pack, Tb, Tb\_next, V, d, dv, inertial\_power\_w, k, loss\_int, out, prog, s\_km, soc\_line, t\_end, t\_next, t\_start, v\_k, v\_prev
              \item 読み取る主なインスタンス属性: 特になし。
              \item 更新する主なインスタンス属性: 特になし。
              \item 明示的に送出する例外: 明示的raiseはない。
              \item 制御構造の規模: 条件分岐 11、ループ 1、try 0
            \end{itemize}
            \paragraph{この定義を読むためのPython構文}
            \begin{itemize}[leftmargin=1.6em]
            \item 特別な構文上の補足は少ない。
            \end{itemize}
            \paragraph{上から順の処理}
            \begin{enumerate}[leftmargin=1.8em]
            \item z に float(z0) の結果を代入する。
\item Tb に float(Tb0) の結果を代入する。
\item s\_km に float(s0\_km) の結果を代入する。
\item v\_prev に float(v0\_ms) の結果を代入する。
\item J に 0.0 の結果を代入する。
\item range(Np) を順に走査し、各要素を k に入れて処理する。
\item   d に data[k] の結果を代入する。
\item   v\_k に float(v[k]) の結果を代入する。
\item   inertial\_power\_w に 0.5 * p.m * (v\_k * v\_k - v\_prev * v\_prev) / max(p.dt, 1e-09) の結果を代入する。
\item   out に model.electrical\_balance(v\_k, d['slope\_pct'], z, Tb, d['G\_poa'], d['Tcell\_C'], headwind\_ms=d.get('headwind\_ms', 0.0), inertial\_power\_w=inertial\_power\_w, ambient\_temp\_c=d.get('Tamb\_C'), elevation\_m=d.get('elevation\_m', 0.0)) の結果を代入する。
\item   I に float(out['I']) の結果を代入する。
\item   V に float(out['V']) の結果を代入する。
\item   P\_pack に float(out['P\_pack']) の結果を代入する。
\item   loss\_int に float(out['losses\_int']) の結果を代入する。
\item   z\_next に model.soc\_step(z, P\_pack, p.dt, current\_a=I, Tbat\_C=Tb) の結果を代入する。
\item   Tb\_next に Tb + p.dt / 1800.0 * (d['Tamb\_C'] - Tb) + loss\_int * p.dt / 50000.0 の結果を代入する。
\item   s\_km に s\_km + v\_k * (p.dt / 1000.0) の結果を代入する。
\item   J を Add で更新する。
\item   J を Add で更新する。
\item   条件 schedule is not None and soc\_day\_end\_target > 0.0 and ('t\_utc' in d) を判定し、真なら内部処理を行う。
\item     win に schedule.current\_drive\_window(d['t\_utc']) の結果を代入する。
\item     条件 win is not None を判定し、真なら内部処理を行う。
\item       (t\_start, t\_end) に win の結果を代入する。
\item       条件 t\_end > t\_start を判定し、真なら内部処理を行う。
\item   条件 schedule is not None and soc\_day\_end\_max > 0.0 and ('t\_utc' in d) を判定し、真なら内部処理を行う。
\item     t\_next に d['t\_utc'] + timedelta(seconds=p.dt) の結果を代入する。
\item     条件 schedule.is\_drive\_time(d['t\_utc']) and (not schedule.is\_drive\_time(t\_next)) を判定し、真なら内部処理を行う。
\item       J を Add で更新する。
\item   dv に (v\_k - v\_prev) / max(p.dt, 0.001) の結果を代入する。
\item   条件 dv\_max\_msps > 0.0 を判定し、真なら内部処理を行う。
\item     J を Add で更新する。
\item   vmax\_local に v\_max\_kmh の結果を代入する。
\item   条件 speed\_profile is not None を判定し、真なら内部処理を行う。
\item     vmax\_local に get\_profile\_val(speed\_profile, s\_km, 'v\_max\_kmh', v\_max\_kmh) の結果を代入する。
\item   条件 vmax\_local < v\_max\_kmh を判定し、真なら内部処理を行う。
\item     J を Add で更新する。
\item   J を Add で更新する。
\item   J を Add で更新する。
\item   J を Add で更新する。
\item   J を Add で更新する。
\item   J を Add で更新する。
\item   J を Add で更新する。
\item   J を Add で更新する。
\item   J を Add で更新する。
\item   (z, Tb) に (z\_next, Tb\_next) の結果を代入する。
\item   v\_prev に v\_k の結果を代入する。
\item J を Add で更新する。
\item 条件 soc\_finish\_target > 0.0 を判定し、真なら内部処理を行う。
\item   J を Add で更新する。
\item J を返す。
            \end{enumerate}
            \paragraph{代表コード断片}
            \begin{lstlisting}[language=Python]
    def cost(v):
        z = float(z0)
        Tb = float(Tb0)
        s_km = float(s0_km)
        v_prev = float(v0_ms)
        J = 0.0
        for k in range(Np):
            d = data[k]
            v_k = float(v[k])
            inertial_power_w = 0.5 * p.m * (v_k * v_k - v_prev * v_prev) / max(p.dt, 1.0e-9)
            out = model.electrical_balance(
                v_k,
                d['slope_pct'],
                z,
                Tb,
                d['G_poa'],
                d['Tcell_C'],
                headwind_ms=d.get('headwind_ms', 0.0),
                inertial_power_w=inertial_power_w,
                ambient_temp_c=d.get('Tamb_C'),
                elevation_m=d.get('elevation_m', 0.0),
            )
            I = float(out['I'])
            V = float(out['V'])
            P_pack = float(out['P_pack'])
            loss_int = float(out['losses_int'])

            z_next = model.soc_step(
                z,
                P_pack,
                p.dt,
                current_a=I,
                Tbat_C=Tb,
            )
            Tb_next = Tb + (p.dt / 1800.0) * (d['Tamb_C'] - Tb) + (loss_int * p.dt) / 50000.0
...
\end{lstlisting}

\subsection{L1402 関数 soc\_guard\_speed}
            \begin{itemize}[leftmargin=1.6em]
              \item 行範囲: L1402--L1448
              \item 所属: module直下
              \item 定義: \path|soc_guard_speed(model, v_kmh, z, Tb, s_km, d0, route_profile, mode, soc_guard)|
              \item docstring: 明示されていない。
              \item このブロックが直接呼ぶ主な関数・メソッド: electrical\_balance, float, get, get\_profile\_val, lower, max, range, soc\_step, str, z\_next\_for
              \item 戻り値の要点: float(lo) / model.soc\_step(z, P\_pack, model.p.dt, current\_a=float(out['I']), Tbat\_C=Tb) / float(lo) / v\_kmh
              \item 主なローカル代入名: P\_pack, \_, elevation\_m, headwind\_ms, hi, lo, mid, mode, out, slope\_pct, target
              \item 読み取る主なインスタンス属性: 特になし。
              \item 更新する主なインスタンス属性: 特になし。
              \item 明示的に送出する例外: 明示的raiseはない。
              \item 制御構造の規模: 条件分岐 6、ループ 2、try 0
            \end{itemize}
            \paragraph{この定義を読むためのPython構文}
            \begin{itemize}[leftmargin=1.6em]
            \item 特別な構文上の補足は少ない。
            \end{itemize}
            \paragraph{上から順の処理}
            \begin{enumerate}[leftmargin=1.8em]
            \item mode に str(mode).lower() の結果を代入する。
\item target に model.p.soc\_min + soc\_guard の結果を代入する。
\item slope\_pct に get\_profile\_val(route\_profile, s\_km, 'slope\_pct', d0.get('slope\_pct', 0.0)) の結果を代入する。
\item headwind\_ms に d0.get('headwind\_ms', 0.0) の結果を代入する。
\item elevation\_m に get\_profile\_val(route\_profile, s\_km, 'elev\_m', d0.get('elevation\_m', 0.0)) の結果を代入する。
\item 関数 z\_next\_for を定義する。
\item 条件 z <= target を判定し、真なら内部処理を行う。
\item   条件 mode == 'stop' を判定し、真なら内部処理を行う。
\item     0.0 を返す。
\item   条件 mode != 'pv\_only' を判定し、真なら内部処理を行う。
\item     v\_kmh を返す。
\item   lo に 0.0 の結果を代入する。
\item   hi に max(0.0, float(v\_kmh)) の結果を代入する。
\item   range(20) を順に走査し、各要素を \_ に入れて処理する。
\item     mid に 0.5 * (lo + hi) の結果を代入する。
\item     条件 z\_next\_for(mid) < z を判定し、真なら内部処理を行う。
\item       hi に mid の結果を代入する。
\item       上の条件が偽の場合:
\item       lo に mid の結果を代入する。
\item   float(lo) を返す。
\item 条件 z\_next\_for(v\_kmh) >= target を判定し、真なら内部処理を行う。
\item   v\_kmh を返す。
\item lo に 0.0 の結果を代入する。
\item hi に max(0.0, float(v\_kmh)) の結果を代入する。
\item range(25) を順に走査し、各要素を \_ に入れて処理する。
\item   mid に 0.5 * (lo + hi) の結果を代入する。
\item   条件 z\_next\_for(mid) < target を判定し、真なら内部処理を行う。
\item     hi に mid の結果を代入する。
\item     上の条件が偽の場合:
\item     lo に mid の結果を代入する。
\item float(lo) を返す。
            \end{enumerate}
            \paragraph{代表コード断片}
            \begin{lstlisting}[language=Python]
def soc_guard_speed(model, v_kmh, z, Tb, s_km, d0, route_profile, mode, soc_guard):
    mode = str(mode).lower()
    target = model.p.soc_min + soc_guard
    slope_pct = get_profile_val(route_profile, s_km, 'slope_pct', d0.get('slope_pct', 0.0))
    headwind_ms = d0.get('headwind_ms', 0.0)
    elevation_m = get_profile_val(route_profile, s_km, 'elev_m', d0.get('elevation_m', 0.0))

    def z_next_for(v_kmh_local):
        out = model.electrical_balance(
            v_kmh_local / 3.6, slope_pct, z, Tb, d0['G_poa'], d0['Tcell_C'],
            headwind_ms=headwind_ms, ambient_temp_c=d0.get('Tamb_C'), elevation_m=elevation_m,
        )
        P_pack = float(out['P_pack'])
        return model.soc_step(
            z,
            P_pack,
            model.p.dt,
            current_a=float(out['I']),
            Tbat_C=Tb,
        )

    if z <= target:
        if mode == 'stop':
            return 0.0
        if mode != 'pv_only':
            return v_kmh
        lo = 0.0
        hi = max(0.0, float(v_kmh))
        for _ in range(20):
            mid = 0.5 * (lo + hi)
            if z_next_for(mid) < z:
                hi = mid
            else:
                lo = mid
        return float(lo)
...
\end{lstlisting}

\subsection{L1409 関数 soc\_guard\_speed.z\_next\_for}
            \begin{itemize}[leftmargin=1.6em]
              \item 行範囲: L1409--L1421
              \item 所属: \path|soc_guard_speed|
              \item 定義: \path|z_next_for(v_kmh_local)|
              \item docstring: 明示されていない。
              \item このブロックが直接呼ぶ主な関数・メソッド: electrical\_balance, float, get, soc\_step
              \item 戻り値の要点: model.soc\_step(z, P\_pack, model.p.dt, current\_a=float(out['I']), Tbat\_C=Tb)
              \item 主なローカル代入名: P\_pack, out
              \item 読み取る主なインスタンス属性: 特になし。
              \item 更新する主なインスタンス属性: 特になし。
              \item 明示的に送出する例外: 明示的raiseはない。
              \item 制御構造の規模: 条件分岐 0、ループ 0、try 0
            \end{itemize}
            \paragraph{この定義を読むためのPython構文}
            \begin{itemize}[leftmargin=1.6em]
            \item 特別な構文上の補足は少ない。
            \end{itemize}
            \paragraph{上から順の処理}
            \begin{enumerate}[leftmargin=1.8em]
            \item out に model.electrical\_balance(v\_kmh\_local / 3.6, slope\_pct, z, Tb, d0['G\_poa'], d0['Tcell\_C'], headwind\_ms=headwind\_ms, ambient\_temp\_c=d0.get('Tamb\_C'), elevation\_m=elevation\_m) の結果を代入する。
\item P\_pack に float(out['P\_pack']) の結果を代入する。
\item model.soc\_step(z, P\_pack, model.p.dt, current\_a=float(out['I']), Tbat\_C=Tb) を返す。
            \end{enumerate}
            \paragraph{代表コード断片}
            \begin{lstlisting}[language=Python]
    def z_next_for(v_kmh_local):
        out = model.electrical_balance(
            v_kmh_local / 3.6, slope_pct, z, Tb, d0['G_poa'], d0['Tcell_C'],
            headwind_ms=headwind_ms, ambient_temp_c=d0.get('Tamb_C'), elevation_m=elevation_m,
        )
        P_pack = float(out['P_pack'])
        return model.soc_step(
            z,
            P_pack,
            model.p.dt,
            current_a=float(out['I']),
            Tbat_C=Tb,
        )
\end{lstlisting}

\subsection{L1451 関数 soc\_day\_guard\_speed}
            \begin{itemize}[leftmargin=1.6em]
              \item 行範囲: L1451--L1482
              \item 所属: module直下
              \item 定義: \path|soc_day_guard_speed(model, v_kmh, z, Tb, s_km, d0, route_profile, target_soc, tol)|
              \item docstring: 明示されていない。
              \item このブロックが直接呼ぶ主な関数・メソッド: electrical\_balance, float, get, get\_profile\_val, max, range, soc\_step, z\_next\_for
              \item 戻り値の要点: float(lo) / v\_kmh / model.soc\_step(z, P\_pack, model.p.dt, current\_a=float(out['I']), Tbat\_C=Tb) / v\_kmh
              \item 主なローカル代入名: P\_pack, \_, elevation\_m, headwind\_ms, hi, lo, mid, out, slope\_pct
              \item 読み取る主なインスタンス属性: 特になし。
              \item 更新する主なインスタンス属性: 特になし。
              \item 明示的に送出する例外: 明示的raiseはない。
              \item 制御構造の規模: 条件分岐 3、ループ 1、try 0
            \end{itemize}
            \paragraph{この定義を読むためのPython構文}
            \begin{itemize}[leftmargin=1.6em]
            \item 特別な構文上の補足は少ない。
            \end{itemize}
            \paragraph{上から順の処理}
            \begin{enumerate}[leftmargin=1.8em]
            \item 条件 target\_soc <= 0.0 を判定し、真なら内部処理を行う。
\item   v\_kmh を返す。
\item slope\_pct に get\_profile\_val(route\_profile, s\_km, 'slope\_pct', d0.get('slope\_pct', 0.0)) の結果を代入する。
\item headwind\_ms に d0.get('headwind\_ms', 0.0) の結果を代入する。
\item elevation\_m に get\_profile\_val(route\_profile, s\_km, 'elev\_m', d0.get('elevation\_m', 0.0)) の結果を代入する。
\item 関数 z\_next\_for を定義する。
\item 条件 z\_next\_for(v\_kmh) >= target\_soc - tol を判定し、真なら内部処理を行う。
\item   v\_kmh を返す。
\item lo に 0.0 の結果を代入する。
\item hi に max(0.0, float(v\_kmh)) の結果を代入する。
\item range(25) を順に走査し、各要素を \_ に入れて処理する。
\item   mid に 0.5 * (lo + hi) の結果を代入する。
\item   条件 z\_next\_for(mid) < target\_soc - tol を判定し、真なら内部処理を行う。
\item     hi に mid の結果を代入する。
\item     上の条件が偽の場合:
\item     lo に mid の結果を代入する。
\item float(lo) を返す。
            \end{enumerate}
            \paragraph{代表コード断片}
            \begin{lstlisting}[language=Python]
def soc_day_guard_speed(model, v_kmh, z, Tb, s_km, d0, route_profile, target_soc, tol):
    if target_soc <= 0.0:
        return v_kmh
    slope_pct = get_profile_val(route_profile, s_km, 'slope_pct', d0.get('slope_pct', 0.0))
    headwind_ms = d0.get('headwind_ms', 0.0)
    elevation_m = get_profile_val(route_profile, s_km, 'elev_m', d0.get('elevation_m', 0.0))

    def z_next_for(v_kmh_local):
        out = model.electrical_balance(
            v_kmh_local / 3.6, slope_pct, z, Tb, d0['G_poa'], d0['Tcell_C'],
            headwind_ms=headwind_ms, ambient_temp_c=d0.get('Tamb_C'), elevation_m=elevation_m,
        )
        P_pack = float(out['P_pack'])
        return model.soc_step(
            z,
            P_pack,
            model.p.dt,
            current_a=float(out['I']),
            Tbat_C=Tb,
        )

    if z_next_for(v_kmh) >= (target_soc - tol):
        return v_kmh
    lo = 0.0
    hi = max(0.0, float(v_kmh))
    for _ in range(25):
        mid = 0.5 * (lo + hi)
        if z_next_for(mid) < (target_soc - tol):
            hi = mid
        else:
            lo = mid
    return float(lo)
\end{lstlisting}

\subsection{L1458 関数 soc\_day\_guard\_speed.z\_next\_for}
            \begin{itemize}[leftmargin=1.6em]
              \item 行範囲: L1458--L1470
              \item 所属: \path|soc_day_guard_speed|
              \item 定義: \path|z_next_for(v_kmh_local)|
              \item docstring: 明示されていない。
              \item このブロックが直接呼ぶ主な関数・メソッド: electrical\_balance, float, get, soc\_step
              \item 戻り値の要点: model.soc\_step(z, P\_pack, model.p.dt, current\_a=float(out['I']), Tbat\_C=Tb)
              \item 主なローカル代入名: P\_pack, out
              \item 読み取る主なインスタンス属性: 特になし。
              \item 更新する主なインスタンス属性: 特になし。
              \item 明示的に送出する例外: 明示的raiseはない。
              \item 制御構造の規模: 条件分岐 0、ループ 0、try 0
            \end{itemize}
            \paragraph{この定義を読むためのPython構文}
            \begin{itemize}[leftmargin=1.6em]
            \item 特別な構文上の補足は少ない。
            \end{itemize}
            \paragraph{上から順の処理}
            \begin{enumerate}[leftmargin=1.8em]
            \item out に model.electrical\_balance(v\_kmh\_local / 3.6, slope\_pct, z, Tb, d0['G\_poa'], d0['Tcell\_C'], headwind\_ms=headwind\_ms, ambient\_temp\_c=d0.get('Tamb\_C'), elevation\_m=elevation\_m) の結果を代入する。
\item P\_pack に float(out['P\_pack']) の結果を代入する。
\item model.soc\_step(z, P\_pack, model.p.dt, current\_a=float(out['I']), Tbat\_C=Tb) を返す。
            \end{enumerate}
            \paragraph{代表コード断片}
            \begin{lstlisting}[language=Python]
    def z_next_for(v_kmh_local):
        out = model.electrical_balance(
            v_kmh_local / 3.6, slope_pct, z, Tb, d0['G_poa'], d0['Tcell_C'],
            headwind_ms=headwind_ms, ambient_temp_c=d0.get('Tamb_C'), elevation_m=elevation_m,
        )
        P_pack = float(out['P_pack'])
        return model.soc_step(
            z,
            P_pack,
            model.p.dt,
            current_a=float(out['I']),
            Tbat_C=Tb,
        )
\end{lstlisting}

\subsection{L1485 関数 main}
            \begin{itemize}[leftmargin=1.6em]
              \item 行範囲: L1485--L4016
              \item 所属: module直下
              \item 定義: \path|main()|
              \item docstring: 明示されていない。
              \item このブロックが直接呼ぶ主な関数・メソッド: ArgumentParser, DataFrame, DetailCsvStream, Params, Path, SmoothRateLimiter, SolarCarModel, Timestamp, ValueError, ZoneInfo, \_deep\_copy\_cfg, abs
              \item 戻り値の要点: choose\_integration\_step\_seconds(args.dt, forecast\_native\_dt\_sec) / t\_utc >= forecast\_start\_time / total\_sec / s0 + alpha * (s1 - s0)
              \item 主なローカル代入名: E\_batt, E\_pv, G\_control\_stop\_poa, G\_drive\_poa, G\_raw, G\_stop\_poa, I, J, Nc, Np, P\_allow, P\_mech\_wheel, P\_pack, P\_pv, Tb, Tb\_before, Tb\_cur, Tb\_next, Tb\_seed, Tb\_step
              \item 読み取る主なインスタンス属性: 特になし。
              \item 更新する主なインスタンス属性: 特になし。
              \item 明示的に送出する例外: ValueError('Either --profile\_yaml or --params\_yaml is required.'), ValueError('forecast\_csv, route\_profile\_csv, and stop\_yaml must be provided.'), ValueError('mpc.upper\_lock\_initial\_policy requires paths.initial\_upper\_policy\_csv')
              \item 制御構造の規模: 条件分岐 207、ループ 23、try 10
            \end{itemize}
            \paragraph{この定義を読むためのPython構文}
            \begin{itemize}[leftmargin=1.6em]
            \item lambdaは名前を付けずに短い関数オブジェクトを作る構文である。
\item 内包表記は、反復・条件・値生成を一つの式で記述してcollectionを作る。
\item with文はcontext managerに開始・終了処理を任せ、ファイルなどの資源を確実に解放する。
\item try文は例外が起き得る処理と、異常時または終了時の経路を分ける。
            \end{itemize}
            \paragraph{上から順の処理}
            \begin{enumerate}[leftmargin=1.8em]
            \item ap に argparse.ArgumentParser() の結果を代入する。
\item ap.add\_argument(...) を実行する。
\item ap.add\_argument(...) を実行する。
\item ap.add\_argument(...) を実行する。
\item ap.add\_argument(...) を実行する。
\item ap.add\_argument(...) を実行する。
\item ap.add\_argument(...) を実行する。
\item ap.add\_argument(...) を実行する。
\item ap.add\_argument(...) を実行する。
\item ap.add\_argument(...) を実行する。
\item ap.add\_argument(...) を実行する。
\item ap.add\_argument(...) を実行する。
\item ap.add\_argument(...) を実行する。
\item ap.add\_argument(...) を実行する。
\item ap.add\_argument(...) を実行する。
\item ap.add\_argument(...) を実行する。
\item ap.add\_argument(...) を実行する。
\item ap.add\_argument(...) を実行する。
\item ap.add\_argument(...) を実行する。
\item ap.add\_argument(...) を実行する。
\item ap.add\_argument(...) を実行する。
\item ap.add\_argument(...) を実行する。
\item ap.add\_argument(...) を実行する。
\item ap.add\_argument(...) を実行する。
\item ap.add\_argument(...) を実行する。
\item ap.add\_argument(...) を実行する。
\item ap.add\_argument(...) を実行する。
\item ap.add\_argument(...) を実行する。
\item ap.add\_argument(...) を実行する。
\item ap.add\_argument(...) を実行する。
\item ap.add\_argument(...) を実行する。
\item ap.add\_argument(...) を実行する。
\item ap.add\_argument(...) を実行する。
\item ap.add\_argument(...) を実行する。
\item ap.add\_argument(...) を実行する。
\item ap.add\_argument(...) を実行する。
\item ap.add\_argument(...) を実行する。
\item ap.add\_argument(...) を実行する。
\item ap.add\_argument(...) を実行する。
\item ap.add\_argument(...) を実行する。
\item ap.add\_argument(...) を実行する。
\item ap.add\_argument(...) を実行する。
\item ap.add\_argument(...) を実行する。
\item ap.add\_argument(...) を実行する。
\item ap.add\_argument(...) を実行する。
\item ap.add\_argument(...) を実行する。
\item ap.add\_argument(...) を実行する。
\item ap.add\_argument(...) を実行する。
\item ap.add\_argument(...) を実行する。
\item ap.add\_argument(...) を実行する。
\item ap.add\_argument(...) を実行する。
\item ap.add\_argument(...) を実行する。
\item ap.add\_argument(...) を実行する。
\item ap.add\_argument(...) を実行する。
\item ap.add\_argument(...) を実行する。
\item ap.add\_argument(...) を実行する。
\item ap.add\_argument(...) を実行する。
\item ap.add\_argument(...) を実行する。
\item ap.add\_argument(...) を実行する。
\item ap.add\_argument(...) を実行する。
\item ap.add\_argument(...) を実行する。
\item ap.add\_argument(...) を実行する。
\item ap.add\_argument(...) を実行する。
\item ap.add\_argument(...) を実行する。
\item ap.add\_argument(...) を実行する。
\item ap.add\_argument(...) を実行する。
\item ap.add\_argument(...) を実行する。
\item ap.add\_argument(...) を実行する。
\item ap.add\_argument(...) を実行する。
\item ap.add\_argument(...) を実行する。
\item ap.add\_argument(...) を実行する。
\item ap.add\_argument(...) を実行する。
\item ap.add\_argument(...) を実行する。
\item ap.add\_argument(...) を実行する。
\item ap.add\_argument(...) を実行する。
\item ap.add\_argument(...) を実行する。
\item ap.add\_argument(...) を実行する。
\item ap.add\_argument(...) を実行する。
\item ap.add\_argument(...) を実行する。
\item ap.add\_argument(...) を実行する。
            \end{enumerate}
            \paragraph{代表コード断片}
            \begin{lstlisting}[language=Python]
def main():
    ap = argparse.ArgumentParser()
    ap.add_argument('--profile_yaml', default='')
    ap.add_argument('--forecast_csv', default='')
    ap.add_argument('--forecast_fill_csv', default='')
    ap.add_argument('--progress_reference_csv', default='')
    ap.add_argument('--forecast_time_mode', default='auto')
    ap.add_argument('--forecast_time_tz', default='UTC')
    ap.add_argument('--route_profile_csv', default='')
    ap.add_argument('--speed_profile_csv', required=False, default='')
    ap.add_argument('--params_yaml', default='')
    ap.add_argument('--stop_yaml', default='')
    ap.add_argument('--drive_schedule_yaml', required=False, default='')
    ap.add_argument('--panel_eff_map', default='')
    ap.add_argument('--mppt_eff_map', default='')
    ap.add_argument('--ocv_soc_map', default='')
    ap.add_argument('--drive_eff_map', default='')
    ap.add_argument('--regen_eff_map', default='')
    ap.add_argument('--rint_map', default='')
    ap.add_argument('--drive_map_eco', default='')
    ap.add_argument('--drive_map_power', default='')
    ap.add_argument('--regen_map_eco', default='')
    ap.add_argument('--regen_map_power', default='')
    ap.add_argument('--dt', type=float, default=600.0)
    ap.add_argument('--horizon_steps', type=int, default=9)
    ap.add_argument('--soc0', type=float, default=0.99)
    ap.add_argument('--Tb0', type=float, default=30.0)
    ap.add_argument('--v0_kmh', type=float, default=40.0)
    ap.add_argument('--start_utc', default='')
    ap.add_argument('--forecast_start_time_utc', default='')
    ap.add_argument('--start_index', type=int, default=-1)
    ap.add_argument('--start_s_km', type=float, default=-1.0)
    ap.add_argument('--resume_csv', default='')
    ap.add_argument('--resume_s_km', type=float, default=-1.0)
    ap.add_argument('--out_csv', default='solar_sim_log.csv')
...
\end{lstlisting}

\subsection{L1852 関数 main.integration\_step\_seconds}
            \begin{itemize}[leftmargin=1.6em]
              \item 行範囲: L1852--L1857
              \item 所属: \path|main|
              \item 定義: \path|integration_step_seconds()|
              \item docstring: 明示されていない。
              \item このブロックが直接呼ぶ主な関数・メソッド: choose\_integration\_step\_seconds
              \item 戻り値の要点: choose\_integration\_step\_seconds(args.dt, forecast\_native\_dt\_sec)
              \item 主なローカル代入名: 特になし。
              \item 読み取る主なインスタンス属性: 特になし。
              \item 更新する主なインスタンス属性: 特になし。
              \item 明示的に送出する例外: 明示的raiseはない。
              \item 制御構造の規模: 条件分岐 0、ループ 0、try 0
            \end{itemize}
            \paragraph{この定義を読むためのPython構文}
            \begin{itemize}[leftmargin=1.6em]
            \item 特別な構文上の補足は少ない。
            \end{itemize}
            \paragraph{上から順の処理}
            \begin{enumerate}[leftmargin=1.8em]
            \item choose\_integration\_step\_seconds(args.dt, forecast\_native\_dt\_sec) を返す。
            \end{enumerate}
            \paragraph{代表コード断片}
            \begin{lstlisting}[language=Python]
    def integration_step_seconds():
        # Prediction must not integrate more coarsely than either the execution
        # state update or the available weather samples.  Using max() here made
        # hourly weather grids turn the planner into an hourly nonlinear state
        # integrator while execution still advanced every 10 minutes.
        return choose_integration_step_seconds(args.dt, forecast_native_dt_sec)
\end{lstlisting}

\subsection{L1859 関数 main.forecast\_has\_coverage}
            \begin{itemize}[leftmargin=1.6em]
              \item 行範囲: L1859--L1866
              \item 所属: \path|main|
              \item 定義: \path|forecast_has_coverage(t_utc: datetime) -> bool|
              \item docstring: 明示されていない。
              \item このブロックが直接呼ぶ主な関数・メソッド: len
              \item 戻り値の要点: t\_utc >= forecast\_start\_time / False / False / t\_utc >= forecast\_start\_data\_utc
              \item 主なローカル代入名: 特になし。
              \item 読み取る主なインスタンス属性: 特になし。
              \item 更新する主なインスタンス属性: 特になし。
              \item 明示的に送出する例外: 明示的raiseはない。
              \item 制御構造の規模: 条件分岐 3、ループ 0、try 0
            \end{itemize}
            \paragraph{この定義を読むためのPython構文}
            \begin{itemize}[leftmargin=1.6em]
            \item `->`以後は戻り値の型注釈であり、実行時の値を自動変換する命令ではない。
            \end{itemize}
            \paragraph{上から順の処理}
            \begin{enumerate}[leftmargin=1.8em]
            \item 条件 len(df) == 0 を判定し、真なら内部処理を行う。
\item   False を返す。
\item 条件 forecast\_coverage\_end\_utc is not None and t\_utc > forecast\_coverage\_end\_utc を判定し、真なら内部処理を行う。
\item   False を返す。
\item 条件 forecast\_mode == 'absolute' and forecast\_start\_data\_utc is not None を判定し、真なら内部処理を行う。
\item   t\_utc >= forecast\_start\_data\_utc を返す。
\item t\_utc >= forecast\_start\_time を返す。
            \end{enumerate}
            \paragraph{代表コード断片}
            \begin{lstlisting}[language=Python]
    def forecast_has_coverage(t_utc: datetime) -> bool:
        if len(df) == 0:
            return False
        if forecast_coverage_end_utc is not None and t_utc > forecast_coverage_end_utc:
            return False
        if forecast_mode == 'absolute' and forecast_start_data_utc is not None:
            return t_utc >= forecast_start_data_utc
        return t_utc >= forecast_start_time
\end{lstlisting}

\subsection{L1868 関数 main.integrate\_drive\_time\_between}
            \begin{itemize}[leftmargin=1.6em]
              \item 行範囲: L1868--L1883
              \item 所属: \path|main|
              \item 定義: \path|integrate_drive_time_between(t_start_utc: datetime, t_end_utc: datetime) -> float|
              \item docstring: 明示されていない。
              \item このブロックが直接呼ぶ主な関数・メソッド: integration\_step\_seconds, is\_drive\_time, max, min, timedelta, total\_seconds
              \item 戻り値の要点: total\_sec / 0.0 / (t\_end\_utc - t\_start\_utc).total\_seconds()
              \item 主なローカル代入名: dt\_local, dt\_sec, t\_cursor, total\_sec
              \item 読み取る主なインスタンス属性: 特になし。
              \item 更新する主なインスタンス属性: 特になし。
              \item 明示的に送出する例外: 明示的raiseはない。
              \item 制御構造の規模: 条件分岐 4、ループ 1、try 0
            \end{itemize}
            \paragraph{この定義を読むためのPython構文}
            \begin{itemize}[leftmargin=1.6em]
            \item `->`以後は戻り値の型注釈であり、実行時の値を自動変換する命令ではない。
            \end{itemize}
            \paragraph{上から順の処理}
            \begin{enumerate}[leftmargin=1.8em]
            \item 条件 t\_end\_utc <= t\_start\_utc を判定し、真なら内部処理を行う。
\item   0.0 を返す。
\item 条件 schedule is None を判定し、真なら内部処理を行う。
\item   (t\_end\_utc - t\_start\_utc).total\_seconds() を返す。
\item dt\_sec に integration\_step\_seconds() の結果を代入する。
\item total\_sec に 0.0 の結果を代入する。
\item t\_cursor に t\_start\_utc の結果を代入する。
\item 条件 t\_cursor < t\_end\_utc が成り立つ間くり返す。
\item   dt\_local に min(dt\_sec, max(0.0, (t\_end\_utc - t\_cursor).total\_seconds())) の結果を代入する。
\item   条件 dt\_local <= 0.0 を判定し、真なら内部処理を行う。
\item     Break 文を実行する。
\item   条件 schedule.is\_drive\_time(t\_cursor) を判定し、真なら内部処理を行う。
\item     total\_sec を Add で更新する。
\item   t\_cursor を Add で更新する。
\item total\_sec を返す。
            \end{enumerate}
            \paragraph{代表コード断片}
            \begin{lstlisting}[language=Python]
    def integrate_drive_time_between(t_start_utc: datetime, t_end_utc: datetime) -> float:
        if t_end_utc <= t_start_utc:
            return 0.0
        if schedule is None:
            return (t_end_utc - t_start_utc).total_seconds()
        dt_sec = integration_step_seconds()
        total_sec = 0.0
        t_cursor = t_start_utc
        while t_cursor < t_end_utc:
            dt_local = min(dt_sec, max(0.0, (t_end_utc - t_cursor).total_seconds()))
            if dt_local <= 0.0:
                break
            if schedule.is_drive_time(t_cursor):
                total_sec += dt_local
            t_cursor += timedelta(seconds=dt_local)
        return total_sec
\end{lstlisting}

\subsection{L1885 関数 main.reference\_value\_at\_time}
            \begin{itemize}[leftmargin=1.6em]
              \item 行範囲: L1885--L1905
              \item 所属: \path|main|
              \item 定義: \path!reference_value_at_time(t_utc: datetime, field: str) -> float | None!
              \item docstring: 明示されていない。
              \item このブロックが直接呼ぶ主な関数・メソッド: float, int, len, max, min, searchsorted, timestamp\_ns
              \item 戻り値の要点: s0 + alpha * (s1 - s0) / None / None / float(progress\_ref\_df[field].iloc[0])
              \item 主なローカル代入名: alpha, idx\_hi, idx\_lo, s0, s1, t0\_ns, t1\_ns, ts\_ns
              \item 読み取る主なインスタンス属性: 特になし。
              \item 更新する主なインスタンス属性: 特になし。
              \item 明示的に送出する例外: 明示的raiseはない。
              \item 制御構造の規模: 条件分岐 5、ループ 0、try 0
            \end{itemize}
            \paragraph{この定義を読むためのPython構文}
            \begin{itemize}[leftmargin=1.6em]
            \item `->`以後は戻り値の型注釈であり、実行時の値を自動変換する命令ではない。
            \end{itemize}
            \paragraph{上から順の処理}
            \begin{enumerate}[leftmargin=1.8em]
            \item 条件 progress\_ref\_df is None or progress\_ref\_time\_ns is None or len(progress\_ref\_df) == 0 を判定し、真なら内部処理を行う。
\item   None を返す。
\item 条件 field not in progress\_ref\_df.columns を判定し、真なら内部処理を行う。
\item   None を返す。
\item ts\_ns に timestamp\_ns(t\_utc) の結果を代入する。
\item idx\_hi に int(np.searchsorted(progress\_ref\_time\_ns, ts\_ns, side='left')) の結果を代入する。
\item 条件 idx\_hi <= 0 を判定し、真なら内部処理を行う。
\item   float(progress\_ref\_df[field].iloc[0]) を返す。
\item 条件 idx\_hi >= len(progress\_ref\_df) を判定し、真なら内部処理を行う。
\item   float(progress\_ref\_df[field].iloc[-1]) を返す。
\item idx\_lo に idx\_hi - 1 の結果を代入する。
\item t0\_ns に int(progress\_ref\_time\_ns[idx\_lo]) の結果を代入する。
\item t1\_ns に int(progress\_ref\_time\_ns[idx\_hi]) の結果を代入する。
\item s0 に float(progress\_ref\_df[field].iloc[idx\_lo]) の結果を代入する。
\item s1 に float(progress\_ref\_df[field].iloc[idx\_hi]) の結果を代入する。
\item 条件 t1\_ns <= t0\_ns を判定し、真なら内部処理を行う。
\item   s1 を返す。
\item alpha に float((ts\_ns - t0\_ns) / max(t1\_ns - t0\_ns, 1)) の結果を代入する。
\item alpha に max(0.0, min(1.0, alpha)) の結果を代入する。
\item s0 + alpha * (s1 - s0) を返す。
            \end{enumerate}
            \paragraph{代表コード断片}
            \begin{lstlisting}[language=Python]
    def reference_value_at_time(t_utc: datetime, field: str) -> float | None:
        if progress_ref_df is None or progress_ref_time_ns is None or len(progress_ref_df) == 0:
            return None
        if field not in progress_ref_df.columns:
            return None
        ts_ns = timestamp_ns(t_utc)
        idx_hi = int(np.searchsorted(progress_ref_time_ns, ts_ns, side='left'))
        if idx_hi <= 0:
            return float(progress_ref_df[field].iloc[0])
        if idx_hi >= len(progress_ref_df):
            return float(progress_ref_df[field].iloc[-1])
        idx_lo = idx_hi - 1
        t0_ns = int(progress_ref_time_ns[idx_lo])
        t1_ns = int(progress_ref_time_ns[idx_hi])
        s0 = float(progress_ref_df[field].iloc[idx_lo])
        s1 = float(progress_ref_df[field].iloc[idx_hi])
        if t1_ns <= t0_ns:
            return s1
        alpha = float((ts_ns - t0_ns) / max(t1_ns - t0_ns, 1))
        alpha = max(0.0, min(1.0, alpha))
        return s0 + alpha * (s1 - s0)
\end{lstlisting}

\subsection{L1907 関数 main.reference\_distance\_at\_time}
            \begin{itemize}[leftmargin=1.6em]
              \item 行範囲: L1907--L1908
              \item 所属: \path|main|
              \item 定義: \path!reference_distance_at_time(t_utc: datetime) -> float | None!
              \item docstring: 明示されていない。
              \item このブロックが直接呼ぶ主な関数・メソッド: reference\_value\_at\_time
              \item 戻り値の要点: reference\_value\_at\_time(t\_utc, 's\_km')
              \item 主なローカル代入名: 特になし。
              \item 読み取る主なインスタンス属性: 特になし。
              \item 更新する主なインスタンス属性: 特になし。
              \item 明示的に送出する例外: 明示的raiseはない。
              \item 制御構造の規模: 条件分岐 0、ループ 0、try 0
            \end{itemize}
            \paragraph{この定義を読むためのPython構文}
            \begin{itemize}[leftmargin=1.6em]
            \item `->`以後は戻り値の型注釈であり、実行時の値を自動変換する命令ではない。
            \end{itemize}
            \paragraph{上から順の処理}
            \begin{enumerate}[leftmargin=1.8em]
            \item reference\_value\_at\_time(t\_utc, 's\_km') を返す。
            \end{enumerate}
            \paragraph{代表コード断片}
            \begin{lstlisting}[language=Python]
    def reference_distance_at_time(t_utc: datetime) -> float | None:
        return reference_value_at_time(t_utc, 's_km')
\end{lstlisting}

\subsection{L1958 関数 main.time\_to\_index}
            \begin{itemize}[leftmargin=1.6em]
              \item 行範囲: L1958--L1969
              \item 所属: \path|main|
              \item 定義: \path|time_to_index(dt_utc)|
              \item docstring: 明示されていない。
              \item このブロックが直接呼ぶ主な関数・メソッド: all, astimezone, clip, datetime64, int, isna, len, replace, searchsorted
              \item 戻り値の要点: int(np.clip(idx, 0, len(df) - 1)) / 0
              \item 主なローカル代入名: idx, lookup\_naive, lookup\_utc, t\_series
              \item 読み取る主なインスタンス属性: 特になし。
              \item 更新する主なインスタンス属性: 特になし。
              \item 明示的に送出する例外: 明示的raiseはない。
              \item 制御構造の規模: 条件分岐 1、ループ 0、try 0
            \end{itemize}
            \paragraph{この定義を読むためのPython構文}
            \begin{itemize}[leftmargin=1.6em]
            \item 特別な構文上の補足は少ない。
            \end{itemize}
            \paragraph{上から順の処理}
            \begin{enumerate}[leftmargin=1.8em]
            \item 条件 'time' not in df.columns or df['time'].isna().all() を判定し、真なら内部処理を行う。
\item   0 を返す。
\item t\_series に df['time'].values の結果を代入する。
\item lookup\_utc に dt\_utc.replace(tzinfo=timezone.utc) if dt\_utc.tzinfo is None else dt\_utc.astimezone(timezone.utc) の結果を代入する。
\item lookup\_naive に lookup\_utc.replace(tzinfo=None) の結果を代入する。
\item idx に int(np.searchsorted(t\_series, np.datetime64(lookup\_naive, 'ns')) - 1) の結果を代入する。
\item int(np.clip(idx, 0, len(df) - 1)) を返す。
            \end{enumerate}
            \paragraph{代表コード断片}
            \begin{lstlisting}[language=Python]
    def time_to_index(dt_utc):
        if 'time' not in df.columns or df['time'].isna().all():
            return 0
        t_series = df['time'].values
        lookup_utc = (
            dt_utc.replace(tzinfo=timezone.utc)
            if dt_utc.tzinfo is None
            else dt_utc.astimezone(timezone.utc)
        )
        lookup_naive = lookup_utc.replace(tzinfo=None)
        idx = int(np.searchsorted(t_series, np.datetime64(lookup_naive, 'ns')) - 1)
        return int(np.clip(idx, 0, len(df) - 1))
\end{lstlisting}

\subsection{L2095 関数 main.remaining\_day\_budget}
            \begin{itemize}[leftmargin=1.6em]
              \item 行範囲: L2095--L2115
              \item 所属: \path|main|
              \item 定義: \path|remaining_day_budget(sim_t_local, _k_idx)|
              \item docstring: 明示されていない。
              \item このブロックが直接呼ぶ主な関数・メソッド: current\_drive\_window, float, forecast\_at\_time, forecast\_has\_coverage, integration\_step\_seconds, is\_drive\_time, max, min, pv\_power\_mppt, timedelta, total\_seconds
              \item 戻り値の要点: (E\_pv, t\_remain) / (None, None) / (None, None)
              \item 主なローカル代入名: E\_pv, P\_pv, \_, dt\_local, dt\_sample, env, t, t\_end, t\_remain, win
              \item 読み取る主なインスタンス属性: 特になし。
              \item 更新する主なインスタンス属性: 特になし。
              \item 明示的に送出する例外: 明示的raiseはない。
              \item 制御構造の規模: 条件分岐 3、ループ 1、try 0
            \end{itemize}
            \paragraph{この定義を読むためのPython構文}
            \begin{itemize}[leftmargin=1.6em]
            \item 特別な構文上の補足は少ない。
            \end{itemize}
            \paragraph{上から順の処理}
            \begin{enumerate}[leftmargin=1.8em]
            \item 条件 schedule is None or not schedule.is\_drive\_time(sim\_t\_local) を判定し、真なら内部処理を行う。
\item   (None, None) を返す。
\item win に schedule.current\_drive\_window(sim\_t\_local) の結果を代入する。
\item 条件 win is None を判定し、真なら内部処理を行う。
\item   (None, None) を返す。
\item (\_, t\_end) に win の結果を代入する。
\item dt\_sample に integration\_step\_seconds() の結果を代入する。
\item t に sim\_t\_local の結果を代入する。
\item E\_pv に 0.0 の結果を代入する。
\item t\_remain に 0.0 の結果を代入する。
\item 条件 t < t\_end and forecast\_has\_coverage(t) が成り立つ間くり返す。
\item   dt\_local に min(dt\_sample, max(0.0, (t\_end - t).total\_seconds())) の結果を代入する。
\item   条件 dt\_local <= 0.0 を判定し、真なら内部処理を行う。
\item     Break 文を実行する。
\item   env に forecast\_at\_time(t, s\_km, drive=True) の結果を代入する。
\item   P\_pv に float(model.pv\_power\_mppt(env['G\_poa'], env['Tcell\_C'])) の結果を代入する。
\item   E\_pv を Add で更新する。
\item   t\_remain を Add で更新する。
\item   t を Add で更新する。
\item (E\_pv, t\_remain) を返す。
            \end{enumerate}
            \paragraph{代表コード断片}
            \begin{lstlisting}[language=Python]
    def remaining_day_budget(sim_t_local, _k_idx):
        if schedule is None or not schedule.is_drive_time(sim_t_local):
            return None, None
        win = schedule.current_drive_window(sim_t_local)
        if win is None:
            return None, None
        _, t_end = win
        dt_sample = integration_step_seconds()
        t = sim_t_local
        E_pv = 0.0
        t_remain = 0.0
        while t < t_end and forecast_has_coverage(t):
            dt_local = min(dt_sample, max(0.0, (t_end - t).total_seconds()))
            if dt_local <= 0.0:
                break
            env = forecast_at_time(t, s_km, drive=True)
            P_pv = float(model.pv_power_mppt(env['G_poa'], env['Tcell_C']))
            E_pv += P_pv * (dt_local / 3600.0)
            t_remain += dt_local
            t += timedelta(seconds=dt_local)
        return E_pv, t_remain
\end{lstlisting}

\subsection{L2117 関数 main.budget\_speed\_limit}
            \begin{itemize}[leftmargin=1.6em]
              \item 行範囲: L2117--L2134
              \item 所属: \path|main|
              \item 定義: \path|budget_speed_limit(P_allow, d0)|
              \item docstring: 明示されていない。
              \item このブロックが直接呼ぶ主な関数・メソッド: electrical\_balance, float, get, range
              \item 戻り値の要点: float(lo) / None
              \item 主なローカル代入名: P\_pack, \_, hi, lo, mid, out
              \item 読み取る主なインスタンス属性: 特になし。
              \item 更新する主なインスタンス属性: 特になし。
              \item 明示的に送出する例外: 明示的raiseはない。
              \item 制御構造の規模: 条件分岐 2、ループ 1、try 0
            \end{itemize}
            \paragraph{この定義を読むためのPython構文}
            \begin{itemize}[leftmargin=1.6em]
            \item 特別な構文上の補足は少ない。
            \end{itemize}
            \paragraph{上から順の処理}
            \begin{enumerate}[leftmargin=1.8em]
            \item 条件 P\_allow is None を判定し、真なら内部処理を行う。
\item   None を返す。
\item lo に 0.0 の結果を代入する。
\item hi に v\_max\_kmh の結果を代入する。
\item range(24) を順に走査し、各要素を \_ に入れて処理する。
\item   mid に 0.5 * (lo + hi) の結果を代入する。
\item   out に model.electrical\_balance(mid / 3.6, d0['slope\_pct'], z, Tb, d0['G\_poa'], d0['Tcell\_C'], headwind\_ms=d0['headwind\_ms'], ambient\_temp\_c=d0.get('Tamb\_C'), elevation\_m=d0.get('elevation\_m', 0.0)) の結果を代入する。
\item   P\_pack に float(out['P\_pack']) の結果を代入する。
\item   条件 P\_pack > P\_allow を判定し、真なら内部処理を行う。
\item     hi に mid の結果を代入する。
\item     上の条件が偽の場合:
\item     lo に mid の結果を代入する。
\item float(lo) を返す。
            \end{enumerate}
            \paragraph{代表コード断片}
            \begin{lstlisting}[language=Python]
    def budget_speed_limit(P_allow, d0):
        if P_allow is None:
            return None
        lo = 0.0
        hi = v_max_kmh
        for _ in range(24):
            mid = 0.5 * (lo + hi)
            out = model.electrical_balance(
                mid / 3.6, d0['slope_pct'], z, Tb, d0['G_poa'], d0['Tcell_C'],
                headwind_ms=d0['headwind_ms'], ambient_temp_c=d0.get('Tamb_C'),
                elevation_m=d0.get('elevation_m', 0.0),
            )
            P_pack = float(out['P_pack'])
            if P_pack > P_allow:
                hi = mid
            else:
                lo = mid
        return float(lo)
\end{lstlisting}

\subsection{L2136 関数 main.forecast\_at\_time}
            \begin{itemize}[leftmargin=1.6em]
              \item 行範囲: L2136--L2234
              \item 所属: \path|main|
              \item 定義: \path|forecast_at_time(t_utc, s_query_km = None, drive = True)|
              \item docstring: 明示されていない。
              \item このブロックが直接呼ぶ主な関数・メソッド: clip, dict, float, get, get\_profile\_val, int, interp\_forecast\_grid, len, max, searchsorted, timestamp\_ns, total\_seconds
              \item 戻り値の要点: dict(G\_raw=G\_raw, G\_drive\_poa=G\_drive\_poa, G\_stop\_poa=G\_stop\_poa, G\_control\_stop\_poa=G\_control\_stop\_poa, G\_poa=G\_drive\_poa if drive else G\_stop\_poa, Tcell\_drive\_C=tcell\_drive, Tcell\_stop\_C=tcell\_stop, Tcell\_control\_stop\_C=tcell\_control\_stop, Tcell\_C=tcell\_drive if drive else tcell\_stop, Tamb\_C=float(row.get('Tamb\_C', 30.0)), headwind\_ms=float(row.get('headwind\_ms', 0.0)), elevation\_m=get\_profile\_val(route\_profile, float(s\_query\_km or 0.0), 'elev\_m', 0.0)) / dict(G\_raw=0.0, G\_drive\_poa=0.0, G\_stop\_poa=0.0, G\_control\_stop\_poa=0.0, G\_poa=0.0, Tcell\_drive\_C=40.0, Tcell\_stop\_C=40.0, Tcell\_control\_stop\_C=40.0, Tcell\_C=40.0, Tamb\_C=30.0, headwind\_ms=0.0, elevation\_m=0.0) / dict(G\_raw=G\_raw, G\_drive\_poa=G\_drive\_poa, G\_stop\_poa=G\_stop\_poa, G\_control\_stop\_poa=G\_control\_stop\_poa, G\_poa=G\_drive\_poa if drive else G\_stop\_poa, Tcell\_drive\_C=float(tcell\_drive), Tcell\_stop\_C=float(tcell\_stop), Tcell\_control\_stop\_C=float(tcell\_control\_stop), Tcell\_C=float(tcell\_drive if drive else tcell\_stop), Tamb\_C=float(tamb), headwind\_ms=float(headwind), elevation\_m=get\_profile\_val(route\_profile, float(s\_query\_km or 0.0), 'elev\_m', 0.0))
              \item 主なローカル代入名: G\_control\_stop\_poa, G\_drive\_poa, G\_raw, G\_stop\_poa, elapsed, ghi, headwind, idx, poa\_control\_stop\_raw, poa\_drive\_raw, poa\_stop\_ideal\_raw, poa\_stop\_raw, row, tamb, tcell\_control\_stop, tcell\_drive, tcell\_stop, tcell\_stop\_ideal
              \item 読み取る主なインスタンス属性: 特になし。
              \item 更新する主なインスタンス属性: 特になし。
              \item 明示的に送出する例外: 明示的raiseはない。
              \item 制御構造の規模: 条件分岐 3、ループ 0、try 0
            \end{itemize}
            \paragraph{この定義を読むためのPython構文}
            \begin{itemize}[leftmargin=1.6em]
            \item 特別な構文上の補足は少ない。
            \end{itemize}
            \paragraph{上から順の処理}
            \begin{enumerate}[leftmargin=1.8em]
            \item 条件 len(df) == 0 を判定し、真なら内部処理を行う。
\item   dict(G\_raw=0.0, G\_drive\_poa=0.0, G\_stop\_poa=0.0, G\_control\_stop\_poa=0.0, G\_poa=0.0, Tcell\_drive\_C=40.0, Tcell\_stop\_C=40.0, Tcell\_control\_stop\_C=40.0, Tcell\_C=40.0, Tamb\_C=30.0, headwind\_ms=0.0, elevation\_m=0.0) を返す。
\item 条件 forecast\_grid is not None を判定し、真なら内部処理を行う。
\item   ghi に interp\_forecast\_grid(forecast\_grid, 'GHI', t\_utc, s\_query\_km, 0.0) の結果を代入する。
\item   poa\_drive\_raw に interp\_forecast\_grid(forecast\_grid, 'POA\_drive', t\_utc, s\_query\_km, ghi) の結果を代入する。
\item   poa\_stop\_ideal\_raw に interp\_forecast\_grid(forecast\_grid, 'POA\_stop\_ideal', t\_utc, s\_query\_km, poa\_drive\_raw) の結果を代入する。
\item   poa\_stop\_raw に poa\_drive\_raw + args.stop\_tilt\_fraction * (poa\_stop\_ideal\_raw - poa\_drive\_raw) の結果を代入する。
\item   poa\_control\_stop\_raw に poa\_drive\_raw + args.control\_stop\_tilt\_fraction * (poa\_stop\_ideal\_raw - poa\_drive\_raw) の結果を代入する。
\item   tamb に interp\_forecast\_grid(forecast\_grid, 'Tamb\_C', t\_utc, s\_query\_km, 30.0) の結果を代入する。
\item   tcell\_drive に interp\_forecast\_grid(forecast\_grid, 'Tcell\_drive\_C', t\_utc, s\_query\_km, interp\_forecast\_grid(forecast\_grid, 'Tcell\_C', t\_utc, s\_query\_km, tamb)) の結果を代入する。
\item   tcell\_stop\_ideal に interp\_forecast\_grid(forecast\_grid, 'Tcell\_stop\_ideal\_C', t\_utc, s\_query\_km, tcell\_drive) の結果を代入する。
\item   tcell\_stop に tcell\_drive + args.stop\_tilt\_fraction * (tcell\_stop\_ideal - tcell\_drive) の結果を代入する。
\item   tcell\_control\_stop に tcell\_drive + args.control\_stop\_tilt\_fraction * (tcell\_stop\_ideal - tcell\_drive) の結果を代入する。
\item   headwind に interp\_forecast\_grid(forecast\_grid, 'headwind\_ms', t\_utc, s\_query\_km, 0.0) の結果を代入する。
\item   G\_raw に float(ghi) * args.solar\_gain の結果を代入する。
\item   G\_drive\_poa に float(poa\_drive\_raw) * args.solar\_gain * args.poa\_gain\_drive の結果を代入する。
\item   G\_stop\_poa に float(poa\_stop\_raw) * args.solar\_gain * args.poa\_gain\_stop の結果を代入する。
\item   G\_control\_stop\_poa に float(poa\_control\_stop\_raw) * args.solar\_gain * args.poa\_gain\_stop の結果を代入する。
\item   dict(G\_raw=G\_raw, G\_drive\_poa=G\_drive\_poa, G\_stop\_poa=G\_stop\_poa, G\_control\_stop\_poa=G\_control\_stop\_poa, G\_poa=G\_drive\_poa if drive else G\_stop\_poa, Tcell\_drive\_C=float(tcell\_drive), Tcell\_stop\_C=float(tcell\_stop), Tcell\_control\_stop\_C=float(tcell\_control\_stop), Tcell\_C=float(tcell\_drive if drive else tcell\_stop), Tamb\_C=float(tamb), headwind\_ms=float(headwind), elevation\_m=get\_profile\_val(route\_profile, float(s\_query\_km or 0.0), 'elev\_m', 0.0)) を返す。
\item 条件 forecast\_time\_ns is not None を判定し、真なら内部処理を行う。
\item   idx に int(np.searchsorted(forecast\_time\_ns, timestamp\_ns(t\_utc)) - 1) の結果を代入する。
\item   idx に int(np.clip(idx, 0, len(df) - 1)) の結果を代入する。
\item   上の条件が偽の場合:
\item   elapsed に (t\_utc - start\_utc).total\_seconds() の結果を代入する。
\item   idx に int(np.clip(elapsed / max(args.dt, 0.001), 0, len(df) - 1)) の結果を代入する。
\item row に df.iloc[idx] の結果を代入する。
\item G\_raw に float(row.get('GHI', 0.0)) * args.solar\_gain の結果を代入する。
\item poa\_drive\_raw に float(row.get('POA\_drive', row.get('GHI', 0.0))) の結果を代入する。
\item poa\_stop\_ideal\_raw に float(row.get('POA\_stop\_ideal', poa\_drive\_raw)) の結果を代入する。
\item poa\_stop\_raw に poa\_drive\_raw + args.stop\_tilt\_fraction * (poa\_stop\_ideal\_raw - poa\_drive\_raw) の結果を代入する。
\item poa\_control\_stop\_raw に poa\_drive\_raw + args.control\_stop\_tilt\_fraction * (poa\_stop\_ideal\_raw - poa\_drive\_raw) の結果を代入する。
\item G\_drive\_poa に poa\_drive\_raw * args.solar\_gain * args.poa\_gain\_drive の結果を代入する。
\item G\_stop\_poa に poa\_stop\_raw * args.solar\_gain * args.poa\_gain\_stop の結果を代入する。
\item G\_control\_stop\_poa に poa\_control\_stop\_raw * args.solar\_gain * args.poa\_gain\_stop の結果を代入する。
\item tcell\_drive に float(row.get('Tcell\_drive\_C', row.get('Tcell\_C', 40.0))) の結果を代入する。
\item tcell\_stop\_ideal に float(row.get('Tcell\_stop\_ideal\_C', tcell\_drive)) の結果を代入する。
\item tcell\_stop に tcell\_drive + args.stop\_tilt\_fraction * (tcell\_stop\_ideal - tcell\_drive) の結果を代入する。
\item tcell\_control\_stop に tcell\_drive + args.control\_stop\_tilt\_fraction * (tcell\_stop\_ideal - tcell\_drive) の結果を代入する。
\item dict(G\_raw=G\_raw, G\_drive\_poa=G\_drive\_poa, G\_stop\_poa=G\_stop\_poa, G\_control\_stop\_poa=G\_control\_stop\_poa, G\_poa=G\_drive\_poa if drive else G\_stop\_poa, Tcell\_drive\_C=tcell\_drive, Tcell\_stop\_C=tcell\_stop, Tcell\_control\_stop\_C=tcell\_control\_stop, Tcell\_C=tcell\_drive if drive else tcell\_stop, Tamb\_C=float(row.get('Tamb\_C', 30.0)), headwind\_ms=float(row.get('headwind\_ms', 0.0)), elevation\_m=get\_profile\_val(route\_profile, float(s\_query\_km or 0.0), 'elev\_m', 0.0)) を返す。
            \end{enumerate}
            \paragraph{代表コード断片}
            \begin{lstlisting}[language=Python]
    def forecast_at_time(t_utc, s_query_km=None, drive=True):
        if len(df) == 0:
            return dict(
                G_raw=0.0,
                G_drive_poa=0.0,
                G_stop_poa=0.0,
                G_control_stop_poa=0.0,
                G_poa=0.0,
                Tcell_drive_C=40.0,
                Tcell_stop_C=40.0,
                Tcell_control_stop_C=40.0,
                Tcell_C=40.0,
                Tamb_C=30.0,
                headwind_ms=0.0,
                elevation_m=0.0,
            )
        if forecast_grid is not None:
            ghi = interp_forecast_grid(forecast_grid, 'GHI', t_utc, s_query_km, 0.0)
            poa_drive_raw = interp_forecast_grid(forecast_grid, 'POA_drive', t_utc, s_query_km, ghi)
            poa_stop_ideal_raw = interp_forecast_grid(
                forecast_grid, 'POA_stop_ideal', t_utc, s_query_km, poa_drive_raw
            )
            poa_stop_raw = poa_drive_raw + args.stop_tilt_fraction * (poa_stop_ideal_raw - poa_drive_raw)
            poa_control_stop_raw = poa_drive_raw + args.control_stop_tilt_fraction * (
                poa_stop_ideal_raw - poa_drive_raw
            )
            tamb = interp_forecast_grid(forecast_grid, 'Tamb_C', t_utc, s_query_km, 30.0)
            tcell_drive = interp_forecast_grid(
                forecast_grid,
                'Tcell_drive_C',
                t_utc,
                s_query_km,
                interp_forecast_grid(forecast_grid, 'Tcell_C', t_utc, s_query_km, tamb),
            )
            tcell_stop_ideal = interp_forecast_grid(
...
\end{lstlisting}

\subsection{L2236 関数 main.apply\_stop\_poa}
            \begin{itemize}[leftmargin=1.6em]
              \item 行範囲: L2236--L2242
              \item 所属: \path|main|
              \item 定義: \path|apply_stop_poa(env, *, control_stop = False)|
              \item docstring: 明示されていない。
              \item このブロックが直接呼ぶ主な関数・メソッド: float, get
              \item 戻り値の要点: env
              \item 主なローカル代入名: irradiance\_key, temperature\_key
              \item 読み取る主なインスタンス属性: 特になし。
              \item 更新する主なインスタンス属性: 特になし。
              \item 明示的に送出する例外: 明示的raiseはない。
              \item 制御構造の規模: 条件分岐 0、ループ 0、try 0
            \end{itemize}
            \paragraph{この定義を読むためのPython構文}
            \begin{itemize}[leftmargin=1.6em]
            \item 特別な構文上の補足は少ない。
            \end{itemize}
            \paragraph{上から順の処理}
            \begin{enumerate}[leftmargin=1.8em]
            \item irradiance\_key に 'G\_control\_stop\_poa' if control\_stop else 'G\_stop\_poa' の結果を代入する。
\item temperature\_key に 'Tcell\_control\_stop\_C' if control\_stop else 'Tcell\_stop\_C' の結果を代入する。
\item env['G\_poa'] に float(env.get(irradiance\_key, env.get('G\_raw', env.get('G\_poa', 0.0)))) の結果を代入する。
\item env['Tcell\_C'] に float(env.get(temperature\_key, env.get('Tcell\_C', env.get('Tamb\_C', 30.0)))) の結果を代入する。
\item env['panel\_stop\_mode'] に 'control\_horizontal' if control\_stop else 'camp\_or\_strategy\_tilt' の結果を代入する。
\item env を返す。
            \end{enumerate}
            \paragraph{代表コード断片}
            \begin{lstlisting}[language=Python]
    def apply_stop_poa(env, *, control_stop=False):
        irradiance_key = 'G_control_stop_poa' if control_stop else 'G_stop_poa'
        temperature_key = 'Tcell_control_stop_C' if control_stop else 'Tcell_stop_C'
        env['G_poa'] = float(env.get(irradiance_key, env.get('G_raw', env.get('G_poa', 0.0))))
        env['Tcell_C'] = float(env.get(temperature_key, env.get('Tcell_C', env.get('Tamb_C', 30.0))))
        env['panel_stop_mode'] = 'control_horizontal' if control_stop else 'camp_or_strategy_tilt'
        return env
\end{lstlisting}

\subsection{L2244 関数 main.reference\_speed\_command}
            \begin{itemize}[leftmargin=1.6em]
              \item 行範囲: L2244--L2262
              \item 所属: \path|main|
              \item 定義: \path!reference_speed_command(t_utc: datetime, s_now_km: float) -> float | None!
              \item docstring: 明示されていない。
              \item このブロックが直接呼ぶ主な関数・メソッド: clip, float, max, reference\_distance\_at\_time, reference\_value\_at\_time, timedelta
              \item 戻り値の要点: max(0.0, float(base\_speed\_kmh) * ref\_speed\_tracking\_gain + correction\_kmh) / None / None
              \item 主なローカル代入名: base\_speed\_kmh, correction\_kmh, lag\_km, lead\_km, ref\_s\_future, ref\_s\_now, t\_future
              \item 読み取る主なインスタンス属性: 特になし。
              \item 更新する主なインスタンス属性: 特になし。
              \item 明示的に送出する例外: 明示的raiseはない。
              \item 制御構造の規模: 条件分岐 4、ループ 0、try 0
            \end{itemize}
            \paragraph{この定義を読むためのPython構文}
            \begin{itemize}[leftmargin=1.6em]
            \item `->`以後は戻り値の型注釈であり、実行時の値を自動変換する命令ではない。
            \end{itemize}
            \paragraph{上から順の処理}
            \begin{enumerate}[leftmargin=1.8em]
            \item 条件 not ref\_speed\_tracking\_enabled を判定し、真なら内部処理を行う。
\item   None を返す。
\item ref\_s\_now に reference\_distance\_at\_time(t\_utc) の結果を代入する。
\item 条件 ref\_s\_now is None を判定し、真なら内部処理を行う。
\item   None を返す。
\item base\_speed\_kmh に reference\_value\_at\_time(t\_utc, 'speed\_kmh') の結果を代入する。
\item 条件 base\_speed\_kmh is None を判定し、真なら内部処理を行う。
\item   t\_future に t\_utc + timedelta(seconds=max(ref\_speed\_tracking\_lookahead\_sec, 1.0)) の結果を代入する。
\item   ref\_s\_future に reference\_distance\_at\_time(t\_future) の結果を代入する。
\item   条件 ref\_s\_future is not None を判定し、真なら内部処理を行う。
\item     base\_speed\_kmh に max(0.0, (ref\_s\_future - ref\_s\_now) * 3600.0 / max(ref\_speed\_tracking\_lookahead\_sec, 1.0)) の結果を代入する。
\item     上の条件が偽の場合:
\item     base\_speed\_kmh に 0.0 の結果を代入する。
\item lag\_km に max(0.0, float(ref\_s\_now) - float(s\_now\_km)) の結果を代入する。
\item lead\_km に max(0.0, float(s\_now\_km) - float(ref\_s\_now)) の結果を代入する。
\item correction\_kmh に ref\_speed\_tracking\_lag\_gain * lag\_km - ref\_speed\_tracking\_lead\_gain * lead\_km の結果を代入する。
\item correction\_kmh に float(np.clip(correction\_kmh, -ref\_speed\_tracking\_max\_correction, ref\_speed\_tracking\_max\_correction)) の結果を代入する。
\item max(0.0, float(base\_speed\_kmh) * ref\_speed\_tracking\_gain + correction\_kmh) を返す。
            \end{enumerate}
            \paragraph{代表コード断片}
            \begin{lstlisting}[language=Python]
    def reference_speed_command(t_utc: datetime, s_now_km: float) -> float | None:
        if not ref_speed_tracking_enabled:
            return None
        ref_s_now = reference_distance_at_time(t_utc)
        if ref_s_now is None:
            return None
        base_speed_kmh = reference_value_at_time(t_utc, 'speed_kmh')
        if base_speed_kmh is None:
            t_future = t_utc + timedelta(seconds=max(ref_speed_tracking_lookahead_sec, 1.0))
            ref_s_future = reference_distance_at_time(t_future)
            if ref_s_future is not None:
                base_speed_kmh = max(0.0, (ref_s_future - ref_s_now) * 3600.0 / max(ref_speed_tracking_lookahead_sec, 1.0))
            else:
                base_speed_kmh = 0.0
        lag_km = max(0.0, float(ref_s_now) - float(s_now_km))
        lead_km = max(0.0, float(s_now_km) - float(ref_s_now))
        correction_kmh = ref_speed_tracking_lag_gain * lag_km - ref_speed_tracking_lead_gain * lead_km
        correction_kmh = float(np.clip(correction_kmh, -ref_speed_tracking_max_correction, ref_speed_tracking_max_correction))
        return max(0.0, float(base_speed_kmh) * ref_speed_tracking_gain + correction_kmh)
\end{lstlisting}

\subsection{L2264 関数 main.resolved\_headwind\_ms}
            \begin{itemize}[leftmargin=1.6em]
              \item 行範囲: L2264--L2272
              \item 所属: \path|main|
              \item 定義: \path|resolved_headwind_ms(env: dict, s_query_km: float) -> float|
              \item docstring: 明示されていない。
              \item このブロックが直接呼ぶ主な関数・メソッド: float, get, get\_profile\_val, isfinite
              \item 戻り値の要点: get\_profile\_val(route\_profile, s\_query\_km, 'headwind\_ms', 0.0) / weather\_value
              \item 主なローカル代入名: weather\_value
              \item 読み取る主なインスタンス属性: 特になし。
              \item 更新する主なインスタンス属性: 特になし。
              \item 明示的に送出する例外: 明示的raiseはない。
              \item 制御構造の規模: 条件分岐 1、ループ 0、try 1
            \end{itemize}
            \paragraph{この定義を読むためのPython構文}
            \begin{itemize}[leftmargin=1.6em]
            \item `->`以後は戻り値の型注釈であり、実行時の値を自動変換する命令ではない。
\item try文は例外が起き得る処理と、異常時または終了時の経路を分ける。
            \end{itemize}
            \paragraph{上から順の処理}
            \begin{enumerate}[leftmargin=1.8em]
            \item weather\_value に env.get('headwind\_ms') の結果を代入する。
\item 例外処理を伴う try ブロックを実行する。
\item   weather\_value に float(weather\_value) の結果を代入する。
\item   (TypeError, ValueError)を捕捉した場合:
\item   weather\_value に float('nan') の結果を代入する。
\item 条件 np.isfinite(weather\_value) を判定し、真なら内部処理を行う。
\item   weather\_value を返す。
\item get\_profile\_val(route\_profile, s\_query\_km, 'headwind\_ms', 0.0) を返す。
            \end{enumerate}
            \paragraph{代表コード断片}
            \begin{lstlisting}[language=Python]
    def resolved_headwind_ms(env: dict, s_query_km: float) -> float:
        weather_value = env.get('headwind_ms')
        try:
            weather_value = float(weather_value)
        except (TypeError, ValueError):
            weather_value = float('nan')
        if np.isfinite(weather_value):
            return weather_value
        return get_profile_val(route_profile, s_query_km, 'headwind_ms', 0.0)
\end{lstlisting}

\subsection{L2274 関数 main.integrate\_stationary\_duration}
            \begin{itemize}[leftmargin=1.6em]
              \item 行範囲: L2274--L2323
              \item 所属: \path|main|
              \item 定義: \path|integrate_stationary_duration(t_utc, z, Tb, s_km, dt_wait, *, control_stop = False)|
              \item docstring: 明示されていない。
              \item このブロックが直接呼ぶ主な関数・メソッド: apply\_stop\_poa, choose\_stationary\_integration\_step\_seconds, clip, electrical\_balance, float, forecast\_at\_time, forecast\_has\_coverage, get, get\_profile\_val, max, min, resolved\_headwind\_ms
              \item 戻り値の要点: (t\_cursor, float(z\_cur), float(Tb\_cur), float(dt\_wait - remaining)) / (t\_utc, z, Tb, 0.0)
              \item 主なローカル代入名: P\_pack, Tb\_cur, dt\_local, dt\_sample, dt\_wait, env, headwind\_ms, loss\_int, out, remaining, slope\_pct, t\_cursor, z\_cur
              \item 読み取る主なインスタンス属性: 特になし。
              \item 更新する主なインスタンス属性: 特になし。
              \item 明示的に送出する例外: 明示的raiseはない。
              \item 制御構造の規模: 条件分岐 2、ループ 1、try 0
            \end{itemize}
            \paragraph{この定義を読むためのPython構文}
            \begin{itemize}[leftmargin=1.6em]
            \item 特別な構文上の補足は少ない。
            \end{itemize}
            \paragraph{上から順の処理}
            \begin{enumerate}[leftmargin=1.8em]
            \item dt\_wait に max(0.0, float(dt\_wait)) の結果を代入する。
\item 条件 dt\_wait <= 0.0 を判定し、真なら内部処理を行う。
\item   (t\_utc, z, Tb, 0.0) を返す。
\item slope\_pct に get\_profile\_val(route\_profile, s\_km, 'slope\_pct', 0.0) の結果を代入する。
\item dt\_sample に choose\_stationary\_integration\_step\_seconds(args.dt, forecast\_native\_dt\_sec, stationary\_prediction\_step\_sec) の結果を代入する。
\item remaining に dt\_wait の結果を代入する。
\item t\_cursor に t\_utc の結果を代入する。
\item z\_cur に float(z) の結果を代入する。
\item Tb\_cur に float(Tb) の結果を代入する。
\item 条件 remaining > 1e-09 and forecast\_has\_coverage(t\_cursor) が成り立つ間くり返す。
\item   dt\_local に min(dt\_sample, remaining) の結果を代入する。
\item   env に forecast\_at\_time(t\_cursor, s\_km, drive=False) の結果を代入する。
\item   条件 control\_stop を判定し、真なら内部処理を行う。
\item     apply\_stop\_poa(...) を実行する。
\item   headwind\_ms に resolved\_headwind\_ms(env, s\_km) の結果を代入する。
\item   out に model.electrical\_balance(0.0, slope\_pct, z\_cur, Tb\_cur, env['G\_poa'], env['Tcell\_C'], headwind\_ms=headwind\_ms, aux\_power\_w=model.scheduled\_auxiliary\_power(is\_driving=False, irradiance\_wm2=env['G\_poa']), ambient\_temp\_c=env.get('Tamb\_C'), elevation\_m=env.get('elevation\_m', 0.0)) の結果を代入する。
\item   P\_pack に float(out['P\_pack']) の結果を代入する。
\item   loss\_int に float(out['losses\_int']) の結果を代入する。
\item   z\_cur に model.soc\_step(z\_cur, P\_pack, dt\_local, current\_a=float(out['I']), Tbat\_C=Tb\_cur) の結果を代入する。
\item   Tb\_cur に Tb\_cur + dt\_local / 1800.0 * (env['Tamb\_C'] - Tb\_cur) + loss\_int * dt\_local / 50000.0 の結果を代入する。
\item   z\_cur に float(np.clip(z\_cur, p.soc\_min, p.soc\_max)) の結果を代入する。
\item   Tb\_cur に float(np.clip(Tb\_cur, p.T\_min, p.T\_max)) の結果を代入する。
\item   t\_cursor を Add で更新する。
\item   remaining を Sub で更新する。
\item (t\_cursor, float(z\_cur), float(Tb\_cur), float(dt\_wait - remaining)) を返す。
            \end{enumerate}
            \paragraph{代表コード断片}
            \begin{lstlisting}[language=Python]
    def integrate_stationary_duration(t_utc, z, Tb, s_km, dt_wait, *, control_stop=False):
        dt_wait = max(0.0, float(dt_wait))
        if dt_wait <= 0.0:
            return t_utc, z, Tb, 0.0
        slope_pct = get_profile_val(route_profile, s_km, 'slope_pct', 0.0)
        dt_sample = choose_stationary_integration_step_seconds(
            args.dt,
            forecast_native_dt_sec,
            stationary_prediction_step_sec,
        )
        remaining = dt_wait
        t_cursor = t_utc
        z_cur = float(z)
        Tb_cur = float(Tb)
        while remaining > 1.0e-9 and forecast_has_coverage(t_cursor):
            dt_local = min(dt_sample, remaining)
            env = forecast_at_time(t_cursor, s_km, drive=False)
            if control_stop:
                apply_stop_poa(env, control_stop=True)
            headwind_ms = resolved_headwind_ms(env, s_km)
            out = model.electrical_balance(
                0.0,
                slope_pct,
                z_cur,
                Tb_cur,
                env['G_poa'],
                env['Tcell_C'],
                headwind_ms=headwind_ms,
                aux_power_w=model.scheduled_auxiliary_power(
                    is_driving=False,
                    irradiance_wm2=env['G_poa'],
                ),
                ambient_temp_c=env.get('Tamb_C'),
                elevation_m=env.get('elevation_m', 0.0),
            )
...
\end{lstlisting}

\subsection{L2325 関数 main.step\_wait}
            \begin{itemize}[leftmargin=1.6em]
              \item 行範囲: L2325--L2330
              \item 所属: \path|main|
              \item 定義: \path|step_wait(t_utc, z, Tb, s_km)|
              \item docstring: 明示されていない。
              \item このブロックが直接呼ぶ主な関数・メソッド: integrate\_stationary\_duration, is\_drive\_time, max, next\_drive\_start, total\_seconds
              \item 戻り値の要点: integrate\_stationary\_duration(t\_utc, z, Tb, s\_km, dt\_wait) / (t\_utc, z, Tb, 0.0)
              \item 主なローカル代入名: dt\_wait, t\_start
              \item 読み取る主なインスタンス属性: 特になし。
              \item 更新する主なインスタンス属性: 特になし。
              \item 明示的に送出する例外: 明示的raiseはない。
              \item 制御構造の規模: 条件分岐 1、ループ 0、try 0
            \end{itemize}
            \paragraph{この定義を読むためのPython構文}
            \begin{itemize}[leftmargin=1.6em]
            \item 特別な構文上の補足は少ない。
            \end{itemize}
            \paragraph{上から順の処理}
            \begin{enumerate}[leftmargin=1.8em]
            \item 条件 schedule is None or schedule.is\_drive\_time(t\_utc) を判定し、真なら内部処理を行う。
\item   (t\_utc, z, Tb, 0.0) を返す。
\item t\_start に schedule.next\_drive\_start(t\_utc) の結果を代入する。
\item dt\_wait に max(0.0, (t\_start - t\_utc).total\_seconds()) の結果を代入する。
\item integrate\_stationary\_duration(t\_utc, z, Tb, s\_km, dt\_wait) を返す。
            \end{enumerate}
            \paragraph{代表コード断片}
            \begin{lstlisting}[language=Python]
    def step_wait(t_utc, z, Tb, s_km):
        if schedule is None or schedule.is_drive_time(t_utc):
            return t_utc, z, Tb, 0.0
        t_start = schedule.next_drive_start(t_utc)
        dt_wait = max(0.0, (t_start - t_utc).total_seconds())
        return integrate_stationary_duration(t_utc, z, Tb, s_km, dt_wait)
\end{lstlisting}

\subsection{L2332 関数 main.mpc\_solve\_distance}
            \begin{itemize}[leftmargin=1.6em]
              \item 行範囲: L2332--L3020
              \item 所属: \path|main|
              \item 定義: \path|mpc_solve_distance(t0_utc, s0_km, z0, Tb0, v0_kmh, v_init = None, v_init_s_km = None)|
              \item docstring: 明示されていない。
              \item このブロックが直接呼ぶ主な関数・メソッド: DataFrame, Path, ValueError, abs, absolute\_control\_distances, append, array, asarray, average\_profile, bool, build\_balance\_seed, build\_upper\_distance\_horizon
              \item 戻り値の要点: (v0, segments, u\_seq, ctrl\_s\_abs) / (v0\_kmh, [\{'v\_kmh': v0\_kmh, 'dt\_sec': args.dt\}], np.array([v0\_kmh], dtype=float), np.array([float(s0\_km)], dtype=float)) / (1.0 - alpha) * u\_vec[idx] + alpha * u\_vec[idx\_next] / np.array(u\_seed, dtype=float)
              \item 主なローカル代入名: I, J, Nc, Np, P\_mech\_wheel, P\_pack, P\_pv, Tb, Tb\_next, Tb\_seed, V, \_, alpha, best\_score, best\_v, bounds, candidate\_grid, checkpoint\_base, checkpoint\_path, checkpoint\_payload
              \item 読み取る主なインスタンス属性: 特になし。
              \item 更新する主なインスタンス属性: 特になし。
              \item 明示的に送出する例外: ValueError('mpc.upper\_lock\_initial\_policy requires paths.initial\_upper\_policy\_csv')
              \item 制御構造の規模: 条件分岐 38、ループ 7、try 0
            \end{itemize}
            \paragraph{この定義を読むためのPython構文}
            \begin{itemize}[leftmargin=1.6em]
            \item 内包表記は、反復・条件・値生成を一つの式で記述してcollectionを作る。
\item with文はcontext managerに開始・終了処理を任せ、ファイルなどの資源を確実に解放する。
            \end{itemize}
            \paragraph{上から順の処理}
            \begin{enumerate}[leftmargin=1.8em]
            \item horizon に build\_upper\_distance\_horizon(mode=upper\_horizon\_mode, s0\_km=s0\_km, race\_km=race\_km, ds\_km=upper\_ds\_km, horizon\_km=upper\_horizon\_km, max\_steps=upper\_max\_steps, ctrl\_km=upper\_ctrl\_km, adaptive\_min\_ds\_km=upper\_adaptive\_min\_ds\_km, adaptive\_max\_ds\_km=upper\_adaptive\_max\_ds\_km, adaptive\_growth=upper\_adaptive\_growth) の結果を代入する。
\item ds\_seq に np.array(horizon.ds\_seq\_km, dtype=float) の結果を代入する。
\item seg\_s に np.array(horizon.seg\_s\_km, dtype=float) の結果を代入する。
\item Np に int(len(ds\_seq)) の結果を代入する。
\item 条件 Np <= 0 を判定し、真なら内部処理を行う。
\item   (v0\_kmh, [\{'v\_kmh': v0\_kmh, 'dt\_sec': args.dt\}], np.array([v0\_kmh], dtype=float), np.array([float(s0\_km)], dtype=float)) を返す。
\item v\_min\_solver に max(0.1, float(upper\_vmin\_kmh)) の結果を代入する。
\item ctrl\_s に np.array(horizon.ctrl\_s\_km, dtype=float) の結果を代入する。
\item ctrl\_s\_abs に absolute\_control\_distances(s0\_km, ctrl\_s) の結果を代入する。
\item Nc に int(len(ctrl\_s)) の結果を代入する。
\item 条件 v\_init is not None and v\_init\_s\_km is not None and (len(v\_init) == len(v\_init\_s\_km)) を判定し、真なら内部処理を行う。
\item   x0 に shift\_upper\_policy\_warm\_start(v\_init\_s\_km, v\_init, ctrl\_s\_abs, minimum\_speed\_kmh=v\_min\_solver, maximum\_speed\_kmh=v\_max\_kmh) の結果を代入する。
\item   上の条件が偽の場合:
\item   条件 v\_init is not None and len(v\_init) == Nc を判定し、真なら内部処理を行う。
\item     x0 に np.array(v\_init, dtype=float) の結果を代入する。
\item     上の条件が偽の場合:
\item     条件 initial\_upper\_policy\_df is not None and (not initial\_upper\_policy\_df.empty) を判定し、真なら内部処理を行う。
\item       x0 に np.interp(ctrl\_s\_abs, initial\_upper\_policy\_df['s\_km'].to\_numpy(dtype=float), initial\_upper\_policy\_df['v\_kmh'].to\_numpy(dtype=float)) の結果を代入する。
\item       x0 に np.clip(x0, v\_min\_solver, v\_max\_kmh) の結果を代入する。
\item       上の条件が偽の場合:
\item       x0 に np.array([v\_max\_kmh] * Nc, dtype=float) の結果を代入する。
\item bounds に [(v\_min\_solver, v\_max\_kmh)] * Nc の結果を代入する。
\item idx に np.searchsorted(ctrl\_s, seg\_s, side='right') - 1 の結果を代入する。
\item idx に np.clip(idx, 0, Nc - 1) の結果を代入する。
\item idx\_next に np.clip(idx + 1, 0, Nc - 1) の結果を代入する。
\item denom に np.maximum(ctrl\_s[idx\_next] - ctrl\_s[idx], 1e-06) の結果を代入する。
\item alpha に (seg\_s - ctrl\_s[idx]) / denom の結果を代入する。
\item 関数 expand\_ctrl を定義する。
\item 関数 build\_balance\_seed を定義する。
\item selected\_prediction に \{\} の結果を代入する。
\item 関数 cost を定義する。
\item sim\_log(...) を実行する。
\item checkpoint\_base に Path(args.latest\_manifest\_json) if args.latest\_manifest\_json else Path(args.out\_csv) の結果を代入する。
\item progress\_path に checkpoint\_base.with\_name('upper\_solver\_progress.json') の結果を代入する。
\item 関数 save\_solver\_progress を定義する。
\item 条件 upper\_lock\_initial\_policy を判定し、真なら内部処理を行う。
\item   条件 initial\_upper\_policy\_df is None or initial\_upper\_policy\_df.empty を判定し、真なら内部処理を行う。
\item     ValueError('mpc.upper\_lock\_initial\_policy requires paths.initial\_upper\_policy\_csv') を送出する。
\item   u\_seq に np.asarray(x0, dtype=float) の結果を代入する。
\item   solve\_info に \{'success': True, 'method': 'locked\_external\_policy', 'label': 'mesh\_validation\_policy', 'message': 'optimization intentionally bypassed for fixed-policy mesh validation', 'fun': float('nan'), 'nit': 0, 'cem\_used': False, 'shgo\_used': False, 'shgo\_nfev': 0, 'shgo\_local\_minima': 0, 'discrete\_global\_proof': False, 'discrete\_grid\_levels': 0, 'discrete\_grid\_values': [], 'discrete\_grid\_candidates': 0, 'discrete\_grid\_nonfinite': 0, 'discrete\_grid\_best\_fun': float('inf'), 'discrete\_grid\_best\_x': [], 'deterministic\_seed\_candidates': 1, 'deterministic\_seed\_nonfinite': 0, 'finite\_library\_candidates': 1, 'selected\_x': u\_seq.tolist(), 'selected\_no\_worse\_than\_grid': False, 'finite\_library\_global\_proof': False, 'continuous\_global\_proof': False, 'certificate\_scope': 'fixed-policy mesh validation only', 'candidates\_evaluated': 1\} の結果を代入する。
\item   上の条件が偽の場合:
\item   (u\_seq, solve\_info) に hybrid\_bounded\_minimize(cost, x0, bounds, maxiter=int(upper\_max\_iter), structured\_seeds=[('balance\_seed', build\_balance\_seed())], cem\_enabled=bool(upper\_global\_search\_enabled), cem\_mode=upper\_global\_search\_mode, cem\_generations=int(upper\_cem\_generations), cem\_population=int(upper\_cem\_population), cem\_elite=int(upper\_cem\_elite), local\_refine\_topk=int(upper\_local\_refine\_topk), rng\_seed=int(max(0.0, round(s0\_km * 10.0))), shgo\_samples=int(upper\_shgo\_samples), shgo\_iters=int(upper\_shgo\_iters), cert\_grid\_levels=int(upper\_cert\_grid\_levels), cert\_grid\_values=[float(value) for value in upper\_cert\_grid\_values], cert\_max\_evaluations=int(upper\_cert\_max\_evaluations), cert\_workers=int(upper\_cert\_workers), progress\_callback=save\_solver\_progress, cert\_progress\_interval=int(cfg.get('mpc', \{\}).get('upper\_cert\_progress\_interval', 25))) の結果を代入する。
\item selected\_cost\_recheck に cost(np.asarray(u\_seq, dtype=float), capture\_trace=True) の結果を代入する。
\item 条件 upper\_lock\_initial\_policy を判定し、真なら内部処理を行う。
\item   solve\_info['fun'] に float(selected\_cost\_recheck) の結果を代入する。
\item prediction\_trace\_rows に list(selected\_prediction.pop('trace', [])) の結果を代入する。
\item prediction\_trace\_path に checkpoint\_base.with\_name(f'upper\_selected\_prediction\_trace\_\{int(round(float(s0\_km) * 1000.0)):010d\}.csv') の結果を代入する。
\item 条件 prediction\_trace\_rows を判定し、真なら内部処理を行う。
\item   ensure\_parent\_dir(...) を実行する。
\item   pd.DataFrame(prediction\_trace\_rows).to\_csv(...) を実行する。
\item   solve\_info['selected\_prediction\_trace\_csv'] に str(prediction\_trace\_path) の結果を代入する。
\item   上の条件が偽の場合:
\item   solve\_info['selected\_prediction\_trace\_csv'] に '' の結果を代入する。
\item solve\_info['selected\_cost\_recheck'] に float(selected\_cost\_recheck) の結果を代入する。
\item solve\_info['selected\_prediction'] に dict(selected\_prediction) の結果を代入する。
\item solve\_info['selected\_mission\_feasible'] に bool(selected\_prediction.get('mission\_feasible', False)) の結果を代入する。
\item sim\_log(...) を実行する。
\item checkpoint\_path に checkpoint\_base.with\_name('upper\_solver\_checkpoint.json') の結果を代入する。
\item ensure\_parent\_dir(...) を実行する。
\item checkpoint\_payload に \{'generated\_utc': datetime.now(timezone.utc).isoformat(), 'profile\_yaml': args.profile\_yaml, 's\_km': float(s0\_km), 'prediction\_steps': int(Np), 'control\_dimensions': int(Nc), 'u\_seq\_kmh': np.asarray(u\_seq, dtype=float).tolist(), 'solve\_info': \{key: solve\_info.get(key) for key in ('success', 'method', 'label', 'message', 'fun', 'nit', 'cem\_used', 'shgo\_used', 'shgo\_nfev', 'shgo\_local\_minima', 'discrete\_global\_proof', 'discrete\_grid\_levels', 'discrete\_grid\_values', 'discrete\_grid\_candidates', 'discrete\_grid\_nonfinite', 'discrete\_grid\_best\_fun', 'discrete\_grid\_best\_x', 'deterministic\_seed\_candidates', 'finite\_library\_candidates', 'deterministic\_seed\_nonfinite', 'selected\_no\_worse\_than\_grid', 'selected\_x', 'finite\_library\_global\_proof', 'selected\_cost\_recheck', 'selected\_prediction', 'selected\_prediction\_trace\_csv', 'selected\_mission\_feasible', 'continuous\_global\_proof', 'certificate\_scope', 'candidates\_evaluated')\}\} の結果を代入する。
\item with 文で checkpoint\_path.open('w', encoding='utf-8') を管理しながら処理する。
\item   json.dump(...) を実行する。
\item sim\_log(...) を実行する。
\item upper\_solve\_log.append(...) を実行する。
\item v\_seq に expand\_ctrl(u\_seq) の結果を代入する。
\item segments に [] の結果を代入する。
\item t\_utc に t0\_utc の結果を代入する。
\item s\_km に float(s0\_km) の結果を代入する。
\item z に float(z0) の結果を代入する。
\item Tb に float(Tb0) の結果を代入する。
\item enumerate(v\_seq) を順に走査し、各要素を (idx\_seg, v\_k) に入れて処理する。
\item   (t\_utc, z, Tb, \_) に step\_wait(t\_utc, z, Tb, s\_km) の結果を代入する。
\item   ds\_step\_km に min(float(ds\_seq[idx\_seg]), max(0.0, race\_km - s\_km)) の結果を代入する。
\item   条件 ds\_step\_km <= 1e-09 を判定し、真なら内部処理を行う。
\item     Break 文を実行する。
\item   v\_k に float(np.clip(v\_k, 0.0, v\_max\_kmh)) の結果を代入する。
\item   vmax\_local に get\_profile\_val(speed\_profile, s\_km, 'v\_max\_kmh', v\_max\_kmh) の結果を代入する。
\item   条件 vmax\_local >= v\_min\_solver を判定し、真なら内部処理を行う。
\item     v\_k に max(v\_min\_solver, min(v\_k, vmax\_local)) の結果を代入する。
\item     上の条件が偽の場合:
            \end{enumerate}
            \paragraph{代表コード断片}
            \begin{lstlisting}[language=Python]
    def mpc_solve_distance(
        t0_utc,
        s0_km,
        z0,
        Tb0,
        v0_kmh,
        v_init=None,
        v_init_s_km=None,
    ):
        horizon = build_upper_distance_horizon(
            mode=upper_horizon_mode,
            s0_km=s0_km,
            race_km=race_km,
            ds_km=upper_ds_km,
            horizon_km=upper_horizon_km,
            max_steps=upper_max_steps,
            ctrl_km=upper_ctrl_km,
            adaptive_min_ds_km=upper_adaptive_min_ds_km,
            adaptive_max_ds_km=upper_adaptive_max_ds_km,
            adaptive_growth=upper_adaptive_growth,
        )
        ds_seq = np.array(horizon.ds_seq_km, dtype=float)
        seg_s = np.array(horizon.seg_s_km, dtype=float)
        Np = int(len(ds_seq))
        if Np <= 0:
            return (
                v0_kmh,
                [{'v_kmh': v0_kmh, 'dt_sec': args.dt}],
                np.array([v0_kmh], dtype=float),
                np.array([float(s0_km)], dtype=float),
            )

        v_min_solver = max(0.1, float(upper_vmin_kmh))
        ctrl_s = np.array(horizon.ctrl_s_km, dtype=float)
        ctrl_s_abs = absolute_control_distances(s0_km, ctrl_s)
...
\end{lstlisting}

\subsection{L2400 関数 main.mpc\_solve\_distance.expand\_ctrl}
            \begin{itemize}[leftmargin=1.6em]
              \item 行範囲: L2400--L2401
              \item 所属: \path|main.mpc_solve_distance|
              \item 定義: \path|expand_ctrl(u_vec)|
              \item docstring: 明示されていない。
              \item このブロックが直接呼ぶ主な関数・メソッド: 特に明示されていない。
              \item 戻り値の要点: (1.0 - alpha) * u\_vec[idx] + alpha * u\_vec[idx\_next]
              \item 主なローカル代入名: 特になし。
              \item 読み取る主なインスタンス属性: 特になし。
              \item 更新する主なインスタンス属性: 特になし。
              \item 明示的に送出する例外: 明示的raiseはない。
              \item 制御構造の規模: 条件分岐 0、ループ 0、try 0
            \end{itemize}
            \paragraph{この定義を読むためのPython構文}
            \begin{itemize}[leftmargin=1.6em]
            \item 特別な構文上の補足は少ない。
            \end{itemize}
            \paragraph{上から順の処理}
            \begin{enumerate}[leftmargin=1.8em]
            \item (1.0 - alpha) * u\_vec[idx] + alpha * u\_vec[idx\_next] を返す。
            \end{enumerate}
            \paragraph{代表コード断片}
            \begin{lstlisting}[language=Python]
        def expand_ctrl(u_vec):
            return (1.0 - alpha) * u_vec[idx] + alpha * u_vec[idx_next]
\end{lstlisting}

\subsection{L2403 関数 main.mpc\_solve\_distance.build\_balance\_seed}
            \begin{itemize}[leftmargin=1.6em]
              \item 行範囲: L2403--L2503
              \item 所属: \path|main.mpc_solve_distance|
              \item 定義: \path|build_balance_seed()|
              \item docstring: 明示されていない。
              \item このブロックが直接呼ぶ主な関数・メソッド: abs, array, average\_profile, clip, electrical\_balance, float, forecast\_at\_time, get, get\_profile\_val, int, integrate\_stationary\_duration, len
              \item 戻り値の要点: np.array(u\_seed, dtype=float)
              \item 主なローカル代入名: Tb\_seed, \_, best\_score, best\_v, candidate\_grid, crossed\_stops, ds\_ctrl, ds\_leg\_km, dt\_seed, dwell\_done, dwell\_sec, edge\_s\_km, env, env\_leg, headwind\_leg, headwind\_ms, idx\_ctrl, limits, out, p\_pack
              \item 読み取る主なインスタンス属性: 特になし。
              \item 更新する主なインスタンス属性: 特になし。
              \item 明示的に送出する例外: 明示的raiseはない。
              \item 制御構造の規模: 条件分岐 7、ループ 3、try 0
            \end{itemize}
            \paragraph{この定義を読むためのPython構文}
            \begin{itemize}[leftmargin=1.6em]
            \item 内包表記は、反復・条件・値生成を一つの式で記述してcollectionを作る。
            \end{itemize}
            \paragraph{上から順の処理}
            \begin{enumerate}[leftmargin=1.8em]
            \item u\_seed に np.array(x0, dtype=float) の結果を代入する。
\item z\_seed に float(z0) の結果を代入する。
\item Tb\_seed に float(Tb0) の結果を代入する。
\item s\_seed に float(s0\_km) の結果を代入する。
\item t\_seed に t0\_utc の結果を代入する。
\item v\_prev\_seed に float(v0\_kmh) の結果を代入する。
\item range(Nc) を順に走査し、各要素を idx\_ctrl に入れて処理する。
\item   (t\_seed, z\_seed, Tb\_seed, \_) に step\_wait(t\_seed, z\_seed, Tb\_seed, s\_seed) の結果を代入する。
\item   vmax\_local に get\_profile\_val(speed\_profile, s\_seed, 'v\_max\_kmh', v\_max\_kmh) の結果を代入する。
\item   vmin\_local に v\_min\_solver の結果を代入する。
\item   条件 schedule is not None を判定し、真なら内部処理を行う。
\item     limits に schedule.speed\_limits(t\_seed) の結果を代入する。
\item     条件 limits is not None を判定し、真なら内部処理を行う。
\item       vmin\_local に max(vmin\_local, float(limits[0])) の結果を代入する。
\item       vmax\_local に min(vmax\_local, float(limits[1])) の結果を代入する。
\item   条件 vmax\_local < vmin\_local を判定し、真なら内部処理を行う。
\item     u\_seed[idx\_ctrl] に max(0.0, vmax\_local) の結果を代入する。
\item     Continue 文を実行する。
\item   candidate\_grid に np.linspace(vmin\_local, vmax\_local, num=max(5, min(13, int((vmax\_local - vmin\_local) / 4.0) + 2))) の結果を代入する。
\item   env に forecast\_at\_time(t\_seed, s\_seed, drive=True) の結果を代入する。
\item   slope\_pct に get\_profile\_val(route\_profile, s\_seed, 'slope\_pct', 0.0) の結果を代入する。
\item   headwind\_ms に resolved\_headwind\_ms(env, s\_seed) の結果を代入する。
\item   best\_v に float(np.clip(v\_prev\_seed, vmin\_local, vmax\_local)) の結果を代入する。
\item   best\_score に float('inf') の結果を代入する。
\item   candidate\_grid を順に走査し、各要素を v\_test に入れて処理する。
\item     out に model.electrical\_balance(v\_test / 3.6, slope\_pct, z\_seed, Tb\_seed, env['G\_poa'], env['Tcell\_C'], headwind\_ms=headwind\_ms, ambient\_temp\_c=env.get('Tamb\_C'), elevation\_m=env.get('elevation\_m', 0.0)) の結果を代入する。
\item     p\_pack に float(out['P\_pack']) の結果を代入する。
\item     score に abs(p\_pack) + 0.35 * (float(v\_test) - float(v\_prev\_seed)) ** 2 の結果を代入する。
\item     条件 score < best\_score を判定し、真なら内部処理を行う。
\item       best\_score に score の結果を代入する。
\item       best\_v に float(v\_test) の結果を代入する。
\item   u\_seed[idx\_ctrl] に best\_v の結果を代入する。
\item   条件 idx\_ctrl + 1 < len(ctrl\_s) を判定し、真なら内部処理を行う。
\item     ds\_ctrl に max(0.0, float(ctrl\_s[idx\_ctrl + 1] - ctrl\_s[idx\_ctrl])) の結果を代入する。
\item     上の条件が偽の場合:
\item     ds\_ctrl に max(float(ds\_seq[-1]), float(upper\_ds\_km)) の結果を代入する。
\item   target\_s\_km に s\_seed + ds\_ctrl の結果を代入する。
\item   crossed\_stops に [stop for stop in stop\_queue if s\_seed + 1e-09 < float(stop.get('s\_km', 0.0)) <= target\_s\_km + 1e-09] の結果を代入する。
\item   stop\_by\_distance に \{float(stop.get('s\_km', 0.0)): stop for stop in crossed\_stops\} の結果を代入する。
\item   seed\_edges に sorted(set([*stop\_by\_distance, float(target\_s\_km)])) の結果を代入する。
\item   stopped\_in\_chunk に False の結果を代入する。
\item   seed\_edges を順に走査し、各要素を edge\_s\_km に入れて処理する。
\item     ds\_leg\_km に max(0.0, float(edge\_s\_km) - s\_seed) の結果を代入する。
\item     条件 ds\_leg\_km > 1e-09 を判定し、真なら内部処理を行う。
\item       dt\_seed に ds\_leg\_km / max(best\_v, 0.001) * 3600.0 の結果を代入する。
\item       env\_leg に forecast\_at\_time(t\_seed, s\_seed, drive=True) の結果を代入する。
\item       slope\_leg に average\_profile(route\_profile, s\_seed, edge\_s\_km, 'slope\_pct', 0.0) の結果を代入する。
\item       headwind\_leg に resolved\_headwind\_ms(env\_leg, s\_seed) の結果を代入する。
\item       out に model.electrical\_balance(best\_v / 3.6, slope\_leg, z\_seed, Tb\_seed, env\_leg['G\_poa'], env\_leg['Tcell\_C'], headwind\_ms=headwind\_leg, ambient\_temp\_c=env\_leg.get('Tamb\_C'), elevation\_m=env\_leg.get('elevation\_m', 0.0)) の結果を代入する。
\item       z\_seed に float(np.clip(model.soc\_step(z\_seed, float(out['P\_pack']), dt\_seed, current\_a=float(out['I']), Tbat\_C=Tb\_seed), p.soc\_min, p.soc\_max)) の結果を代入する。
\item       Tb\_seed に float(np.clip(Tb\_seed + dt\_seed / 1800.0 * (env\_leg['Tamb\_C'] - Tb\_seed) + float(out['losses\_int']) * dt\_seed / 50000.0, p.T\_min, p.T\_max)) の結果を代入する。
\item       t\_seed を Add で更新する。
\item       s\_seed に float(edge\_s\_km) の結果を代入する。
\item     stop に stop\_by\_distance.get(float(edge\_s\_km)) の結果を代入する。
\item     条件 stop is not None を判定し、真なら内部処理を行う。
\item       dwell\_sec に max(0.0, float(stop.get('dwell\_sec', stop.get('dwell\_s', 0.0)))) の結果を代入する。
\item       (t\_seed, z\_seed, Tb\_seed, dwell\_done) に integrate\_stationary\_duration(t\_seed, z\_seed, Tb\_seed, s\_seed, dwell\_sec, control\_stop=True) の結果を代入する。
\item       stopped\_in\_chunk に stopped\_in\_chunk or dwell\_done > 1e-09 の結果を代入する。
\item   v\_prev\_seed に 0.0 if stopped\_in\_chunk else best\_v の結果を代入する。
\item np.array(u\_seed, dtype=float) を返す。
            \end{enumerate}
            \paragraph{代表コード断片}
            \begin{lstlisting}[language=Python]
        def build_balance_seed():
            u_seed = np.array(x0, dtype=float)
            z_seed = float(z0)
            Tb_seed = float(Tb0)
            s_seed = float(s0_km)
            t_seed = t0_utc
            v_prev_seed = float(v0_kmh)
            for idx_ctrl in range(Nc):
                t_seed, z_seed, Tb_seed, _ = step_wait(t_seed, z_seed, Tb_seed, s_seed)
                vmax_local = get_profile_val(speed_profile, s_seed, 'v_max_kmh', v_max_kmh)
                vmin_local = v_min_solver
                if schedule is not None:
                    limits = schedule.speed_limits(t_seed)
                    if limits is not None:
                        vmin_local = max(vmin_local, float(limits[0]))
                        vmax_local = min(vmax_local, float(limits[1]))
                if vmax_local < vmin_local:
                    u_seed[idx_ctrl] = max(0.0, vmax_local)
                    continue
                candidate_grid = np.linspace(vmin_local, vmax_local, num=max(5, min(13, int((vmax_local - vmin_local) / 4.0) + 2)))
                env = forecast_at_time(t_seed, s_seed, drive=True)
                slope_pct = get_profile_val(route_profile, s_seed, 'slope_pct', 0.0)
                headwind_ms = resolved_headwind_ms(env, s_seed)
                best_v = float(np.clip(v_prev_seed, vmin_local, vmax_local))
                best_score = float("inf")
                for v_test in candidate_grid:
                    out = model.electrical_balance(
                        v_test / 3.6, slope_pct, z_seed, Tb_seed, env['G_poa'], env['Tcell_C'],
                        headwind_ms=headwind_ms, ambient_temp_c=env.get('Tamb_C'),
                        elevation_m=env.get('elevation_m', 0.0),
                    )
                    p_pack = float(out['P_pack'])
                    score = abs(p_pack) + 0.35 * ((float(v_test) - float(v_prev_seed)) ** 2)
                    if score < best_score:
                        best_score = score
...
\end{lstlisting}

\subsection{L2507 関数 main.mpc\_solve\_distance.cost}
            \begin{itemize}[leftmargin=1.6em]
              \item 行範囲: L2507--L2812
              \item 所属: \path|main.mpc_solve_distance|
              \item 定義: \path|cost(u_vec, capture_trace = False)|
              \item docstring: 明示されていない。
              \item このブロックが直接呼ぶ主な関数・メソッド: append, average\_profile, bool, clear, clip, current\_drive\_window, electrical\_balance, expand\_ctrl, finish\_cost, float, forecast\_at\_time, forecast\_has\_coverage
              \item 戻り値の要点: finish\_cost(J, 'horizon\_complete') / final\_value / finish\_cost(J + upper\_cost\_cfg.constraint\_penalty * (1.0 + missing\_km) ** 2, 'forecast\_exhausted') / finish\_cost(J + upper\_cost\_cfg.constraint\_penalty * (1.0 + missing\_km) ** 2, 'zero\_speed\_before\_horizon\_end')
              \item 主なローカル代入名: I, J, P\_mech\_wheel, P\_pack, P\_pv, Tb, Tb\_next, V, day\_end\_crossing, ds\_chunk\_km, dt\_travel, dt\_wait, dwell\_done, dwell\_sec, elapsed\_plan\_sec, env, final\_value, forces, headwind\_ms, horizon\_target\_s\_km
              \item 読み取る主なインスタンス属性: 特になし。
              \item 更新する主なインスタンス属性: 特になし。
              \item 明示的に送出する例外: 明示的raiseはない。
              \item 制御構造の規模: 条件分岐 21、ループ 3、try 0
            \end{itemize}
            \paragraph{この定義を読むためのPython構文}
            \begin{itemize}[leftmargin=1.6em]
            \item 特別な構文上の補足は少ない。
            \end{itemize}
            \paragraph{上から順の処理}
            \begin{enumerate}[leftmargin=1.8em]
            \item z に float(z0) の結果を代入する。
\item Tb に float(Tb0) の結果を代入する。
\item s\_km に float(s0\_km) の結果を代入する。
\item t\_utc に t0\_utc の結果を代入する。
\item v\_prev に float(v0\_kmh) の結果を代入する。
\item p\_pack\_prev に None の結果を代入する。
\item elapsed\_plan\_sec に 0.0 の結果を代入する。
\item J に 0.0 の結果を代入する。
\item soc\_violation\_sq に 0.0 の結果を代入する。
\item min\_predicted\_soc に float(z) の結果を代入する。
\item prediction\_trace に [] の結果を代入する。
\item v\_seq に expand\_ctrl(u\_vec) の結果を代入する。
\item horizon\_target\_s\_km に float(s0\_km + np.sum(ds\_seq)) の結果を代入する。
\item predicted\_stop\_idx に 0 の結果を代入する。
\item 条件 predicted\_stop\_idx < len(stop\_queue) and float(stop\_queue[predicted\_stop\_idx].get('s\_km', 0.0)) <= s\_km + 1e-09 が成り立つ間くり返す。
\item   predicted\_stop\_idx を Add で更新する。
\item 関数 finish\_cost を定義する。
\item range(Np) を順に走査し、各要素を k\_i に入れて処理する。
\item   remaining\_segment\_km に float(ds\_seq[k\_i]) の結果を代入する。
\item   条件 remaining\_segment\_km > 1e-09 が成り立つ間くり返す。
\item     wait\_start\_time に t\_utc の結果を代入する。
\item     wait\_start\_soc に float(z) の結果を代入する。
\item     wait\_start\_tb に float(Tb) の結果を代入する。
\item     dt\_wait に 0.0 の結果を代入する。
\item     条件 predicted\_stop\_idx < len(stop\_queue) を判定し、真なら内部処理を行う。
\item       stop\_item に stop\_queue[predicted\_stop\_idx] の結果を代入する。
\item       stop\_s\_km に float(stop\_item.get('s\_km', 0.0)) の結果を代入する。
\item       条件 s\_km >= stop\_s\_km - 1e-07 を判定し、真なら内部処理を行う。
\item     (t\_utc, z, Tb, schedule\_wait) に step\_wait(t\_utc, z, Tb, s\_km) の結果を代入する。
\item     dt\_wait を Add で更新する。
\item     elapsed\_plan\_sec を Add で更新する。
\item     条件 dt\_wait > 1e-09 を判定し、真なら内部処理を行う。
\item       v\_prev に 0.0 の結果を代入する。
\item     条件 capture\_trace and dt\_wait > 1e-09 を判定し、真なら内部処理を行う。
\item       prediction\_trace.append(...) を実行する。
\item     条件 not forecast\_has\_coverage(t\_utc) を判定し、真なら内部処理を行う。
\item       missing\_km に max(0.0, horizon\_target\_s\_km - s\_km) の結果を代入する。
\item       finish\_cost(J + upper\_cost\_cfg.constraint\_penalty * (1.0 + missing\_km) ** 2, 'forecast\_exhausted') を返す。
\item     v\_k に float(v\_seq[k\_i]) の結果を代入する。
\item     vmax\_local に get\_profile\_val(speed\_profile, s\_km, 'v\_max\_kmh', v\_max\_kmh) の結果を代入する。
\item     条件 vmax\_local >= v\_min\_solver を判定し、真なら内部処理を行う。
\item       v\_k に max(v\_min\_solver, min(v\_k, vmax\_local)) の結果を代入する。
\item       上の条件が偽の場合:
\item       v\_k に max(0.0, min(v\_k, vmax\_local)) の結果を代入する。
\item     limits に schedule.speed\_limits(t\_utc) if schedule is not None else None の結果を代入する。
\item     条件 limits is not None を判定し、真なら内部処理を行う。
\item       v\_k に max(float(limits[0]), min(v\_k, float(limits[1]))) の結果を代入する。
\item     条件 v\_k <= 1e-06 を判定し、真なら内部処理を行う。
\item       missing\_km に max(0.0, horizon\_target\_s\_km - s\_km) の結果を代入する。
\item       finish\_cost(J + upper\_cost\_cfg.constraint\_penalty * (1.0 + missing\_km) ** 2, 'zero\_speed\_before\_horizon\_end') を返す。
\item     max\_chunk\_sec に integration\_step\_seconds() の結果を代入する。
\item     条件 schedule is not None を判定し、真なら内部処理を行う。
\item       win に schedule.current\_drive\_window(t\_utc) の結果を代入する。
\item       条件 win is None を判定し、真なら内部処理を行う。
\item       max\_chunk\_sec に min(max\_chunk\_sec, max(0.0, (win[1] - t\_utc).total\_seconds())) の結果を代入する。
\item     条件 max\_chunk\_sec <= 1e-09 を判定し、真なら内部処理を行う。
\item       t\_utc を Add で更新する。
\item       Continue 文を実行する。
\item     ds\_chunk\_km に min(remaining\_segment\_km, v\_k * max\_chunk\_sec / 3600.0) の結果を代入する。
\item     条件 predicted\_stop\_idx < len(stop\_queue) を判定し、真なら内部処理を行う。
\item       stop\_s\_km に float(stop\_queue[predicted\_stop\_idx].get('s\_km', 0.0)) の結果を代入する。
\item       条件 s\_km < stop\_s\_km < s\_km + ds\_chunk\_km を判定し、真なら内部処理を行う。
\item     dt\_travel に ds\_chunk\_km / v\_k * 3600.0 の結果を代入する。
\item     env に forecast\_at\_time(t\_utc, s\_km, drive=True) の結果を代入する。
\item     slope\_pct に average\_profile(route\_profile, s\_km, s\_km + ds\_chunk\_km, 'slope\_pct', 0.0) の結果を代入する。
\item     headwind\_ms に resolved\_headwind\_ms(env, s\_km) の結果を代入する。
\item     kinetic\_delta\_wh に 0.5 * p.m * ((v\_k / 3.6) ** 2 - (v\_prev / 3.6) ** 2) / 3600.0 の結果を代入する。
\item     inertial\_power\_w に kinetic\_delta\_wh * 3600.0 / max(dt\_travel, 1e-09) の結果を代入する。
\item     out に model.electrical\_balance(v\_k / 3.6, slope\_pct, z, Tb, env['G\_poa'], env['Tcell\_C'], headwind\_ms=headwind\_ms, inertial\_power\_w=inertial\_power\_w, ambient\_temp\_c=env.get('Tamb\_C'), elevation\_m=env.get('elevation\_m', 0.0)) の結果を代入する。
\item     I に float(out['I']) の結果を代入する。
\item     V に float(out['V']) の結果を代入する。
\item     P\_pv に float(out.get('P\_pv', 0.0)) の結果を代入する。
\item     P\_pack に float(out['P\_pack']) の結果を代入する。
\item     loss\_int に float(out['losses\_int']) の結果を代入する。
\item     loss\_line に float(out.get('losses\_line', 0.0)) の結果を代入する。
\item     P\_mech\_wheel に float(out.get('P\_mech\_wheel', 0.0)) の結果を代入する。
\item     kinetic\_step\_wh に max(0.0, kinetic\_delta\_wh) の結果を代入する。
\item     z\_next\_raw に model.soc\_step(z, P\_pack, dt\_travel, current\_a=I, Tbat\_C=Tb) の結果を代入する。
\item     min\_predicted\_soc に min(min\_predicted\_soc, float(z\_next\_raw)) の結果を代入する。
\item     soc\_violation\_sq を Add で更新する。
            \end{enumerate}
            \paragraph{代表コード断片}
            \begin{lstlisting}[language=Python]
        def cost(u_vec, capture_trace=False):
            z = float(z0)
            Tb = float(Tb0)
            s_km = float(s0_km)
            t_utc = t0_utc
            v_prev = float(v0_kmh)
            p_pack_prev = None
            elapsed_plan_sec = 0.0
            J = 0.0
            soc_violation_sq = 0.0
            min_predicted_soc = float(z)
            prediction_trace = []
            v_seq = expand_ctrl(u_vec)
            horizon_target_s_km = float(s0_km + np.sum(ds_seq))
            predicted_stop_idx = 0
            while (
                predicted_stop_idx < len(stop_queue)
                and float(stop_queue[predicted_stop_idx].get('s_km', 0.0)) <= s_km + 1.0e-9
            ):
                predicted_stop_idx += 1

            def finish_cost(value, status):
                status_text = str(status)
                missing_distance_km = max(0.0, horizon_target_s_km - s_km)
                mission_violation_sq = float(soc_violation_sq)
                if status_text != 'horizon_complete':
                    missing_fraction = missing_distance_km / max(1.0, horizon_target_s_km - s0_km)
                    # A truncated prediction is always lexicographically worse than
                    # every full-horizon feasible candidate, irrespective of soft weights.
                    mission_violation_sq += 1.0 + missing_fraction * missing_fraction
                mission_feasible = bool(
                    status_text == 'horizon_complete' and mission_violation_sq <= 1.0e-12
                )
                final_value = float(value)
                if upper_cost_cfg.objective_mode in {'fastest', 'fastest_feasible', 'minimum_time'}:
...
\end{lstlisting}

\subsection{L2528 関数 main.mpc\_solve\_distance.cost.finish\_cost}
            \begin{itemize}[leftmargin=1.6em]
              \item 行範囲: L2528--L2567
              \item 所属: \path|main.mpc_solve_distance.cost|
              \item 定義: \path|finish_cost(value, status)|
              \item docstring: 明示されていない。
              \item このブロックが直接呼ぶ主な関数・メソッド: bool, clear, float, isoformat, max, str, update
              \item 戻り値の要点: final\_value
              \item 主なローカル代入名: final\_value, missing\_distance\_km, missing\_fraction, mission\_feasible, mission\_violation\_sq, status\_text
              \item 読み取る主なインスタンス属性: 特になし。
              \item 更新する主なインスタンス属性: 特になし。
              \item 明示的に送出する例外: 明示的raiseはない。
              \item 制御構造の規模: 条件分岐 3、ループ 0、try 0
            \end{itemize}
            \paragraph{この定義を読むためのPython構文}
            \begin{itemize}[leftmargin=1.6em]
            \item 特別な構文上の補足は少ない。
            \end{itemize}
            \paragraph{上から順の処理}
            \begin{enumerate}[leftmargin=1.8em]
            \item status\_text に str(status) の結果を代入する。
\item missing\_distance\_km に max(0.0, horizon\_target\_s\_km - s\_km) の結果を代入する。
\item mission\_violation\_sq に float(soc\_violation\_sq) の結果を代入する。
\item 条件 status\_text != 'horizon\_complete' を判定し、真なら内部処理を行う。
\item   missing\_fraction に missing\_distance\_km / max(1.0, horizon\_target\_s\_km - s0\_km) の結果を代入する。
\item   mission\_violation\_sq を Add で更新する。
\item mission\_feasible に bool(status\_text == 'horizon\_complete' and mission\_violation\_sq <= 1e-12) の結果を代入する。
\item final\_value に float(value) の結果を代入する。
\item 条件 upper\_cost\_cfg.objective\_mode in \{'fastest', 'fastest\_feasible', 'minimum\_time'\} を判定し、真なら内部処理を行う。
\item   final\_value に float(elapsed\_plan\_sec) if mission\_feasible else 1e+18 + 1000000000000000.0 * mission\_violation\_sq + float(elapsed\_plan\_sec) の結果を代入する。
\item 条件 capture\_trace を判定し、真なら内部処理を行う。
\item   selected\_prediction.clear(...) を実行する。
\item   selected\_prediction.update(...) を実行する。
\item   selected\_prediction['trace'] に prediction\_trace の結果を代入する。
\item final\_value を返す。
            \end{enumerate}
            \paragraph{代表コード断片}
            \begin{lstlisting}[language=Python]
            def finish_cost(value, status):
                status_text = str(status)
                missing_distance_km = max(0.0, horizon_target_s_km - s_km)
                mission_violation_sq = float(soc_violation_sq)
                if status_text != 'horizon_complete':
                    missing_fraction = missing_distance_km / max(1.0, horizon_target_s_km - s0_km)
                    # A truncated prediction is always lexicographically worse than
                    # every full-horizon feasible candidate, irrespective of soft weights.
                    mission_violation_sq += 1.0 + missing_fraction * missing_fraction
                mission_feasible = bool(
                    status_text == 'horizon_complete' and mission_violation_sq <= 1.0e-12
                )
                final_value = float(value)
                if upper_cost_cfg.objective_mode in {'fastest', 'fastest_feasible', 'minimum_time'}:
                    final_value = (
                        float(elapsed_plan_sec)
                        if mission_feasible
                        else 1.0e18 + 1.0e15 * mission_violation_sq + float(elapsed_plan_sec)
                    )
                if capture_trace:
                    selected_prediction.clear()
                    selected_prediction.update(
                        {
                        'status': status_text,
                        'objective': final_value,
                        'terminal_time_utc': t_utc.isoformat(),
                        'terminal_distance_km': float(s_km),
                        'target_distance_km': float(horizon_target_s_km),
                        'missing_distance_km': float(missing_distance_km),
                        'terminal_soc': float(z),
                        'minimum_soc': float(min_predicted_soc),
                        'soc_violation_sq': float(soc_violation_sq),
                        'mission_violation_sq': float(mission_violation_sq),
                        'mission_feasible': mission_feasible,
                        'terminal_tb_c': float(Tb),
...
\end{lstlisting}

\subsection{L2818 関数 main.mpc\_solve\_distance.save\_solver\_progress}
            \begin{itemize}[leftmargin=1.6em]
              \item 行範囲: L2818--L2836
              \item 所属: \path|main.mpc_solve_distance|
              \item 定義: \path|save_solver_progress(progress: dict) -> None|
              \item docstring: 明示されていない。
              \item このブロックが直接呼ぶ主な関数・メソッド: dict, dump, ensure\_parent\_dir, float, get, int, isoformat, now, open, replace, sim\_log, str
              \item 戻り値の要点: 戻り値が無いか、return 文が明示されていない。
              \item 主なローカル代入名: handle, payload, tmp\_path
              \item 読み取る主なインスタンス属性: 特になし。
              \item 更新する主なインスタンス属性: 特になし。
              \item 明示的に送出する例外: 明示的raiseはない。
              \item 制御構造の規模: 条件分岐 0、ループ 0、try 0
            \end{itemize}
            \paragraph{この定義を読むためのPython構文}
            \begin{itemize}[leftmargin=1.6em]
            \item `->`以後は戻り値の型注釈であり、実行時の値を自動変換する命令ではない。
\item with文はcontext managerに開始・終了処理を任せ、ファイルなどの資源を確実に解放する。
            \end{itemize}
            \paragraph{上から順の処理}
            \begin{enumerate}[leftmargin=1.8em]
            \item payload に \{'generated\_utc': datetime.now(timezone.utc).isoformat(), 'profile\_yaml': args.profile\_yaml, 's\_km': float(s0\_km), 'prediction\_steps': int(Np), 'control\_dimensions': int(Nc), **dict(progress)\} の結果を代入する。
\item ensure\_parent\_dir(...) を実行する。
\item tmp\_path に progress\_path.with\_suffix(progress\_path.suffix + '.tmp') の結果を代入する。
\item with 文で tmp\_path.open('w', encoding='utf-8') を管理しながら処理する。
\item   json.dump(...) を実行する。
\item tmp\_path.replace(...) を実行する。
\item sim\_log(...) を実行する。
            \end{enumerate}
            \paragraph{代表コード断片}
            \begin{lstlisting}[language=Python]
        def save_solver_progress(progress: dict) -> None:
            payload = {
                'generated_utc': datetime.now(timezone.utc).isoformat(),
                'profile_yaml': args.profile_yaml,
                's_km': float(s0_km),
                'prediction_steps': int(Np),
                'control_dimensions': int(Nc),
                **dict(progress),
            }
            ensure_parent_dir(str(progress_path))
            tmp_path = progress_path.with_suffix(progress_path.suffix + '.tmp')
            with tmp_path.open('w', encoding='utf-8') as handle:
                json.dump(payload, handle, ensure_ascii=False, indent=2, default=str)
            tmp_path.replace(progress_path)
            sim_log(
                f"upper solve progress: stage={payload.get('stage', '')} "
                f"completed={payload.get('completed', '')}/{payload.get('total', '')} "
                f"best={payload.get('best_fun', '')}"
            )
\end{lstlisting}

\subsection{L3022 関数 main.propagate\_execution\_step}
            \begin{itemize}[leftmargin=1.6em]
              \item 行範囲: L3022--L3199
              \item 所属: \path|main|
              \item 定義: \path|propagate_execution_step(cmd_kmh, z0, Tb0, s0_km, env, *, force_stop = False, race_distance_km = None, max_elapsed_sec = None, env_provider = None)|
              \item docstring: 明示されていない。
              \item このブロックが直接呼ぶ主な関数・メソッド: advance\_rate\_limiter\_to\_distance\_boundary, append, bool, clear, clip, dict, electrical\_balance, env\_provider, float, get, limit\_step\_duration\_to\_distance, max
              \item 戻り値の要点: (z, Tb, s, \{'elapsed\_sec': float(elapsed\_sec), 'v\_exec\_kmh': float(sum\_v\_exec\_dt / denom), 'v\_exec\_target\_kmh': float(sum\_v\_target\_dt / denom), 'v\_exec\_ms': float(sum\_v\_exec\_dt / denom / 3.6), 'out': avg\_out, 'forces': avg\_forces, 'substeps': substeps\}) / (z0, Tb0, float(race\_distance\_km), \{'elapsed\_sec': 0.0, 'v\_exec\_kmh': 0.0, 'v\_exec\_target\_kmh': 0.0, 'v\_exec\_ms': 0.0, 'out': \{\}, 'forces': \{\}, 'substeps': []\})
              \item 主なローカル代入名: Tb, Tb\_before, \_, acceleration\_ms2, avg\_forces, avg\_out, cmd\_kmh, delayed\_cmd, denom, dt\_sub, elapsed\_sec, env\_sub, exec\_sim\_time\_sec, exec\_target\_kmh, force\_metric\_keys, forces, headwind\_ms, inertial\_power\_w, key, kinetic\_delta\_j
              \item 読み取る主なインスタンス属性: 特になし。
              \item 更新する主なインスタンス属性: 特になし。
              \item 明示的に送出する例外: 明示的raiseはない。
              \item 制御構造の規模: 条件分岐 10、ループ 4、try 0
            \end{itemize}
            \paragraph{この定義を読むためのPython構文}
            \begin{itemize}[leftmargin=1.6em]
            \item 内包表記は、反復・条件・値生成を一つの式で記述してcollectionを作る。
            \end{itemize}
            \paragraph{上から順の処理}
            \begin{enumerate}[leftmargin=1.8em]
            \item Nonlocal 文を実行する。
\item slope\_pct に float(env['slope\_pct']) の結果を代入する。
\item headwind\_ms に float(env['headwind\_ms']) の結果を代入する。
\item step\_sec に max(1e-06, float(args.dt)) の結果を代入する。
\item 条件 max\_elapsed\_sec is not None を判定し、真なら内部処理を行う。
\item   step\_sec に min(step\_sec, max(1e-06, float(max\_elapsed\_sec))) の結果を代入する。
\item cmd\_kmh に float(np.clip(cmd\_kmh, 0.0, v\_max\_kmh)) の結果を代入する。
\item 条件 race\_distance\_km is not None and s0\_km >= float(race\_distance\_km) - DISTANCE\_EPS\_KM を判定し、真なら内部処理を行う。
\item   (z0, Tb0, float(race\_distance\_km), \{'elapsed\_sec': 0.0, 'v\_exec\_kmh': 0.0, 'v\_exec\_target\_kmh': 0.0, 'v\_exec\_ms': 0.0, 'out': \{\}, 'forces': \{\}, 'substeps': []\}) を返す。
\item out\_metric\_keys に ('P\_pv\_raw', 'P\_pv\_unlimited', 'P\_pv\_limit\_loss', 'P\_pv', 'eta\_panel', 'eta\_mppt', 'P\_aux', 'P\_road\_load', 'P\_inertia', 'P\_mech', 'P\_mech\_wheel', 'P\_dc\_to\_drv', 'P\_reg\_to\_dc', 'P\_pack', 'I', 'V', 'OCV', 'Rint', 'Rline', 'Rpolarization', 'Rtotal', 'iv\_discriminant', 'eta\_charge', 'losses\_int', 'losses\_rint', 'losses\_line', 'losses\_polarization', 'eff\_drv', 'eff\_reg', 'torque\_drive\_nm', 'torque\_regen\_nm', 'omega\_motor\_radps', 'omega\_wheel\_radps') の結果を代入する。
\item force\_metric\_keys に ('F\_aero', 'F\_roll', 'F\_grade', 'F\_total', 'theta', 'v\_rel\_ms', 'normal\_force\_n', 'Crr\_eff', 'slope\_scaled\_pct') の結果を代入する。
\item 条件 exec\_model\_enabled を判定し、真なら内部処理を行う。
\item   条件 force\_stop or cmd\_kmh <= 0.1 を判定し、真なら内部処理を行う。
\item     exec\_delay\_queue.clear(...) を実行する。
\item     exec\_delay\_queue.append(...) を実行する。
\item     上の条件が偽の場合:
\item     exec\_delay\_queue.append(...) を実行する。
\item z に float(z0) の結果を代入する。
\item Tb に float(Tb0) の結果を代入する。
\item s に float(s0\_km) の結果を代入する。
\item remaining に step\_sec の結果を代入する。
\item elapsed\_sec に 0.0 の結果を代入する。
\item sum\_v\_exec\_dt に 0.0 の結果を代入する。
\item sum\_v\_target\_dt に 0.0 の結果を代入する。
\item sum\_metrics に \{key: 0.0 for key in out\_metric\_keys + force\_metric\_keys\} の結果を代入する。
\item substeps に [] の結果を代入する。
\item 条件 remaining > 1e-09 が成り立つ間くり返す。
\item   条件 race\_distance\_km is not None and s >= float(race\_distance\_km) - DISTANCE\_EPS\_KM を判定し、真なら内部処理を行う。
\item     s に float(race\_distance\_km) の結果を代入する。
\item     Break 文を実行する。
\item   dt\_sub に min(exec\_inner\_dt\_sec, remaining) の結果を代入する。
\item   v\_previous\_kmh に float(v\_exec) の結果を代入する。
\item   条件 exec\_model\_enabled を判定し、真なら内部処理を行う。
\item     条件 exec\_delay\_queue and float(exec\_delay\_queue[0][0]) <= exec\_sim\_time\_sec + 1e-09 が成り立つ間くり返す。
\item       (\_, delayed\_cmd) に exec\_delay\_queue.popleft() の結果を代入する。
\item       exec\_target\_kmh に float(np.clip(delayed\_cmd, 0.0, v\_max\_kmh)) の結果を代入する。
\item     (v\_exec, dt\_sub) に advance\_rate\_limiter\_to\_distance\_boundary(exec\_limiter, exec\_target\_kmh, start\_time\_sec=exec\_sim\_time\_sec, step\_sec=dt\_sub, s0\_km=s, distance\_limit\_km=race\_distance\_km) の結果を代入する。
\item     (v\_exec, stop\_snapped) に snap\_execution\_stop\_speed\_kmh(v\_exec, exec\_target\_kmh, stop\_requested=bool(force\_stop or cmd\_kmh <= 0.1), deadband\_kmh=exec\_deadband\_kmh, quantize\_step\_kmh=exec\_quantize\_step\_kmh) の結果を代入する。
\item     条件 stop\_snapped を判定し、真なら内部処理を行う。
\item       exec\_limiter.reset(...) を実行する。
\item     上の条件が偽の場合:
\item     exec\_target\_kmh に float(cmd\_kmh) の結果を代入する。
\item     v\_exec に float(cmd\_kmh) の結果を代入する。
\item     dt\_sub に limit\_step\_duration\_to\_distance(dt\_sub, v\_exec, s, race\_distance\_km) の結果を代入する。
\item   条件 dt\_sub <= 1e-09 を判定し、真なら内部処理を行う。
\item     条件 race\_distance\_km is not None を判定し、真なら内部処理を行う。
\item       s に min(float(race\_distance\_km), float(s)) の結果を代入する。
\item     Break 文を実行する。
\item   env\_sub に env\_provider(elapsed\_sec, s) if env\_provider is not None else env の結果を代入する。
\item   slope\_pct に float(env\_sub['slope\_pct']) の結果を代入する。
\item   headwind\_ms に float(env\_sub['headwind\_ms']) の結果を代入する。
\item   v\_ms に v\_exec / 3.6 の結果を代入する。
\item   v\_previous\_ms に v\_previous\_kmh / 3.6 の結果を代入する。
\item   acceleration\_ms2 に (v\_ms - v\_previous\_ms) / max(dt\_sub, 1e-09) の結果を代入する。
\item   kinetic\_delta\_j に 0.5 * model.p.m * (v\_ms * v\_ms - v\_previous\_ms * v\_previous\_ms) の結果を代入する。
\item   inertial\_power\_w に kinetic\_delta\_j / max(dt\_sub, 1e-09) の結果を代入する。
\item   z\_before に float(z) の結果を代入する。
\item   Tb\_before に float(Tb) の結果を代入する。
\item   s\_before に float(s) の結果を代入する。
\item   out に model.electrical\_balance(v\_ms, slope\_pct, z, Tb, env\_sub['G\_poa'], env\_sub['Tcell\_C'], headwind\_ms=headwind\_ms, aux\_power\_w=env\_sub.get('aux\_power\_w'), inertial\_power\_w=inertial\_power\_w, ambient\_temp\_c=env\_sub.get('Tamb\_C'), elevation\_m=env\_sub.get('elevation\_m', 0.0)) の結果を代入する。
\item   loss\_int に float(out['losses\_int']) の結果を代入する。
\item   forces に model.resistive\_forces(v\_ms, slope\_pct, headwind\_ms=headwind\_ms, ambient\_temp\_c=env\_sub.get('Tamb\_C'), elevation\_m=env\_sub.get('elevation\_m', 0.0)) の結果を代入する。
\item   z に model.soc\_step(z, float(out['P\_pack']), dt\_sub, current\_a=float(out['I']), Tbat\_C=Tb) の結果を代入する。
\item   Tb に Tb + dt\_sub / 1800.0 * (env\_sub['Tamb\_C'] - Tb) + loss\_int * dt\_sub / 50000.0 の結果を代入する。
\item   z に float(np.clip(z, model.p.soc\_min, model.p.soc\_max)) の結果を代入する。
\item   Tb に float(np.clip(Tb, model.p.T\_min, model.p.T\_max)) の結果を代入する。
\item   s を Add で更新する。
\item   条件 race\_distance\_km is not None を判定し、真なら内部処理を行う。
\item     race\_distance に float(race\_distance\_km) の結果を代入する。
\item     s に race\_distance if s >= race\_distance - DISTANCE\_EPS\_KM else min(race\_distance, float(s)) の結果を代入する。
\item   elapsed\_sec を Add で更新する。
\item   sum\_v\_exec\_dt を Add で更新する。
\item   sum\_v\_target\_dt を Add で更新する。
\item   out\_metric\_keys を順に走査し、各要素を key に入れて処理する。
\item     sum\_metrics[key] を Add で更新する。
\item   force\_metric\_keys を順に走査し、各要素を key に入れて処理する。
\item     sum\_metrics[key] を Add で更新する。
\item   substeps.append(...) を実行する。
\item   exec\_sim\_time\_sec を Add で更新する。
\item   remaining を Sub で更新する。
            \end{enumerate}
            \paragraph{代表コード断片}
            \begin{lstlisting}[language=Python]
    def propagate_execution_step(
        cmd_kmh,
        z0,
        Tb0,
        s0_km,
        env,
        *,
        force_stop=False,
        race_distance_km=None,
        max_elapsed_sec=None,
        env_provider=None,
    ):
        nonlocal exec_sim_time_sec, exec_target_kmh, v_exec
        slope_pct = float(env['slope_pct'])
        headwind_ms = float(env['headwind_ms'])
        step_sec = max(1.0e-6, float(args.dt))
        if max_elapsed_sec is not None:
            step_sec = min(step_sec, max(1.0e-6, float(max_elapsed_sec)))
        cmd_kmh = float(np.clip(cmd_kmh, 0.0, v_max_kmh))
        if race_distance_km is not None and s0_km >= (float(race_distance_km) - DISTANCE_EPS_KM):
            return z0, Tb0, float(race_distance_km), {
                'elapsed_sec': 0.0,
                'v_exec_kmh': 0.0,
                'v_exec_target_kmh': 0.0,
                'v_exec_ms': 0.0,
                'out': {},
                'forces': {},
                'substeps': [],
            }
        out_metric_keys = (
            'P_pv_raw', 'P_pv_unlimited', 'P_pv_limit_loss', 'P_pv',
            'eta_panel', 'eta_mppt', 'P_aux',
            'P_road_load', 'P_inertia', 'P_mech', 'P_mech_wheel', 'P_dc_to_drv', 'P_reg_to_dc',
            'P_pack', 'I', 'V', 'OCV', 'Rint', 'Rline', 'Rpolarization', 'Rtotal',
            'iv_discriminant', 'eta_charge', 'losses_int', 'losses_rint', 'losses_line', 'losses_polarization',
...
\end{lstlisting}

\subsection{L3201 関数 main.write\_execution\_detail}
            \begin{itemize}[leftmargin=1.6em]
              \item 行範囲: L3201--L3364
              \item 所属: \path|main|
              \item 定義: \path|write_execution_detail(substep, outer_time_utc, *, drive_now, hard_stop, upper_speed_cmd_kmh, speed_safety_override_active, soc_guard_intervened, plan_id, outer_step_index, outer_step_dt_sec, outer_boundary_reason)|
              \item docstring: 明示されていない。
              \item このブロックが直接呼ぶ主な関数・メソッド: abs, append, astimezone, bool, float, get, int, isoformat, max, min, select\_drive\_mode, startswith
              \item 戻り値の要点: row
              \item 主なローカル代入名: acceleration\_ms2\_sub, detail\_min\_soc, dt\_sub, energy\_fields, env\_sub, forces\_sub, label, lower\_command\_index, out\_sub, power\_key, row, s\_end, s\_start, soc\_guard\_intervention\_rows, soc\_guard\_intervention\_sec, step\_end\_utc, step\_time\_utc, torque\_drive\_nm, v\_cmd\_sub, v\_exec\_ms\_sub
              \item 読み取る主なインスタンス属性: 特になし。
              \item 更新する主なインスタンス属性: 特になし。
              \item 明示的に送出する例外: 明示的raiseはない。
              \item 制御構造の規模: 条件分岐 1、ループ 1、try 0
            \end{itemize}
            \paragraph{この定義を読むためのPython構文}
            \begin{itemize}[leftmargin=1.6em]
            \item 特別な構文上の補足は少ない。
            \end{itemize}
            \paragraph{上から順の処理}
            \begin{enumerate}[leftmargin=1.8em]
            \item Nonlocal 文を実行する。
\item Nonlocal 文を実行する。
\item dt\_sub に float(substep['dt\_sec']) の結果を代入する。
\item step\_time\_utc に outer\_time\_utc + timedelta(seconds=float(substep['elapsed\_from\_step\_sec'])) の結果を代入する。
\item step\_end\_utc に step\_time\_utc + timedelta(seconds=dt\_sub) の結果を代入する。
\item out\_sub に substep['out'] の結果を代入する。
\item forces\_sub に substep['forces'] の結果を代入する。
\item env\_sub に substep['env'] の結果を代入する。
\item z\_start に float(substep['z\_start']) の結果を代入する。
\item z\_end に float(substep['z\_end']) の結果を代入する。
\item s\_start に float(substep['s\_start\_km']) の結果を代入する。
\item s\_end に float(substep['s\_end\_km']) の結果を代入する。
\item v\_cmd\_sub に float(substep['v\_cmd\_kmh']) の結果を代入する。
\item v\_exec\_sub に float(substep['v\_exec\_kmh']) の結果を代入する。
\item v\_exec\_previous\_sub に float(substep.get('v\_exec\_previous\_kmh', v\_exec\_sub)) の結果を代入する。
\item v\_target\_sub に float(substep['v\_exec\_target\_kmh']) の結果を代入する。
\item v\_exec\_ms\_sub に v\_exec\_sub / 3.6 の結果を代入する。
\item acceleration\_ms2\_sub に float(substep.get('acceleration\_ms2', 0.0)) の結果を代入する。
\item energy\_fields に \{\} の結果を代入する。
\item energy\_power\_keys を順に走査し、各要素を power\_key に入れて処理する。
\item   label に power\_key[2:] if power\_key.startswith('P\_') else power\_key の結果を代入する。
\item   value\_wh に float(out\_sub.get(power\_key, 0.0)) * dt\_sub / 3600.0 の結果を代入する。
\item   cumulative\_energy\_wh[power\_key] を Add で更新する。
\item   energy\_fields[f'E\_\{label\}\_step\_Wh'] に value\_wh の結果を代入する。
\item   energy\_fields[f'E\_\{label\}\_cumulative\_Wh'] に cumulative\_energy\_wh[power\_key] の結果を代入する。
\item lower\_command\_index を Add で更新する。
\item 条件 soc\_guard\_intervened を判定し、真なら内部処理を行う。
\item   soc\_guard\_intervention\_rows を Add で更新する。
\item   soc\_guard\_intervention\_sec を Add で更新する。
\item detail\_min\_soc に min(detail\_min\_soc, z\_start, z\_end) の結果を代入する。
\item detail\_tracking\_errors.append(...) を実行する。
\item torque\_drive\_nm に float(out\_sub.get('torque\_drive\_nm', 0.0)) の結果を代入する。
\item row に \{'lower\_command\_index': int(lower\_command\_index), 'outer\_step\_index': int(outer\_step\_index), 'time\_utc': step\_time\_utc.isoformat(), 'time\_end\_utc': step\_end\_utc.isoformat(), 'time\_local': step\_time\_utc.astimezone(detail\_local\_tz).isoformat(), 'step\_dt\_sec': dt\_sub, 'detail\_target\_dt\_sec': detail\_interval\_sec, 'detail\_step\_kind': 'lower\_command\_cycle' if abs(dt\_sub - detail\_interval\_sec) <= 1e-09 else 'boundary\_remainder', 'lower\_command\_rate\_hz': 1.0 / max(dt\_sub, 1e-09), 'outer\_step\_requested\_dt\_sec': float(args.dt), 'outer\_step\_actual\_dt\_sec': float(outer\_step\_dt\_sec), 'outer\_step\_boundary\_reason': str(outer\_boundary\_reason or 'planner\_step\_complete'), 'is\_drive\_window': bool(drive\_now), 'is\_forced\_stop': bool(hard\_stop), 'plan\_id': int(plan\_id), 's\_km': s\_start, 's\_end\_km': s\_end, 'ds\_step\_km': s\_end - s\_start, 'upper\_speed\_cmd\_kmh': float(upper\_speed\_cmd\_kmh), 'speed\_safety\_override\_active': bool(speed\_safety\_override\_active), 'soc\_guard\_intervened': bool(soc\_guard\_intervened), 'lower\_reference\_target\_kmh': v\_target\_sub, 'lower\_speed\_cmd\_kmh': v\_exec\_sub, 'v\_cmd\_kmh': v\_cmd\_sub, 'v\_exec\_kmh': v\_exec\_sub, 'v\_exec\_previous\_kmh': v\_exec\_previous\_sub, 'v\_exec\_target\_kmh': v\_target\_sub, 'v\_cmd\_ms': v\_cmd\_sub / 3.6, 'v\_exec\_ms': v\_exec\_ms\_sub, 'acceleration\_ms2': acceleration\_ms2\_sub, 'soc': z\_start, 'soc\_end': z\_end, 'usable\_battery\_start\_Wh': max(0.0, (z\_start - model.p.soc\_min) * model.p.E\_nom\_Wh), 'usable\_battery\_end\_Wh': max(0.0, (z\_end - model.p.soc\_min) * model.p.E\_nom\_Wh), 'battery\_energy\_used\_step\_Wh': float(out\_sub.get('P\_pack', 0.0)) * dt\_sub / 3600.0, 'battery\_charge\_used\_step\_Ah': (z\_start - z\_end) * float(model.p.Q\_nom\_Ah), 'soc\_state\_definition': 'charge' if float(model.p.Q\_nom\_Ah) > 0.0 else 'energy', 'Tb\_C': float(substep['Tb\_start']), 'Tb\_end\_C': float(substep['Tb\_end']), 'slope\_pct': float(env\_sub.get('slope\_pct', 0.0)), 'slope\_scaled\_pct': float(forces\_sub.get('slope\_scaled\_pct', env\_sub.get('slope\_pct', 0.0))), 'headwind\_ms': float(env\_sub.get('headwind\_ms', 0.0)), 'v\_rel\_ms': float(forces\_sub.get('v\_rel\_ms', v\_exec\_ms\_sub + env\_sub.get('headwind\_ms', 0.0))), 'G\_raw': float(env\_sub.get('G\_raw', env\_sub.get('G\_poa', 0.0))), 'G\_drive\_poa': float(env\_sub.get('G\_drive\_poa', env\_sub.get('G\_raw', 0.0))), 'G\_stop\_poa': float(env\_sub.get('G\_stop\_poa', env\_sub.get('G\_raw', 0.0))), 'G\_control\_stop\_poa': float(env\_sub.get('G\_control\_stop\_poa', env\_sub.get('G\_raw', 0.0))), 'G\_poa': float(env\_sub.get('G\_poa', 0.0)), 'Tamb\_C': float(env\_sub.get('Tamb\_C', 30.0)), 'Tcell\_drive\_C': float(env\_sub.get('Tcell\_drive\_C', env\_sub.get('Tcell\_C', 40.0))), 'Tcell\_stop\_C': float(env\_sub.get('Tcell\_stop\_C', env\_sub.get('Tcell\_C', 40.0))), 'Tcell\_control\_stop\_C': float(env\_sub.get('Tcell\_control\_stop\_C', env\_sub.get('Tcell\_C', 40.0))), 'Tcell\_C': float(env\_sub.get('Tcell\_C', 40.0)), 'panel\_stop\_mode': str(env\_sub.get('panel\_stop\_mode', 'drive')), 'P\_pv\_raw': float(out\_sub.get('P\_pv\_raw', 0.0)), 'P\_pv\_unlimited': float(out\_sub.get('P\_pv\_unlimited', out\_sub.get('P\_pv', 0.0))), 'P\_pv\_limit\_loss': float(out\_sub.get('P\_pv\_limit\_loss', 0.0)), 'eta\_panel': float(out\_sub.get('eta\_panel', 0.0)), 'eta\_mppt': float(out\_sub.get('eta\_mppt', 0.0)), 'P\_pv': float(out\_sub.get('P\_pv', 0.0)), 'P\_solar\_w': float(out\_sub.get('P\_pv', 0.0)), 'P\_aux': float(out\_sub.get('P\_aux', 0.0)), 'P\_road\_load': float(out\_sub.get('P\_road\_load', 0.0)), 'P\_inertia': float(out\_sub.get('P\_inertia', 0.0)), 'P\_mech': float(out\_sub.get('P\_mech', 0.0)), 'P\_mech\_wheel': float(out\_sub.get('P\_mech\_wheel', 0.0)), 'P\_dc\_to\_drv': float(out\_sub.get('P\_dc\_to\_drv', 0.0)), 'P\_reg\_to\_dc': float(out\_sub.get('P\_reg\_to\_dc', 0.0)), 'P\_pack': float(out\_sub.get('P\_pack', 0.0)), 'P\_vehicle\_load\_w': float(out\_sub.get('P\_dc\_to\_drv', 0.0)) - float(out\_sub.get('P\_reg\_to\_dc', 0.0)) + float(out\_sub.get('P\_aux', 0.0)), 'P\_net\_battery\_w': float(out\_sub.get('P\_pack', 0.0)), 'P\_battery\_terminal': float(out\_sub.get('I', 0.0)) * float(out\_sub.get('V', 0.0)), 'power\_balance\_residual\_W': float(out\_sub.get('P\_pack', 0.0)) - float(out\_sub.get('I', 0.0)) * float(out\_sub.get('V', 0.0)), 'I': float(out\_sub.get('I', 0.0)), 'V': float(out\_sub.get('V', 0.0)), 'OCV': float(out\_sub.get('OCV', 0.0)), 'Rint': float(out\_sub.get('Rint', 0.0)), 'Rline': float(out\_sub.get('Rline', 0.0)), 'Rpolarization': float(out\_sub.get('Rpolarization', 0.0)), 'Rtotal': float(out\_sub.get('Rtotal', 0.0)), 'iv\_discriminant': float(out\_sub.get('iv\_discriminant', 0.0)), 'eta\_charge': float(out\_sub.get('eta\_charge', 1.0)), 'losses\_int': float(out\_sub.get('losses\_int', 0.0)), 'losses\_rint': float(out\_sub.get('losses\_rint', 0.0)), 'losses\_line': float(out\_sub.get('losses\_line', 0.0)), 'losses\_polarization': float(out\_sub.get('losses\_polarization', 0.0)), 'eff\_drv': float(out\_sub.get('eff\_drv', 0.0)), 'eff\_reg': float(out\_sub.get('eff\_reg', 0.0)), 'torque\_drive\_nm': torque\_drive\_nm, 'torque\_regen\_nm': float(out\_sub.get('torque\_regen\_nm', 0.0)), 'omega\_motor\_radps': float(out\_sub.get('omega\_motor\_radps', 0.0)), 'omega\_wheel\_radps': float(out\_sub.get('omega\_wheel\_radps', 0.0)), 'drive\_mode\_selected': model.select\_drive\_mode(v\_exec\_ms\_sub, torque\_drive\_nm), 'F\_aero': float(forces\_sub.get('F\_aero', 0.0)), 'F\_roll': float(forces\_sub.get('F\_roll', 0.0)), 'F\_grade': float(forces\_sub.get('F\_grade', 0.0)), 'F\_total': float(forces\_sub.get('F\_total', 0.0)), 'normal\_force\_n': float(forces\_sub.get('normal\_force\_n', 0.0)), 'theta\_rad': float(forces\_sub.get('theta', 0.0)), 'Crr\_eff': float(forces\_sub.get('Crr\_eff', model.p.Crr)), 'elevation\_m': float(env\_sub.get('elevation\_m', 0.0)), 'air\_density\_kgm3': float(forces\_sub.get('air\_density\_kgm3', model.p.rho))\} の結果を代入する。
\item row.update(...) を実行する。
\item row.update(...) を実行する。
\item detail\_stream.write(...) を実行する。
\item row を返す。
            \end{enumerate}
            \paragraph{代表コード断片}
            \begin{lstlisting}[language=Python]
    def write_execution_detail(
        substep,
        outer_time_utc,
        *,
        drive_now,
        hard_stop,
        upper_speed_cmd_kmh,
        speed_safety_override_active,
        soc_guard_intervened,
        plan_id,
        outer_step_index,
        outer_step_dt_sec,
        outer_boundary_reason,
    ):
        nonlocal lower_command_index, detail_min_soc
        nonlocal soc_guard_intervention_rows, soc_guard_intervention_sec
        dt_sub = float(substep['dt_sec'])
        step_time_utc = outer_time_utc + timedelta(seconds=float(substep['elapsed_from_step_sec']))
        step_end_utc = step_time_utc + timedelta(seconds=dt_sub)
        out_sub = substep['out']
        forces_sub = substep['forces']
        env_sub = substep['env']
        z_start = float(substep['z_start'])
        z_end = float(substep['z_end'])
        s_start = float(substep['s_start_km'])
        s_end = float(substep['s_end_km'])
        v_cmd_sub = float(substep['v_cmd_kmh'])
        v_exec_sub = float(substep['v_exec_kmh'])
        v_exec_previous_sub = float(substep.get('v_exec_previous_kmh', v_exec_sub))
        v_target_sub = float(substep['v_exec_target_kmh'])
        v_exec_ms_sub = v_exec_sub / 3.6
        acceleration_ms2_sub = float(substep.get('acceleration_ms2', 0.0))

        energy_fields = {}
        for power_key in energy_power_keys:
...
\end{lstlisting}







        \subsection{CLI 引数}
            \begin{longtable}{p{0.07\linewidth} p{0.78\linewidth}}
            \toprule
            行 & 引数 \\
            \midrule
            1487 & \path|--profile_yaml| \\
1488 & \path|--forecast_csv| \\
1489 & \path|--forecast_fill_csv| \\
1490 & \path|--progress_reference_csv| \\
1491 & \path|--forecast_time_mode| \\
1492 & \path|--forecast_time_tz| \\
1493 & \path|--route_profile_csv| \\
1494 & \path|--speed_profile_csv| \\
1495 & \path|--params_yaml| \\
1496 & \path|--stop_yaml| \\
1497 & \path|--drive_schedule_yaml| \\
1498 & \path|--panel_eff_map| \\
1499 & \path|--mppt_eff_map| \\
1500 & \path|--ocv_soc_map| \\
1501 & \path|--drive_eff_map| \\
1502 & \path|--regen_eff_map| \\
1503 & \path|--rint_map| \\
1504 & \path|--drive_map_eco| \\
1505 & \path|--drive_map_power| \\
1506 & \path|--regen_map_eco| \\
1507 & \path|--regen_map_power| \\
1508 & \path|--dt| \\
1509 & \path|--horizon_steps| \\
1510 & \path|--soc0| \\
1511 & \path|--Tb0| \\
1512 & \path|--v0_kmh| \\
1513 & \path|--start_utc| \\
1514 & \path|--forecast_start_time_utc| \\
1515 & \path|--start_index| \\
1516 & \path|--start_s_km| \\
1517 & \path|--resume_csv| \\
1518 & \path|--resume_s_km| \\
1519 & \path|--out_csv| \\
1520 & \path|--out_detail_csv| \\
1521 & \path|--report_html| \\
1522 & \path|--summary_json| \\
1523 & \path|--resolved_yaml| \\
1524 & \path|--latest_manifest_json| \\
1525 & \path|--override| \\
1526 & \path|--soc_guard_margin| \\
1527 & \path|--soc_guard_mode| \\
1528 & \path|--solar_gain| \\
1529 & \path|--poa_gain_drive| \\
1530 & \path|--poa_gain_stop| \\
1531 & \path|--stop_tilt_fraction| \\
1532 & \path|--control_stop_tilt_fraction| \\
1533 & \path|--energy_budget| \\
1534 & \path|--exec_model_enabled| \\
1535 & \path|--exec_inner_dt_sec| \\
1536 & \path|--detail_rate_hz| \\
1537 & \path|--exec_tau_sec| \\
1538 & \path|--exec_accel_limit_kmhps| \\
1539 & \path|--exec_decel_limit_kmhps| \\
1540 & \path|--exec_deadband_kmh| \\
1541 & \path|--exec_quantize_step_kmh| \\
1542 & \path|--exec_reaction_delay_sec| \\
1543 & \path|--upper_mode| \\
1544 & \path|--upper_ds_km| \\
1545 & \path|--upper_horizon_km| \\
1546 & \path|--upper_max_steps| \\
1547 & \path|--upper_horizon_mode| \\
1548 & \path|--upper_adaptive_min_ds_km| \\
1549 & \path|--upper_adaptive_max_ds_km| \\
1550 & \path|--upper_adaptive_growth| \\
1551 & \path|--upper_replan_km| \\
1552 & \path|--upper_replan_sec| \\
1553 & \path|--upper_max_iter| \\
1554 & \path|--upper_global_search_enabled| \\
1555 & \path|--upper_global_search_mode| \\
1556 & \path|--upper_cem_generations| \\
1557 & \path|--upper_cem_population| \\
1558 & \path|--upper_cem_elite| \\
1559 & \path|--upper_local_refine_topk| \\
1560 & \path|--upper_shgo_samples| \\
1561 & \path|--upper_shgo_iters| \\
1562 & \path|--upper_cert_grid_levels| \\
1563 & \path|--upper_cert_max_evaluations| \\
1564 & \path|--upper_cert_workers| \\
1570 & \path|--upper_ctrl_km| \\
1571 & \path|--upper_vmin_kmh| \\
            \bottomrule
            \end{longtable}



        \section{処理の流れ}
        \begin{enumerate}[leftmargin=1.7em]
        \item CLI 引数を解釈し、main() から処理を起動する。
        \end{enumerate}

        \section{このファイルの位置づけ}
        この資料群では、\path|scripts/solar_sim.py| を
        \textbf{offline core}の中核として扱う。
        したがって、ここで扱う処理は単独で完結せず、
        前段の profile / launch / telemetry と後段の planner / logger / report へ繋がる。

        \end{document}
